\documentclass[12pt]{book}

\usepackage{amssymb,amsmath,amsfonts, graphics}
\usepackage{bm}
\usepackage{epsfig}
\usepackage{ccaption}
\usepackage{graphicx}
\usepackage{fancyhdr}
\usepackage[]{mdframed}
\usepackage{float}
\usepackage{pdfpages}
\usepackage{cite}
\usepackage[utf8]{inputenc}
\usepackage[catalan]{babel}

\usepackage[
      colorlinks=true,
      linkcolor=blue,
      urlcolor=blue,
      filecolor=black,
      citecolor=red,
      pdfstartview=FitV,
      pdftitle={},
        pdfauthor={Raimon Luna},
        pdfsubject={},
        pdfkeywords={},
        pdfpagemode=None,
        bookmarksopen=true
      ]{hyperref}


\setlength{\textheight}{21cm}
\setlength{\parindent}{0.5cm}
\setlength{\textwidth}{15cm}
\setlength{\oddsidemargin}{0.5cm}
\setlength{\evensidemargin}{0.5cm}
\linespread{1.25}
\addtolength{\arraycolsep}{-3pt}
\raggedbottom



\title{Física per a Ciències Biomèdiques}
\author{Raimon Luna i Perelló \\ Universitat de Barcelona}

\begin{document}

\maketitle

\null\vfill
\noindent
Copyright (text i il·lustracions)\\
Raimon Luna i Perelló -- 2025 \\
Departament de Física Quàntica i Astrofísica\\
Institut de Ciències del Cosmos\\
Universitat de Barcelona
\newpage

\tableofcontents

\chapter*{Introducció}

Aquest text és una introducció a la física clàssica per al primer curs d'un grau de l'àrea de les ciències biomèdiques. Està dividit en tres parts principals: fluids, termodinàmica i ones. Cada part consta d'un o més capítols. Tot i que els conceptes físics que s'hi tracten són d'ús general en qualsevol ciència o enginyeria, en aquest text s'ha tractat d'emfatitzar la seva rellevància en el món dels éssers vius.

Les explicacions en aquest text intenten ser tant sintètiques i senzilles com sigui possible, i amb el mínim formalisme matemàtic, sense perdre el rigor. S'ha intentat donar molta importància al suport visual de cada concepte, a través d'il·lustracions esquemàtiques però clares de cada situació física. Les figures basades en visualitzacions geomètriques 3D estan fetes amb el programari Mathematica, mentre que les representacions 2D són codis en Python. En alguns casos les imatges són de simulacions realitzades amb eines matemàtiques avançades que escapen als objectius del curs, com ara els {\it lattice Boltzmann methods} o les transformacions de Joukowsky.

Els coneixements previs necessaris per entendre el contingut són els de matemàtiques del batxillerat. Això inclou la manipulació simbòlica d'equacions, trigonometria i nocions de càlcul diferencial i integral. Des del punt de vista de la física, és molt recomanable conèixer les magnituds i unitats de mesura habituals, i els conceptes bàsics de la cinemàtica i la dinàmica newtoniana. Tot i això, la física que apareix en aquest curs no és difícil d'entendre de forma intuïtiva, no només a través del formalisme matemàtic. Durant tot el text s'encoratja el lector a desenvolupar aquesta intuïció física, que a llarg termini és una eina de gran valor.

\newpage

La primera part es destina a la mecànica de fluids, dividida en tres capítols: Estàtica de fluids, Dinàmica de fluids ideals, i Fluids viscosos. Aquesta part és de gran rellevància a la biologia i la biotecnologia, ja que la major part de processos dels éssers vius tenen lloc dins l'aigua. Per tant, és important entendre bé els conceptes de pressió, fluxos, cabal i forces viscoses. La física dels fluids en el món macroscòpic és la responsable del desplaçament dels animals aquàtics i del vol dels ocells i insectes.

La segona part, la més curta, introdueix principis de termodinàmica. Inclou conceptes molt importants com el calor i la temperatura, i els mecanismes de transmissió de calor. També introdueix la noció d'energia des del punt de vista de la física, i els principis de la termodinàmica que en controlen el seu comportament a escala macroscòpica. Algunes de les seves aplicacions immediates són a la química i la bioquímica, però també juga un paper clau en la relació dels éssers vius amb el seu hàbitat i clima.

Finalment, els últims tres capítols es dediquen als fenòmens periòdics i les ones: Ones mecàniques, Ones sonores i Ones electromagnètiques. En aquest cas, trobem que les ones són un fenomen que apareix en aspectes aparentment molt diferents de la vida a la Terra. En particular, dos dels nostres sentits principals es basen completament en la sensibilitat a dos tipus d'ones, l'audició (ones sonores) i la visió (ones electromagnètiques). També són útils en el propi estudi dels sistemes biològics, a través dels microscopis. El fenomen de difracció de les ones electromagnètiques, concretament els raigs X, va conduir al descobriment de l'estructura de l'ADN a través de la famosa Fotografia 51 obtinguda per Rosalind Franklin en col·laboració amb Raymond Gosling.


\part{Física de Fluids}

\chapter{Estàtica de fluids}
\label{chap:Hidrostatica}

\section{Conceptes generals}

Els fluids són un dels tres estats d'agregació clàssics de la matèria, juntament amb els sòlids i els gasos. En el camp de la biologia i biotecnologia, els fluids són probablement l'estat més important, ja que la major part de sistemes biològics tenen lloc en un medi aquós. Això succeeix tant per als organismes unicel·lulars (bacteris, protozous, algues, fongs...) com en l'interior dels teixits dels organismes pluricel·lulars (sang, sistema respiratori, citoplasma cel·lular...). També tenen una gran importància com a medi de locomoció, per un costat per als éssers vius que neden en l'aigua de mars, oceans, rius i llacs; i per l'altre per als que es desplacen per l'aire a través del vol i del moviment dels vents. Al llarg d'aquesta part del llibre, intentarem explicar alguns d'aquests fenòmens i la física que els fa possibles. Abans de començar, però, és important introduir una sèrie de conceptes i magnituds que seran molt importants per poder entendre la física dels fluids. En primer lloc, què és exactament un fluid?

En la nostra experiència quotidiana, veiem tres tipus de matèria: sòlids, líquids i gasos. En general, considerarem que els dos últims, líquids i gasos, són fluids. Una característica molt important que els diferencia és que els sòlids solen mantenir la seva forma i volum, que només es poden modificar fent un esforç considerable per deformar-los. En canvi, els fluids no tenen forma pròpia i s'adapten al recipient que els conté. 

En el cas dels líquids, veiem que ocupen la part inferior d'aquest recipient, i deixen una superfície lliure a la part superior, que anomenem menisc. També veiem que són gairebé incompressibles, és a dir, mantenen molt bé el seu volum. Això es pot comprovar fàcilment tapant amb el dit una xeringa plena d'aigua i intentant apretar l'èmbol.  

Els gasos, en canvi, solen ocupar tot l'espai que tenen disponible, i per això s'han de guardar en recipients tancats, com bombones. A més, són molt compressibles, no costa gaire canviar-ne el volum. Si repetim l'experiment anterior amb una xeringa plena d'aire, veurem que no costa gaire esforç comprimir-lo. Això és perquè les forces d'interacció entre les molècules (o àtoms) que formen un gas solen ser molt més febles que en el cas d'un líquid o un sòlid.

En els dos primers capítols d'aquest llibre, estudiarem el que s'anomena un fluid ideal. Aquest és un concepte que, tot i no ser totalment realista quan el comparem amb els fluids reals, simplifica enormement els nostres càlculs i en molts casos és suficient per treballar. El primer pas que prendrem és ignorar per complet la seva composició a nivell microscòpic. És a dir, no ens importarà si està format per molècules, si té altres components en solució, si hi ha cèl·lules al seu interior... Nosaltres el veurem com un medi continu sense estructura interior. La segona cosa que assumirem és que és totalment incompressible, és a dir que el seu volum (i densitat) no canviaran mai, independentment de la pressió a la que el sotmetem. 

I finalment, mentre que un fluid ideal pot suportar forces de compressió (també anomenades forces normals), no pot suportar forces tangencials (també anomenades de forces de tall o de cisallament). Si apareix una força tangencial, com està il·lustrat a la Figura \ref{fig:ForcesTangencials}, el fluid no oposarà cap resistència i per tant fluirà en la direcció que li imposi la força. Com veurem més endavant, aquesta suposició no és del tot correcta en molts dels fluids biològics, ja que aquests solen oposar una petita resistència a fluir anomenada viscositat. Estudiarem els efectes de la viscositat al capítol \ref{chap:Viscositat}, però de moment ens centrarem en fluids ideals, no viscosos.
%
\begin{figure}[H]
    \centering
    \includegraphics[width=0.8\linewidth]{Figures/ForcesTangencials.pdf}
    \caption{Un fluid ideal oposa resistència a les forces normals, però no a les tangencials.}
    \label{fig:ForcesTangencials}
\end{figure}
%

\subsection{Magnituds característiques}

Abans de començar, anem a introduir algunes magnituds i conceptes que ens seran molt útils en la física dels fluids. La majoria són bastant conegudes, almenys de manera intuitiva, en la nostra vida quotidiana.

\paragraph{Pressió:} La pressió $P$ és la força que fa un fluid contra qualsevol superfície, real o imaginària, que hi estigui en contacte. És a dir, un fluid exerceix una pressió contra les parets del seu recipient, contra els cossos que hi ha submergits, i contra altres parts del propi fluid. La pressió es mesura com una força (mesurada en Newtons en el SI) per unitat de superfície (mesurada en metres quadrats en el SI). En general, si sobre una superfície $S$ hi actua una força normal $F$, la pressió es defineix com el límit de $F/S$ quan $S$ es fa molt petita, és a dir,
%
\begin{equation}
    P = \lim_{S \to 0} \left(\frac{F}{S}\right)\, .
\end{equation}
%
És important recordar que la pressió és una magnitud escalar, és a dir, no té ni direcció ni sentit. Això és perquè només estem considerant la força perpendicular a la superfície, que és la mateixa en qualsevol direcció. Tècnicament, diem que la pressió és isotròpica. La unitat del Sietema Internacional d'unitats (SI) per a la pressió, 1 Newton per metre  quadrat, s'anomena Pascal, però hi ha altres unitats:

\begin{itemize}
    \item Pascal:    1 Pa = 1 N/m$^2$
    \item Atmosfera: 1 atm = 1,01325 $\cdot 10^5$ Pa
    \item Bar:       1 bar = 10$^5$ Pa
    \item Mil·límetres de mercuri: 760 mmHg = 1 atm
\end{itemize}

\paragraph{Densitat:} La densitat $\rho$ és simplement la magnitud que ens indica si el fluid és pesat o lleuger, és a dir, la seva massa per unitat de volum. De manera semblant a la definició de pressió, prenem un element de volum $V$ i en mesurem la massa $M$. La densitat es defineix com el límit de $M/V$ quan $V$ es fa molt petit,
%
\begin{equation}
    \rho = \lim_{V \to 0} \left(\frac{M}{V}\right)\, .
\end{equation}
%
Novament, la densitat és una magnitud escalar. La unitat del (SI) per a la densitat és el quilogram per metre cúbic (kg/m$^3$), tot i que en el laboratori se sol utilitzar el gram per centímetre cúbic (g/cm$^3$). A la Taula \ref{tab:Densitats} es mostren els valors de densitat d'algunes substàncies quotidianes.

\begin{table}[H]
    \centering
    \begin{tabular}{|l|c|}
    \hline
                      & $\rho$ \\ \hline
        Aire sec      & 1,29 $\cdot$ 10$^{-3}$ g/cm$^3$ \\ \hline
        Oli d'oliva   & 0,92 g/cm$^3$ \\ \hline
        Aigua (a 4ºC) & 1,00 g/cm$^3$ \\ \hline
        Ferro         & 7,86 g/cm$^3$ \\ \hline
        Mercuri       & 13,6 g/cm$^3$ \\ \hline
    \end{tabular}
    \caption{Densitats d'alguns materials habituals.}
    \label{tab:Densitats}
\end{table}

\paragraph{Compressibilitat:} La compressibilitat $\kappa$ mesura la facilitat de comprimir, o canviar de volum, el fluid. La definim com la variació (relativa) de volum quan fem un canvi en la pressió. També es pot definir com la variació relativa de densitat, tot i que en aquest cas el signe global serà diferent (si el volum disminueix, la densitat augmenta, i viceversa). Això es pot expressar com una derivada

\begin{equation}
    \kappa = \frac{d\rho/\rho}{dP} = - \frac{dV/V}{dP} \, .
\end{equation}

Típicament, la compressibilitat dels gasos serà molt més alta que la dels líquids, com ja hem comentat. En el cas d'un fluid ideal, ja hem dit que és incompressible, i per tant la seva compressibilitat és zero. És una magnitud escalar, i la les seves dimensions són l'invers d'una pressió. En el SI mesurarem la compressibilitat en Pa$^{-1}$, o metres quadrats per Newton (m$^2$/N).

\subsection{L'equació d'estat}

L'equació d'estat d'un fluid és una llei que relaciona les seves magnituds entre si. Per exemple, ens diu quina és la densitat del fluid en funció de la seva pressió i temperatura. La més famosa és l'equació d'estat dels gasos ideals
%
\begin{equation}
    P V = n R T\, ,
\end{equation}
%
on $n$ és el nombre de mols i $R$ = 8,31 J / K mol és la constant dels gasos ideals. Una extensió d'aquesta equació, que té en compte la interacció entre molècules, és l'equació dels gasos de van der Waals
%
\begin{equation}
    \left(P + \frac{a}{V^2}\right)\left(V - b\right) = nRT\, ,
\end{equation}
%
on $a$ i $b$ són dues constants que depenen del gas. És important recordar que les magnituds que apareixen a l'equació d'estat poden ser intensives o extensives:

\begin{itemize}
    \item Magnitud intensiva: No depèn de la mida del sistema (temperatura, densitat...)
    \item Magnitud extensiva: És proporcional a la mida del sistema (massa, volum...)
\end{itemize}

Una manera senzilla de recordar la distinció entre magnituds intensives i extensives és pensar què passaria si dividim el nostre sistema, per exemple un got d'aigua, en dues parts iguals. La temperatura i la densitat del l'aigua no variaran (són intensives), mentre que la massa i el volum de cada part seran la meitat (són extensives).

\subsection{L'anomalia de l'aigua}

En relació a les equacions d'estat que apareixen a les ciències biològiques, la de l'aigua és particularment important, i té algunes propietats molt peculiars que han fet possible l'aparició de la vida a la Terra. La gran majoria de substàncies augmenten de volum, és a dir disminueixen la seva densitat, en augmentar la temperatura. En el cas de l'aigua, això no sempre es compleix, ja que la seva densitat assoleix un màxim a 4ºC, i disminueix per a temperatures més baixes. Això vol dir que el gel, que apareix a 0ºC, és més lleuger que l'aigua líquida que l'envolta. Ho podeu comprovar si deixeu una ampolla o llauna de refresc al congelador massa temps: en congelar-se, la beguda augmentarà de volum i molt probablement trencarà l'envàs.

Aquest fet ha estat crucial per fer possible la vida sota l'aigua en situacions de fred extrem. Gràcies a que el gel és menys dens que l'aigua líquida, flotarà cap a la superfície i formarà una capa a la part superior. El cas més típic és el d'un llac gelat. Aquesta capa de gel actuarà com un aïllant que evitarà que l'aigua de sota es congeli completament, i per tant la vida aquàtica podrà sobreviure l'hivern. Si la fase sòlida (el gel) s'enfonsés, com passa en la majoria de líquids, els llacs i oceans es podrien congelar sencers, eliminant així la major part d'essers vius que hi visquessin. 

\section{Pressió hidroestàtica}

En aquest punt, ja tenim les eines necessàries per estudiar la física dels fluids! Començarem analitzant els fluids que es troben en repòs, és a dir aquells que no es mouen, normalment perquè estan continguts dins d'un recipient o una cavitat sense cap força que els impulsi a fluir. La teoria que estudia els fluids en repòs s'anomena hidroestàtica. 

El primer concepte que és important entendre és el de pressió hidroestàtica, és a dir, la pressió que exerceix cada part del fluid sobre les parets del recipient i sobre les altres parts del fluid. Aquesta pressió ha de produir un equilibri de forces sobre totes les parts, de manera que el fluid es mantingui aturat. Per poder conèixer aquesta pressió, és molt útil pensar en un dibuix com el de la Figura \ref{fig:PressioHidrostatica}. Imaginem que tenim un got ple d'aigua, que hem deixat reposar durant una estona, de manera que l'aigua està completament aturada. En aquesta situació, podem pensar en un petit cub de fluid al seu interior, i calcular quina ha de ser la pressió per tal que es mantingui en repòs. 

Per simetria, les forces que actuen sobre les seves parets laterals, que hem etiquetat com a $F_{lat}$, tenen totes la mateixa magnitud. Per tant, les forces laterals sobre cares oposades s'anul·laran mútuament i el cub no es desplaçarà de costat.
%
\begin{figure}[H]
    \centering
    \includegraphics[width=0.5\linewidth]{Figures/PressioHidrostatica.pdf}
    \caption{Forces exercides per la pressió hidroestàtica sobre un element cúbic de fluid. Les forces laterals $F_{lat}$ es cancel·len entre elles, mentre que les forces verticals $F_1$ i $F_2$ han de contrarestar el pes del cub.}
    \label{fig:PressioHidrostatica}
\end{figure}
%
 En la direcció vertical és una mica més complicat. Tenim dues forces que empenyen el cub cap a baix, i una força que l'empeny cap a dalt. Si el cub no es mou tampoc verticalment, cal que la suma vectorial d'aquestes forces sigui zero. Anem a analitzar les forces una per una:

\begin{itemize}
    \item La pressió a la part superior del cub, que anomenarem $P_1$. Aquesta pressió produirà una força $F_1 = L^2 P_1$ cap avall, on $L$ és la longitud del costat del cub.
    \item La pressió a la part inferior del cub, que anomenarem $P_2$. Aquesta pressió produirà una força $F_2 = L^2 P_2$ cap amunt.
    \item El pes del cub, que serà una força $m g = \rho L^3 g$ cap avall, on $\rho$ és la densitat de l'aigua i $g$ = 9,81 m/s$^2$ és l'acceleració de la gravetat.
\end{itemize}

Si el cub està en equilibri, s'ha de complir 
%
\begin{equation}
    F_1 + m g = F_2\, ,
\end{equation}
%
és a dir que, substituint els valors de les forces,
%
\begin{equation}
    P_1 + \rho g L = P_2\, .
\end{equation}
%
Ara ja podem saber quin és l'increment de pressió, que anomenarem $\Delta P = P_2 - P_1$, entre les altures de la part superior i inferior del cub, que és precisament la longitud del costat del cub $\Delta h = L$. D'aquesta manera, obtenim una de les equacions més importants del curs:
%
\begin{equation}
    \boxed{\Delta P = \rho g \Delta h}
    \label{eqn:PressioHidroestatica}
\end{equation}
%
Aquesta equació ens diu que la pressió en un fluid en repòs augmenta amb la profunditat. Com més profund ens submergim, major serà la pressió que el fluid exercirà sobre nosaltres. 

Això és una cosa que notem en la nostra vida quotidiana, per exemple en banyar-nos en una piscina. Un cas extrem és el d'un submarí que es desplaça a centenars de metres sota la superfície. En aquest cas, la pressió que ha de suportar la seva estructura pot arribar a ser immensa. Un altre cas bastant intuïtiu és el d'una presa, en un pantà. Com que la part inferior de la presa ha de suportar una pressió molt més gran que la part superior, les preses solen construir-se molt més amples a la part de baix que a la part de dalt. 

\section{Principi de Pascal}

El Principi de Pascal és una llei aparentment senzilla, però que pot donar lloc a fenòmens bastant curiosos. A més, té una gran quantitat d'aplicacions, tant en tecnologia com en la vida de molts organismes, incloent-hi els esquelets d'alguns artròpodes. L'enunciat del Principi és el següent:

\begin{mdframed}
\begin{it}
    La pressió aplicada en un fluid ideal, homogeni, incompressible, en un recipient tancat, es ``transmet'' amb igual intensitat en totes direccions a qualsevol punt, incloent les parets del recipient.
\end{it}
\end{mdframed}
%
\begin{figure}[H]
    \centering
    \includegraphics[width=0.5\linewidth]{Figures/PrincipiDePascal.pdf}
    \caption{Si exercim una força $F$ sobre l'èmbol del cilindre de secció $S$, la pressió de l'oli a una profunditat $\Delta h$ sota l'èmbol serà $P = F/S + \rho g \Delta h$.}
    \label{fig:PrincipiDePascal}
\end{figure}
%
Això vol dir que si exercim una força $F$ amb alguna mena d'èmbol de superfície $S$ com el de la Figura \ref{fig:PrincipiDePascal}, tots els punts del fluid que toquen l'èmbol passaran a tenir una pressió hidroestàtica $P_e = F/S$. I no només això, tots els altres punts del fluid tindran una pressió donada per $P = P_e + \rho g \Delta h$, en el cas que tinguin una diferència d'altura $\Delta h$ respecte a l'èmbol. En el cas que la gravetat $g$ sigui negligible respecte a $P_e$, podem considerar que la pressió és $P_e$ en tot el fluid. Veurem ara alguns exemples que mostren aplicacions d'aquest principi.

\subsection{L'experiment de Torricelli (1643)}

L'any 1642, Evangelista Torricelli va fer el primer experiment destinat a mesurar la pressió de l'aire. Va construir un aparell senzill que actualment es coneix com a baròmetre de mercuri, i que encara s'utilitza en algunes ocasions. Es diu ``de mercuri'' perquè fa servir aquest metall, líquid a temperatura ambient, però amb una densitat molt alta. Això fa que, d'acord amb l'equació (\ref{eqn:PressioHidroestatica}), la pressió augmenti molt ràpidament amb les diferències d'altura.

L'experiment consisteix en omplir de mercuri fins dalt un tub de vidre obert per un extrem, invertir-lo amb l'extrem tapat, sobmergir-lo dins una cubeta amb més mercuri, i finalment destapar-lo. El muntatge es pot veure a la Figura \ref{fig:ExperimentDeTorricelli}. 
%
\begin{figure}[H]
    \centering
    \includegraphics[width=0.4\linewidth]{Figures/ExperimentDeTorricelli.pdf}
    \caption{En l'experiment de Torricelli, el nivell de mercuri dins el tub es manté a una altura de 760 mm per sobre de la cubeta a causa de la pressió atmosfèrica.}
    \label{fig:ExperimentDeTorricelli}
\end{figure}
%
Un cop destapada la part inferior del tub, el nivell del mercuri baixarà a causa del seu pes, fins que quedi equilibrat amb la pressió que l'aire fa en sentit contrari. D'aquesta manera, l'altura final del mercuri correspondrà a la pressió atmosfèrica. Podem aplicar l'equació (\ref{eqn:PressioHidroestatica}) per veure quina serà aquesta altura. En primer lloc, sabem que a la part superior del tub la pressió és 0, perquè l'aire no hi ha pogut entrar i per tant hi tenim el buit. A la superfície del mercuri de la cubeta, la pressió és l'atmosfèrica $P_\text{atm}$, perquè està en contacte amb l'aire. Llavors, tenim
%
\begin{equation}
 P_\text{atm} - 0 = \rho g \Delta h \, .   
\end{equation}
%
Torricelli va obtenir que $\Delta h$ = 760 mm. Per aquesta raó, definim la unitat dels ``mil·límetres de mercuri'' da tal manera que 760 mmHg corresponen a 1 atmosfera de pressió. Quina altura hagués obtingut Torricelli si hagués utilitzat aigua enlloc de mercuri?

\subsection{L'experiment de Pascal (1646)}

%
\begin{figure}[H]
    \centering
    \includegraphics[width=0.33\linewidth]{Figures/ExperimentDePascal.pdf}
    \caption{En l'experiment de Pascal, la gran altura assolida per l'aigua en el tub produeix una gran pressió dins el dipòsit inferior, que el pot arribar a fer esclatar.}
    \label{fig:ExperimentDePascal}
\end{figure}
%
Tot i que la veracitat històrica d'aquest experiment no està totalment clara, els principis en els que es basa són perfectament reals, i de fet l'experiment es pot reproduir en l'actualitat. Es diu que el 1646, Blaise Pascal va decidir demostrar el seu principi de la pressió hidrostàtica fent un experiment altament espectacular, i força perillós. L'experiment es basa en el fet que l'increment de pressió hidrostàtica en un recipient només depèn de la profunditat, i no de la quantitat de fluid que hi hagi. Per tant, si fem servit un tub vertical molt estret, podem en principi aconseguir pressions molt altes amb una quantitat moderada de fluid. 

En el seu muntatge, Pascal suposadament va connectar un tub molt prim de molts metres d'altura a la part superior d'una bóta de vi plena d'aigua fins dalt, com es veu a la Figura \ref{fig:ExperimentDePascal}. Llavors, amb l'ajuda d'una escala, va afegir el contingut d'una gerra d'aigua per omplir el el tub. Això naturalment resultà en un augment molt gran del nivell de l'aigua, de manera que la pressió dins la bóta va ser prou gran per fer-la esclatar. Podeu calcular que, per un tub d'uns 10 metres d'altura i un barril d'un metre de diàmetre, la força que ha se suportar cada tapa del barril és l'equivalent a gairebé 8 tones de pes.

\subsection{La premsa hidràulica}

\begin{figure}[H]
    \centering
    \includegraphics[width=0.65\linewidth]{Figures/PremsaHidraulica.pdf}
    \caption{Un model senzill de premsa hidràulica. El Principi de Pascal ens permet multiplicar la força fent servir èmbols d'àrees diferents, ja que la pressió és la mateixa en tots els punts de l'oli.}
    \label{fig:PremsaHidraulica}
\end{figure}
%
Una aplicació tecnològica molt important de la llei de Pascal és la possibilitat de ``multiplicar'' la força gràcies a l'ús de dos èmbols amb àrees diferents connectats a la mateixa cambra de pressió, com es veu a la Figura \ref{fig:PremsaHidraulica}. Si exercim una força $F_1$ sobre l'èmbol d'àrea $S_1$, això es traduirà en una força $F_2$ sobre l'èmbol d'àrea $S_2$. Si recordem el Principi de Pascal, la pressió sobre els dos èmbols ha de ser la mateixa, és a dir
%
\begin{equation}
    \frac{F_1}{S_1} = \frac{F_2}{S_2}\, ,
\end{equation}
%
i per tant
%
\begin{equation}
    F_2 = \left(\frac{S_2}{S_1}\right) F_1\, .
\end{equation}
%
Si fem que $S_2$ sigui molt més gran que $S_1$, per exemple $S_2/S_1 = 10$, multiplicarem la força per 10. Aquest és el principi de la premsa hidràulica, que es fa servir en tota mena d'aplicacions, des de la metal·lúrgia fins al sistema de frens dels vehicles.  

\subsection{La llei dels vasos comunicants}

Finalment, una conseqüència curiosa del Principi de Pascal és la llei dels vasos comunicants. Aquests són precisament el que el nom indica, una sèrie de recipients oberts que comuniquen entre ells a través d'una conducció per la seva part inferior, com es veu a la Figura \ref{fig:VasosComunicants}.
%
\begin{figure}[H]
    \centering
    \includegraphics[width=0.8\linewidth]{Figures/VasosComunicants.pdf}
    \caption{Els vasos comunicants s'omplen tots fins al mateix nivell, per tal que a pressió a la seva base sigui la mateixa. Això pot resultar paradoxal, ja que el pes del líquid en cada vas pot ser molt diferent a la dels altres.}
    \label{fig:VasosComunicants}
\end{figure}
%
Com que la pressió només depèn de l'altura, tots els recipients s'omplen fins al mateix nivell, corresponent a la pressió atmosfèrica. Això passa independentment de la forma que tinguin els recipients, encara que el pes del líquid que conté cada un d'ells pugui ser molt diferent. Això es fa servir tradicionalment per moltes aplicacions, per exemple per mesurar el nivell del vi en una bota.

\section{Principi d’Arquimedes}

Fins ara hem vist que la pressió hidroestàtica d'un fluid exerceix una força sobre les parets del recipient, i sobre altres parts adjacents del fluid. El Principi d'Arquimedes es basa en l'efecte que la pressió hidroestàtica té sobre els cossos que es troben submergits en aquest fluid. És el que permet el fenomen de la flotació en l'aigua, i fins i tot en l'aire. A la Figura \ref{fig:PrincipiArquimedes} hem delimitat un volum $V$, de forma arbitrària, en l'interior d'un fluid. La pressió hidroestàtica exercirà tot un conjunt de forces en tota la seva superfície, la suma de les quals anomenarem empenyiment, i simbolitzarem com a $\vec E$. Calcular el valor de $\vec E$ directament requereix l'ús d'eines matemàtiques complicades, de manera que la trobarem a través d'un argument d'equilibri.
%
\begin{figure}[H]
    \centering
    \includegraphics[width=0.5\linewidth]{Figures/PrincipiArquimedes.pdf}
    \caption{Tot cos submergit total o parcialment en un fluid experimenta una força vertical cap amunt (empenyiment d'Arquimedes) igual al pes del fluid desallotjat.}
    \label{fig:PrincipiArquimedes}
\end{figure}
%
Si el volum $V$ estigués ple del mateix fluid que la resta del recipient, les forces haurien d'estar equilibrades per tal de mantenir-lo en repòs. En altres paraules, l'empenyiment $\vec E$ ha de tenir igual magnitud i sentit contrari que el pes del fluid en $V$, és a dir
%
\begin{equation}
    E = m_f g = V \rho_\text{f} g \, ,
\end{equation}
%
on $\rho_\text{f}$ és la densitat del fluid. Ara, què passaria si el volum $V$ no estigués ple de fluid, sinó que fos un objecte submergit, per exemple una pedra? El fluid del seu voltant seria el mateix, i amb la mateixa pressió hidroestàtica, i per tant la força d'empenyiment seria la mateixa! Això ens permet enunciar el famós Principi d'Arquimedes:

\begin{mdframed}
\begin{it}
    Tot cos submergit total o parcialment en un fluid experimenta una força vertical cap amunt (empenyiment d'Arquimedes) igual al pes del fluid desallotjat.
\end{it}
\end{mdframed}

Observeu que, si el cos que hem submergit és més dens que el fluid (més ``pesant''), l'empenyiment no serà suficient per compensar el pes i l'objecte s'enfonsarà. Només s'arribarà a l'equilibri quan toqui el fons del recipient. Aquest seria el cas en llançar una pedra al riu.

En canvi, si la densitat de l'objecte és menor (és més ``lleuger'' que el fluid), l'empenyiment serà major que el pes, i l'objecte flotarà. En aquest cas, s'arribarà a l'equilibri quan una part de l'objecte sobresurti del fluid. En aquest cas, el fluid desallotjat per la part submergida tindrà la mateixa massa que el propi objecte. Per exemple, això és el que passa amb un vaixell, o un iceberg (recordeu l'anomalia de l'aigua?). Fins i tot en l'aire, amb els globus aerostàtics i amb els globus d'heli de les festes.

En el món dels éssers vius, la flotació és fonamental en una gran quantitat d'organismes nedadors (meduses, peixos, mol·luscs...), sovint a través d'un òrgan molt enginyós conegut com a bufeta natatòria. Els animals nedadors, com els peixos,  necessiten controlar la seva altura dins de l'aigua i poder enfonsar-se i emergir a voluntat. Per tant, necessiten poder canviar la seva densitat mitjana sempre que vulguin. Atès que no poden canviar fàcilment la seva massa (agafar o deixar anar llast), el mètode que utilitzen és canviar el seu volum total. Això ho fan a través d'un petit ``globus'' que tenen al seu interior, anomenat bufeta natatòria, que poden omplir i buidar de gasos, majoritàriament oxigen i nitrogen. Per pujar, l'omplen de gasos, i per baixar la buiden. Aquest mètode també és utilitzat pels éssers humans en els submarins.

\chapter{Dinàmica de fluids ideals}
\label{chap:Hidrodinamica}

\section{Camps de velocitat}

Un cop hem estudiat la física bàsica dels fluids ideals en repòs, és a dir aquells que no es mouen, podem passar a estudiar els fluids en moviment. Això ens permetrà explicar, per exemple, la circulació de la sang dins les venes, artèries i capil·lars; i fins i tot els principis que permeten el vol de les aus. 

En el cas dels fluids en repòs, podíem descriure la seva física en termes d'un camp de pressió $P(\vec r)$. Per poder tractar amb la dinàmica de fluids, o hidrodinàmica, tindrem també un camp de velocitats: $\vec v(\vec r, t)$, $P(\vec r, t)$. Tot i això, en aquest curs ens restringirem al cas en què aquest camp de velocitats compleixi certes propietats que faran més fàcil el seu tractament. En primer lloc, seguirem considerant que el fluid és incompressible i no viscós, és a dir que és un fluid ideal. Considerarem també que el flux és estacionari, laminar i irrotacional.

\begin{itemize}
    \item Flux estacionari: La velocitat i la pressió no depenen del temps: $P(\vec r)$, $\vec v(\vec r)$.
    \item Flux laminar: Les capes del fluid avancen paral·leles, sense barrejar-se entre elles. És a dir, no hi ha turbulència.
    \item Flux irrotacional: El fluid no dona voltes. És a dir, no té cap tipus de remolí o vòrtex.
\end{itemize}

En aquestes condicions, serà útil treballar amb els conceptes de línia de flux i tub de flux. 
%
\begin{figure}[H]
    \centering
    \includegraphics[width=0.6\linewidth]{Figures/TubDeCorrent.pdf}
    \caption{Representació d'un tub de flux, limitat per una família de línies de flux.}
    \label{fig:TubDeCorrent}
\end{figure}
%
\begin{itemize}
    \item Una línia de flux és simplement la trajectòria que seguiria una partícula, per exemple un bacteri, dins del fluid. Correspon a la línia imaginària que resulta d'unir els vectors velocitat l'un darrere l'altre. Si el flux és estacionari, aquestes línies de flux no canviaran en el temps.
    \item Un tub de flux és un conjunt de línies de flux que envolten una superfície imaginària en el fluid, tal com es veu a la Figura \ref{fig:TubDeCorrent}. En un flux estacionari actuarà com si fos un tub tancat ja que, per definició, el fluid no travessarà les seves parets.
\end{itemize} 

\section{Equació de continuïtat}

Un concepte fonamental en la hidrodinàmica és que la quantitat de fluid es conserva. No es pot generar ni destruir matèria. D'aquesta manera, que en règim estacionari, la massa que entra en una regió de l'espai durant un cert temps ha de sortir per algun altre lloc. Si el fluid és, a més, incompressible, no només la massa s'ha de conservar sinó també el seu volum. La manera més senzilla de visualitzar-ho és pensar en una canonada, com es veu a la Figura \ref{fig:Continuitat}.
%
\begin{figure}[H]
    \centering
    \includegraphics[width=0.9\linewidth]{Figures/Continuitat1.pdf}
    \includegraphics[width=0.9\linewidth]{Figures/Continuitat2.pdf}
    \caption{L'equació de continuïtat ens diu que, en un flux estacionari, la quantitat de fluid dins una secció del tub serà sempre la mateixa. Per això, el fluid que entri per l'esquerra del tub durant un temps $\Delta t$ tindrà la mateixa massa que el que en surti per la dreta.}
    \label{fig:Continuitat}
\end{figure}
%
En aquesta situació, durant un temps $\Delta t$, entrarà per l'esquerra un volum de fluid $\Delta V_1 = S_1 v_1 \Delta t$, on $S_1$ és l'àrea de la secció de l'entrada del tub. D'igual manera, per la dreta del tub  sortirà un volum $\Delta V_2 = S_2 v_2 \Delta t$. Així, podem definir el concepte de cabal ($Q$) com el volum de fluid que passa per un tub per unitat de temps
%
\begin{equation}
    \boxed{Q \equiv \frac{\Delta V}{\Delta t} = S v}
    \label{eqn:Cabal}
\end{equation}
%
De manera similar, el flux ($\Phi$) es defineix com la massa de fluid que passa per un tub per unitat de temps
%
\begin{equation}
    \boxed{\Phi \equiv \frac{\Delta m}{\Delta t} = \rho Q = \rho S v}
    \label{eqn:Flux}
\end{equation}
%
Com hem dit, en un flux estacionari, tant la massa que entra al tub ha de ser idèntica a la que en surt. Per tant, el flux es conserva al llarg de la canonada. Això s'expressa a través de l'equació de continuïtat
%
\begin{equation}
    \boxed{\Phi = \text{const.}} \quad \text{(per fluxos estacionaris)}
    \label{eqn:ContinuitatPhi}
\end{equation}
%
En el cas que el fluid sigui incompressible, el volum entrant i el sortint seran iguals, i per tant es conservarà també el cabal
%
\begin{equation}
    \boxed{Q = \text{const.}} \quad \text{(per fluxos estacionaris incompressibles)}
    \label{eqn:ContinuitatQ}
\end{equation}
%
Una de les conseqüències més importants de l'equació de continuïtat és que la velocitat del fluid en un tub de corrent augmenta quan disminueix l'àrea de la seva secció, i viceversa. Per exemple, el raig de l'aixeta de la cuina es fa més estret a mesura que l'aigua guanya velocitat a causa de la gravetat. En el cas dels vasos sanguinis, la situació pot semblar contradictòria, ja que la velocitat de la sang en els capil·lars és molt més baixa que en les grans venes i artèries. Això passa perquè l'àrea conjunta (sumada) de tots els capil·lars pels quals la sang pot abastir el nostre cos és molt més gran que la dels grans vasos com l'artèria aorta o la vena cava.

\section{Equació de Bernoulli}

Si bé la conservació de la massa ens ha donat l'equació de continuïtat, la conservació de l'energia ens donarà una altra equació molt important: l'equació de Bernoulli. Aquesta llei es complirà només en el cas de fluids no viscosos, laminars i irrotacionals, en els que no hi ha fregament entre les capes del fluid, i per tant no es perd energia en forma de calor. Per entendre l'equació de Bernoulli, hem de recordar que el treball per forces externes a un sistema mecànic serà igual al canvi en la seva energia mecànica (cinètica i potencial). Escrit matemàticament
%
\begin{equation}
    W = \Delta K + \Delta U\, .
    \label{eqn:EnergiaTreball1}
\end{equation}
%
El treball exercit per una força es defineix com el producte de la força pel desplaçament del seu punt d'aplicació en la direcció de la força. L'energia cinètica d'una massa de fluid es defineix com a $K = \frac12 m v^2$, i l'energia potencial (en aquest cas gravitatòria), com a $U = m g h$, on $h$ és l'alçada del fluid. Anem a aplicar aquests conceptes a la Figura \ref{fig:Bernoulli}.
%
\begin{figure}[H]
    \centering
    \includegraphics[width=0.9\linewidth]{Figures/Bernoulli.pdf}
    \caption{L'equació de Bernoulli és equivalent al principi de conservació de l'energia mecànica per fluids no viscosos, laminars i irrotacionals. Ens diu que la quantitat $P + \frac12 \rho v^2 + \rho g h$ es conserva al llarg de cada tub de flux.}
    \label{fig:Bernoulli}
\end{figure}
%
El cilindre de fluid que entra dins el tub per l'esquerra durant el temps $\Delta t$ tindrà una massa total donada per $\rho S_1 v_1 \Delta t$. La seva energia cinètica serà, $K_1 = \frac12 \rho_1 S_1 v_1^3 \Delta t $, i la potencial $U_1 = \rho_1 S_1 v_1 g h_1 \Delta t$. D'igual manera, el cilindre de fluid que surt per la dreta tindrà $K_2 = \frac12 \rho_2 S_2 v_2^3 \Delta t $, i la potencial $U_2 = \rho_2 S_2 v_2 g h_2 \Delta t$. El treball exercit per la pressió durant aquest temps $\Delta t$ serà $W = P_1 S_1 v_1 \Delta t - P_2 S_2 v_2 \Delta t$. Sabent que la variació d'energia mecànica en el temps $\Delta t$ serà l'energia afegida pel fluid que entra, menys l'energia que s'emporta el fluid que surt, tenim
%
\begin{equation}
    W = K_2 + U_2 - K_1 - U_2 \, .
    \label{eqn:EnergiaTreball2}
\end{equation}
%
Això, reescrit en termes de les variables del fluid, queda
%
\begin{equation}
    \left( P_1 S_1 v_1 - P_2 S_2 v_2 \right) \Delta t = \frac12 \left( \rho_2 S_2 v_2^3 - \rho_1 S_1 v_1^3 \right) \Delta t +  g \left( \rho_2 S_2 v_2 h_2 - \rho_1 S_1 v_1 h_1 \right) \Delta t\, .
\end{equation}
%
Finalment, si recordem l'equació de continuïtat (\ref{eqn:ContinuitatPhi}), $\rho_1 S_1 v_1 = \rho_2 S_2 v_2$, fent una mica d'àlgebra arribem a 
%
\begin{equation}
    P_1 + \frac12 \rho_1 v_1^2 + \rho_1 g h_1 = P_2 + \frac12 \rho_2 v_2^2 + \rho_2 g h_2\, .
\end{equation}
%
Si ens hi fixem, això ens dóna una altra equació de conservació, anomenada equació de Bernoulli
%
\begin{equation}
    \boxed{P + \frac12 \rho v^2 + \rho g h = \text{const.}}
    \label{eqn:Bernoulli}
\end{equation}
%
Aquesta equació es compleix per tots els punts dins un tub de corrent estacionari, no viscós i irrotacional. Serà de gran utilitat en l'estudi dels fluids en moviment. A continuació, anem a veure algunes de les aplicacions més usuals.

\subsection{Efecte Venturi}

%
\begin{figure}[H]
    \centering
    \includegraphics[width=0.8\linewidth]{Figures/Venturi.pdf}
    \caption{L'efecte Venturi és una aplicació senzilla de l'equació de Bernoulli. Si hi ha un estretament en un tub, la velocitat del fluid augmentarà ($v_2 > v_1$) i per tant la pressió disminuirà ($P_2 < P_1$).}
    \label{fig:Venturi}
\end{figure}
%
A la Figura \ref{fig:Venturi} es pot veure un fragment de tub amb un petit estretament a la part central. A causa de l'equació de continuïtat, la velocitat del fluid serà més alta a la part estreta que a la part gruixuda. Si combinem això amb l'equació de Bernoulli, tenint en compte que el tub està a una altura constant, veiem que la pressió a la part estreta és més baixa que a la part gruixuda. Concretament, tenim
%
\begin{equation}
    P_1 - P_2 = \frac12 \rho \left(v_2^2 - v_1^2 \right) = \frac12 \rho v_1^2 \left[\left(\frac{S_1}{S_2}\right)^2 - 1 \right] \, .
    \label{eqn:Venturi}
\end{equation}
%
Aquest tipus de muntatge es fa servir en diversos dispositius. Per exemple en els tubs de Venturi, on es mesura la diferència de pressió entre dos trams de tub amb diàmetres diferents, per tal de determinar el cabal del fluid. També s'utilitza en polvoritzadors de perfums, i en els carburadors de motors d'explosió. L'efecte Venturi es feia servir tradicionalment en la farga hidràulica catalana, per comprimir l'aire que després seria utilitzat per avivar el forn on s'escalfava el ferro. D'aquesta manera s'evitava l'ús de manxes, que presentaven un alt risc d'incendi. 

\subsection{Efecte Magnus}

%
\begin{figure}[H]
    \centering
    \includegraphics[width=0.49\linewidth]{Figures/EfecteMagnus.pdf}
    \caption{Quan una pilota es llança en rotació, la velocitat de l'aire respecte a la superfície de la pilota és diferent entre un costat i l'altre. Això provoca una diferència de pressió que en desvia la trajectòria. A la figura, el color vermell indica altes pressions, mentre que el blau correspon a baixes pressions. }
    \label{fig:EfecteMagnus}
\end{figure}
%
Un altre cas típic on l'equació de Bernoulli juga un paper important és en l'efecte Magnus, molt conegut en els esports de pilota simplement com a ``efecte''. Això passa quan un objecte esfèric, com ara una pilota de futbol o tenis, gira molt ràpidament mentre es desplaça per l'aire. Tal com es mostra a la Figura \ref{fig:EfecteMagnus}, la pilota ``arrossega'' part de l'aire en el moviment de rotació al seu voltant. Això fa que la pressió en un dels seus costats, on l'aire es mou més ràpidament respecte de la pilota, tingui una pressió menor. D'aquesta manera, la pilota pateix una força aerodinàmica que la fa seguir una trajectòria corba. Un efecte molt semblant a aquest és el que permet les forces de sustentació en una ala, i per tant el vol dels ocells i els avions.


\subsection{Llei de Torricelli}

%
\begin{figure}[H]
    \centering
    \includegraphics[width=0.6\linewidth]{Figures/LleiDeTorricelli.pdf}
    \caption{Fent servir l'equació de Bernoulli podem obtenir la llei de Torricelli, que ens diu la velocitat del líquid que s'escapa per un forat del dipòsit.}
    \label{fig:LleiDeTorricelli}
\end{figure}
%
Podem també aplicar l'equació de Bernoulli per saber la velocitat a la qual un líquid s'escapa per un forat a la paret d'un dipòsit, com es veu a la Figura \ref{fig:LleiDeTorricelli}. En aquest cas, cal escollir amb cura el tub de flux on la volem aplicar. Agafarem un tub de flux que vagi des de la superfície del líquid fins a la sortida pel forat. Considerarem que el forat és molt més petit que el dipòsit, de manera que el fluid a la seva superfície pràcticament no es mou. Així, tenim
%
\begin{equation}
    P_\text{atm} + 0 + \rho g h = P_\text{atm} + \frac12 \rho v^2 + 0 \, .
    \label{eqn:LLeiTorricelli1}
\end{equation}
%
Aïllant la velocitat de sortida $v$ i simplificant, ens queda
%
\begin{equation}
    v = \sqrt{2 g h} \, .
    \label{eqn:LLeiTorricelli2}
\end{equation}
%
Fixeu-vos que aquesta és exactament la mateixa velocitat que tindria el líquid si hagués caigut des del repòs una alçada $h$ en caiguda lliure.

\section{Forces de sustentació}

\subsection{Efecte Coandă}

Per acabar aquest capítol, veurem els conceptes bàsics que permeten un dels efectes més espectaculars de la dinàmica de fluids: Les forces de sustentació i el vol, tant dels ocells com dels avions i helicòpters. Una peculiaritat dels fluids que fan servir tant els animals com els aparells que es desplacen per l'aire i l'aigua és que els fluids tenen tendència a ser atrets i a ``resseguir'' les superfícies que es troben al seu pas. Ho podeu veure esquematitzat a la Figura \ref{fig:EfecteCoanda}. Això es coneix com a efecte Coandă, i permet dirigir les línies de flux de la manera que ens convingui a través de superfícies aerodinàmiques (o hidrodinàmiques) com ales, alerons i timons. En la cultura popular catalana, l'efecte Coandă s'utilitza per a la tradició de l'ou com balla, en la qual un ou es manté suspès sobre el raig d'aigua d'una font. 
%
\begin{figure}[H]
    \centering
    \includegraphics[width=0.49\linewidth]{Figures/EfecteCoanda.pdf}
    \caption{Les línies de flux tendeixen a seguir el contorn dels objectes que es troben al seu pas. Aquest fenomen s'anomena efecte Coandă.}
    \label{fig:EfecteCoanda}
\end{figure}
%

\subsection{L'aerodinàmica i el vol}

Tant les ales dels ocells com les dels avions, juntament amb els rotors dels helicòpters, i fins i tot les veles dels vaixells, es basen en un perfil aerodinàmic com el que es mostra a la Figura \ref{fig:Sustentacio}. Les superfícies aerodinàmiques amb aquest perfil, la secció del qual té amb una forma similar a la d'una ametlla, experimenten forces quan es desplacen a través de l'aire. Aquestes forces es classifiquen en dues categories:  sustentació i resistència (o arrossegament).
%
\begin{figure}[H]
    \centering
    \includegraphics[width=0.6\linewidth]{Figures/Sustentacio.pdf}
    \caption{Simulació del flux de l'aire al voltant del perfil de l'ala d'un avió. Les diferències de pressió produeixen una força d'arrossegament cap enrere, i una força de sustentació cap amunt que ens permet volar. A la figura, el color vermell indica altes pressions, mentre que el blau correspon a baixes pressions.}
    \label{fig:Sustentacio}
\end{figure}
%
\begin{itemize}
    \item Força de sustentació: Aquesta força és vertical i cap amunt, a causa de la diferència de pressions entre les parts inferior i superior de l'ala. És la força que manté els avions suspesos en l'aire.
    \item Força d'arrossegament: Aquesta és la força d'arrossegament que s'oposa al moviment de l'ala. En el cas d'un avió, està orientada horitzontalment cap enrere, i és la raó per la qual els avions necessiten mantenir els motors en marxa mentre volen.
\end{itemize}
%
Podem fer una anàlisi aproximada de l'origen de la força de sustentació a través de l'equació de Bernoulli, negligint efectes gravitatoris, tal i com hem fet en el cas de l'efecte Venturi a l'equació (\ref{eqn:Venturi}). En aquest cas, la diferència entre les pressions a la part inferior $P_i$ i a la part superior $P_s$ de l'ala es pot obtenir a través de les velocitats de l'aire a les dues regions, $v_i$ i $v_s$, respectivament
%
\begin{equation}
    P_i - P_s = \frac12 \rho \left(v_s^2 - v_i^2 \right) \, .
    \label{eqn:Sustentacio1}
\end{equation}
%
Si considerem que les velocitats sota i sobre l'ala es poden expressar com el producte de dos factors d'escala per la velocitat de l'avió $v$, és a dir $v_i = c_i v$ i $v_s = c_s v$, tenim
%
\begin{equation}
    P_i - P_s = \frac12 \rho v^2 \left( c_s^2 - c_i^2 \right) = \frac12 \rho v^2 C_s\, ,
    \label{eqn:Sustentacio2}
\end{equation}
%
on hem definit el coeficient de sustentació com $C_s \equiv c_s^2 - c_i^2$. Finalment, recordant que la força total és la diferència de pressions multiplicada per la superfície de l'ala $S_\text{ala}$, obtenim l'expressió per a la força de sustentació $F_s$ com 
%
\begin{equation}
    \boxed{F_s = \frac12 \rho \, C_s \,  S_\text{ala} \, v^2 }
    \label{eqn:Sustentacio3}
\end{equation}
%
Podem ara observar aquesta equació i veure l'origen d'alguns efectes molt utilitzats en l'aviació, tot i que apliquen exactament igual en el cas dels ocells. Per tal de poder volar, la força de sustentació generada per les ales s'oposarà al pes de l'avió, evitant d'aquesta manera que l'aparell caigui cap al terra. Si la sustentació és una mica superior al pes, l'avió guanyarà altura, mentre que si la sustentació és inferior al pes l'avió baixarà. 

Com major sigui l'altura a la que es troba l'avió, menor serà la densitat $\rho$ de l'aire, ja que aquesta disminueix amb l'alçada. Per tant, a través de l'equació (\ref{eqn:Sustentacio3}) podem veure que l'avió necessitarà desplaçar-se a més velocitat $v$ per tal de mantenir la mateixa força de sustentació $F_s$. D'aquesta manera, els avions volen molt més ràpid en altitud de creuer que prop del terra, cosa que és molt convenient si volem arribar aviat a la destinació. Per tal d'enlairar-se i aterrar amb seguretat, cal reduir la velocitat el màxim possible. Per això, els avions poden augmentar la superfície $S_\text{ala}$ de les ales desplegant unes extensions anomenades ``flaps'' i ``slats'' en anglès, que permeten mantenir-se en vol a velocitats més baixes.

\subsection{Angle d'atac}

%
\begin{figure}[H]
    \centering
    \includegraphics[width=0.49\linewidth]{Figures/AngleAtac10.pdf}
    \includegraphics[width=0.49\linewidth]{Figures/AngleAtac30.pdf}
    \caption{Perfils aerodinàmics amb angles d'atac de 10$^\circ$ (esquerra) i 30$^\circ$ (dreta). Novament, el color vermell indica altes pressions, mentre que el blau correspon a baixes pressions. Podem veure que la força de sustentació serà més gran per a major angle d'atac.}
    \label{fig:AngleAtac}
\end{figure}
%
També, és possible guanyar i perdre altura variant l'angle del cos de l'avió respecte a la direcció del seu desplaçament, el que es coneix amb el nom de capcineig, o pujar i baixar el nas de l'avió. L'angle que forma el perfil de les ales amb la direcció del vent es coneix com a ``angle d'atac'', i està exemplificat a la Figura \ref{fig:AngleAtac}. Un angle d'atac massa gran cap amunt farà guanyar altura però reduirà molt la velocitat. Això pot provocar que l'avió ``entri en pèrdua'', és a dir, que perdi momentàniament la seva sustentació i es desplomi. Naturalment, entrar en pèrdua prop del terra és extremadament perillós i cal evitar-ho.  

\chapter{Fluids viscosos}
\label{chap:Viscositat}

\section{Viscositat}

%
\begin{figure}[H]
    \centering
    \includegraphics[width=0.6\linewidth]{Figures/Viscositat.pdf}
    \caption{A diferència dels fluids ideals, els fluids viscosos presenten una força de fregament que s'oposa al desplaçament tangencial.}
    \label{fig:Viscositat}
\end{figure}
%
Els fluids ideals que hem descrit fins ara no presenten cap tipus de resistència a les forces tangencials, i per tant flueixen sense cap tipus de fregament ni pèrdua d'energia. Per exemple, si remenéssim una galleda plena de fluid ideal fins que es posés a girar, el fluid seguiria girant per sempre més. Naturalment, la nostra experiència quotidiana ens diu que els fluids real no són així. Si bé alguns fluids es poden aproximar prou bé com a fluids ideals a escales de temps curtes, com l'aigua i l'aire, d'altres tenen clarament un comportament molt diferent, com per exemple la mel o el sabó líquid. Aquests últims presenten molt clarament una resistència a fluir, i s'aturen ràpidament si els posem a girar dins una galleda. Aquesta resistència a fluir es pot descriure força bé com una força de fregament que s'oposa al desplaçament tangencial de les capes de fluid, un fenomen que anomenarem força viscosa. 

A la Figura \ref{fig:Viscositat} es representa un fragment de fluid atrapat entre dues planxes metàl·liques que es desplacen lateralment. El resultat és que les diverses capes de fluid es desplacen a velocitat diferent, on les capes que estan en contacte directe amb les planxes s'hi mantenen enganxades. Esperem que el fluid s'oposi a aquesta deformació, produint una força viscosa en sentit oposat, que anomenarem $F_\eta$. Intuïtivament, aquesta força serà proporcional a la superfície $S$ de les capes de fluid involucrades, i també serà proporcional al ritme de variació de la velocitat de les capes respecte a la seva distància $dv/dz$. Aquesta última derivada la serà molt gran si les planxes estan molt juntes o es mouen a velocitats molt diferents, però serà petita si la diferència de velocitats és petita o les planxes estan molt separades. Matemàticament, definim el coeficient de viscositat $\eta$ com la constant de proporcionalitat
%
\begin{equation}
    \boxed{F_\eta = - \eta \, S \, \frac{dv}{dz} }
    \label{eqn:Viscositat}
\end{equation}
%
on el signe negatiu indica que la força s'oposa al moviment del fluid. El coeficient de viscositat $\eta$ és una nova propietat que depèn del fluid, i es mesura en unitats de Pa·s (Pascals segon) en el SI. També és comú trobar-la expressada en g/cm·s, una unitat que rep el nom de Poise i que correspon a 0.1 Pa·s. En general, la viscositat dels fluids és més gran que la dels gasos. Un fet bastant curiós, relacionat amb l'origen microscòpic de les forces viscoses, és que la viscositat dels líquids sol disminuir amb la temperatura mentre que la dels gasos augmenta. En el cas dels líquids, es pot comprovar fàcilment posant un pot de mel al microones per tornar-la menys viscosa. A la Taula \ref{tab:Viscositats} es mostren els valors de viscositat d'algunes substàncies quotidianes.

\begin{table}[H]
    \centering
    \begin{tabular}{|l|c|}
    \hline
                       & $\eta$ \\ \hline
        Aire sec       & 18,3 $\cdot$ 10$^{-6}$ Pa·s \\ \hline
        Aigua (a 20ºC) & 1,00 $\cdot$ 10$^{-3}$ Pa·s \\ \hline
        Llet           & 2,12 $\cdot$ 10$^{-3}$ Pa·s \\ \hline
        Oli d'oliva    & 56,2 $\cdot$ 10$^{-3}$ Pa·s \\ \hline
        Mel            & 2 - 10 Pa·s \\ \hline
    \end{tabular}
    \caption{Viscositats d'alguns fluids habituals.}
    \label{tab:Viscositats}
\end{table}


\section{Flux viscós en un tub}

Si intentem fer circular un fluid viscós a través d'un tub, això ens costarà una certa energia, i necessitarem mantenir una diferència de pressió entre els seus extrems per tal de mantenir el flux. Aquest és el cas, per exemple, de la sang dins els nostres vasos sanguinis. A causa de la seva viscositat, necessitem que el cor bategui constantment per aplicar una pressió que manté la sang en circulació. Si el cor s'atura, la viscositat de la sang en frena ràpidament la circulació i deixa de subministrar oxigen i nutrients a les nostres cèl·lules, cosa que impedeix mantenir el funcionament del nostre cos. 
%
\begin{figure}[H]
    \centering
    \includegraphics[width=0.9\linewidth]{Figures/FluxEnTub.pdf}
    \caption{Divisió del fluid dins una canonada en cilindres concèntrics, per poder deduir el perfil de velocitats. A causa de la viscositat, necessitarem una diferència de pressió $\Delta P$ entre els extrems del tub per tal de mantenir el fluid en moviment.}
    \label{fig:FluxEnTub}
\end{figure}
%
Per tal d'entendre com afecta la viscositat a la circulació de fluid per un tram de tub de llargada $L$ i radi $R$, sota una diferència de pressió $\Delta P$, anem a aplicar l'equació (\ref{eqn:Viscositat}). Intuïtivament, podem esperar que el fluid estigui aturat a les parets del tub, que estan fixes, i es mogui a màxima velocitat al llarg del seu eix. Per fer el càlcul, descomponem el fluid dins el tram de tub en una sèrie de cilindres coaxials, com es veu a la Figura \ref{fig:FluxEnTub}. Aquest tipus de descomposició ens serà útil al llarg d'aquest capítol. Suposarem que a l'esquerra del tub tenim una pressió $P + \Delta P$, i que a la dreta tenim només $P$, de manera que el flux es veu impulsat cap a la dreta.

Per tal de mantenir un flux estacionari, per a cada cilindre de fluid (de radi $r < R$) la força que exerceix el diferencial de pressió $\Delta P$ s'ha de veure exactament compensada per la força viscosa $F_\eta$ que el fluid fa en sentit contrari. Si coneixem que la superfície lateral del cilindre és $2 \pi r L$ i l'àrea de les tapes és $\pi r^2$, podem escriure
%
\begin{equation}
    \pi r^2 \Delta P = - \eta  2 \pi r L  \frac{dv}{dr} \, .
    \label{eqn:FluxEnTub1}
\end{equation}
%
Això és el que es coneix en matemàtiques com una equació diferencial. En aquest cas, és molt senzill resoldre-la separant les variables $v$ i $r$ i integrant des del radi $r$, on la velocitat serà $v(r)$, fins al radi del tub $R$ on tindrem $v=0$, com
%
\begin{equation}
    \int_{v(r)}^0 dv = - \frac{\Delta P}{2 \eta L} \int_r^R r \, dr\, .
    \label{eqn:FluxEnTub2}
\end{equation}
%
Resolent la integral, obtenim 
%
\begin{equation}
    \boxed{v(r) = \frac{\Delta P}{4 \eta L} \left(R^2 - r^2\right)}
    \label{eqn:FluxEnTub2}
\end{equation}
%
Podem veure que és un perfil en forma de paràbola, que es fa zero a la superfície del tub i arriba a un màxim al llarg del seu eix ($r=0$), com esperàvem.

\subsection{Llei de Poiseuille}

%
\begin{figure}[H]
    \centering
    \includegraphics[width=0.9\linewidth]{Figures/Poiseuille.pdf}
    \caption{Per deduir la llei de Poiseuille, dividim el fluid de la canonada en una successió de closques cilíndriques d'amplada infinitesimal $dr$. Això ens permet considerar que la velocitat del fluid és constant en cada closca.}
    \label{fig:Poiseuille}
\end{figure}
%
Un cop conegut el perfil de velocitats d'un fluid viscós dins un tub, és natural preguntar-se quin és el cabal total de fluid que circularà, depenent del diferencial de pressió que hi apliquem. És a dir, quin volum de fluid passarà pel tub en un temps determinat. Per això, fem servir un esquema com el de la Figura \ref{fig:Poiseuille}, aquest cop dividint el fluid en una sèrie de closques cilíndriques molt primes. En cada closca, de gruix $dr$, podrem considerar que la velocitat del fluid és constant, si fem $dr \to 0$. En aquest límit, la closca contribuirà en un cabal $dQ$. Recordant que el cabal és el producte de la velocitat per l'àrea de la secció transversal, i que la secció transversal de la closca es pot aproximar per $2 \pi r dr$, arribem a una altra equació diferencial
%
\begin{equation}
    dQ = \frac{\Delta P}{4 \eta L} \left(R^2 - r^2\right) 2 \pi r \, dr \, .
    \label{eqn:Poiseuille1}
\end{equation}
%
Aquesta equació és encara més fàcil de resoldre que l'anterior, només hem d'integrar des del centre ($r=0$) fins a la superfície del tub ($r=R$). Ens queda
%
\begin{equation}
    Q = \frac{\pi \Delta P}{2 \eta L} \int_0^R \left(R^2 - r^2\right) r \, dr \, .
    \label{eqn:Poiseuille1}
\end{equation}
%
Resolent la integral, obtenim la Llei de Poiseuille 
%
\begin{equation}
    \boxed{Q = \frac{\pi R^4}{8 \eta L} \Delta P }
    \label{eqn:Poiseuille2}
\end{equation}
%
Aquesta llei té una gran importància, perquè ens diu que el cabal de fluid que circularà pel tub depèn de la quarta potència del radi del tub. Això és crucial, per exemple, en la circulació de la sang i les malalties cardiovasculars. Si el radi interior d'una artèria queda reduït a la meitat a causa de les plaques d'ateroma, la quantitat de sang que hi circularà no es reduirà a la meitat, sinó que es reduirà en un factor $2^4 = 16$. Per tal de compensar aquest dèficit de flux, el cor hauria d'aplicar una pressió sanguínia 16 vegades major. Per aquesta raó, l'ateroesclerosi comporta un greu perill per a la salut, fins i tot quan el gruix de la placa és relativament baix.

\subsection{Resistència hidrodinàmica}

Si ens fixem en l'equació (\ref{eqn:Poiseuille2}), i tenim familiaritat amb conceptes d'electricitat, ens podem adonar que la llei de Poiseuille té una forma molt semblant a la llei d'Ohm per a circuits elèctrics: $V = RI$. En aquest cas, la diferència de pressió $\Delta P$ juga el paper de la diferència de potencial en un circuit $V$, mentre que el cabal $Q$ equivaldria a la intensitat del corrent elèctric $I$. D'aquesta manera, podríem definr la ``resistència hidrodinàmica'' $R_H$ com 
%
\begin{equation}
    R_H = \frac{8 \eta L}{\pi R^4} \, .
    \label{eqn:ResistenciaHidrodinamica}
\end{equation}
%
Quan tenim tot un sistema de tubs amb bifurcacions, podem agrupar-los depenent de si estan connectats en sèrie o en paral·lel, tal com faríem en un circuit elèctric. Quan els tubs estan connectats en sèrie, les resistències hidrodinàmiques se sumen 
%
\begin{equation}
    R_\text{sèrie} = \sum_i r_i \, ,
    \label{eqn:EnSerie}
\end{equation}
%
mentre que si estan en paral·lel, la resistència equivalent serà
%
\begin{equation}
    \frac{1}{R_\text{paral·lel}} = \sum_i \frac{1}{r_i} \, .
    \label{eqn:EnParallel}
\end{equation}
%

\section{Flux laminar i turbulent}

Fins aquest punt en el curs, hem considerat sempre que el flux del fluid és laminar i estacionari, és a dir, on les capes del fluid avancen paral·leles sense barrejar-se entre
elles ni canviar amb el temps. Aquest tipus de flux sol aparèixer quan les velocitats són petites, i el fluid es desplaça sense obstacles ni parets anguloses. El seu aspecte extern és bastant característic, ja que no s'observa cap tipus de canvi en el temps, fins al punt que un raig de fluid pot arribar a semblar sòlid, com una vareta de vidre. 
%
\begin{figure}[H]
    \centering
    \includegraphics[width=\linewidth]{Figures/FluxLaminar.pdf}
    \includegraphics[width=\linewidth]{Figures/FluxTurbulent.pdf}
    \caption{Simulació de flux laminar (a dalt) i flux turbulent (a baix) al voltant d'un obstacle esfèric. Els remolins del cas turbulent s'anomenen vòrtexs de Von Kármán, i són molt habituals a la natura. }
    \label{fig:LaminarTurbulent}
\end{figure}
%
En moltes ocasions, però, els fluids no segueixen un flux laminar, sinó que formen remolins, vòrtexs i canvis bruscos. En aquest cas, les làmines de fluid es barregen entre elles, i el flux canvia constantment, a mesura que apareixen i desapareixen remolins. Això passa quan les velocitats són grans, i el fluid troba obstacles angulosos en el seu pas. L'aspecte extern d'aquest flux és més caòtic, com per exemple un riu d'alta muntanya. 

A la Figura \ref{fig:LaminarTurbulent} es presenten dues simulacions de flux, on s'han representat amb fletxes les línies de corrent. Podem veure que en el cas del flux laminar, les línies estan ordenades suaument al voltant de l'obstacle, mentre que en el cas turbulent es produeixen oscil·lacions caòtiques. 

\subsection{Nombre de Reynolds}

Seria interessant poder quantificar d'alguna manera el grau de ``turbulència'' d'un flux. És a dir, voldríem definir algun número que sigui molt petit quan el flux és perfectament laminar, i gran quan sigui molt turbulent. Voldríem que aquest número sigui adimensional (sense unitats), i que depengui de les magnituds característiques del nostre sistema. Aquestes són la longitud $L$ i velocitat $v$ característiques del nostre sistema, i la densitat $\rho$ i viscositat del fluid $\eta$. Una manera de construir aquest número $N_R$ és 
%
\begin{equation}
    \boxed{N_R = \frac{\rho v L}{\eta}}
    \label{eqn:NombreDeReynolds}
\end{equation}
%
El número $N_R$ s'anomena Nombre de Reynolds. Experimentalment, observem que el flux és aproximadament laminar si $N_R < 1000$, mentre que es torna turbulent quan $N_R > 1300$. Podem comprovar que per sistemes macroscòpics ($v \approx 1$ m/s, $L \approx 1$ cm), el flux dins l'aigua serà turbulent ($N_R \approx 10^4$), mentre que en l'aire serà laminar ($N_R \approx 500$). En canvi, per a sistemes microscòpics més propis de la biotecnologia ($v \approx 1 \mu m$/s, $L \approx 1 \mu m$), tant l'aire com l'aigua exhibiran un flux laminar.

\section{Moviment de sòlids en fluids}

A la natura sovint ens trobem situacions en les quals un objecte sòlid s'ha de desplaçar a través de l'interior d'un fluid. En el món macroscòpic tenim animals que volen i neden, mentre que en el món microscòpic qualsevol microorganisme o cèl·lula que visqui en un medi aquós ha de desplaçar-se dins l'aigua. Aquests objectes sòlids experimentaran una força d'arrossegament o fregament que s'oposarà al seu moviment. En el món microscòpic, en el qual les velocitats solen ser baixes, tindrem típicament $N_R < 1$ i la força d'arrossegament estarà dominada per la viscositat del fluid. En canvi, per a sistemes macroscòpics solem tenir $N_R > 1000$ i l'arrossegament prové de la pròpia energia cinètica del fluid turbulent. Analitzarem cada cas per separat.

\subsection{Baixa velocitat: Llei de Stokes}

%
\begin{figure}[H]
    \centering
    \includegraphics[width=0.6\linewidth]{Figures/LleiDeStokes.pdf}
    \caption{Representació de les línies de flux al voltant d'un objecte esfèric. La viscositat del fluid produirà una força d'arrossegament donada per la llei de Stokes.}
    \label{fig:LleiDeStokes}
\end{figure}
%
Quan la velocitat és baixa ($N_R < 1$), s'observa que la força d'arrossegament és proporcional a la velocitat $v$ del cos. Aquest seria el cas, per exemple, d'un bacteri que es desplaça a través de l'aigua. En general, podem escriure
%
\begin{equation}
    F_\text{arr} = C \eta L v
    \label{eqn:ArrossegamentBaixaV}
\end{equation}
%
on $L$ és alguna longitud característica de l'objecte, com ara la mida del bacteri, i $C$ és un coeficient que dependrà de la forma del cos. En el cas particular en què tinguem un cos esfèric com a la Figura \ref{fig:LleiDeStokes}, sabem que $C = 6\pi$, que ens dóna la llei de Stokes
%
\begin{equation}
    \boxed{F_\text{Stokes} = 6 \pi \eta R v}
    \label{eqn:Stokes}
\end{equation}
%
La potència necessària per tal de moure la partícula a velocitat constant, contrarestant la força de Stokes estarà donada per
%
\begin{equation}
    \boxed{\mathcal{P} = F \, v = 6 \pi \eta R v^2}
    \label{eqn:PotenciaLaminar}
\end{equation}
%
La llei de Stokes ens permet estudiar a quina velocitat sedimentaran les petites partícules, que suposarem esfèriques, a través d'un fluid. Inicialment, si les partícules són més denses que el fluid que les envolta, patiran una acceleració cap avall a causa de la gravetat. Tot i això, quan la seva velocitat sigui suficientment gran, la força d'arrossegament serà suficientment gran per compensar el pes aparent de la partícula, és a dir el pes real menys l'empenyiment d'Arquimedes. A partir d'aquest moment, la partícula caurà a velocitat constant: la {\bf velocitat de sedimentació} $v_s$. Per trobar aquesta velocitat hem de considerar que la suma de les tres forces que actuen sobre la partícula (pes, empenyiment i arrossegament) es cancel·len entre elles. Anomenarem $\rho_p$ a la densitat de la partícula, de radi $R$, i $\rho_f$ a la densitat del fluid que l'envolta.
%
\begin{equation}
    \rho_p V g = 6 \pi \eta R v_s + \rho_f V g\, .
\end{equation}
%
Utilitzant que el volum de l'esfera és $V = \frac43 \pi R^3$, tenim que
%
\begin{equation}
    \boxed{v_s = \frac{2gR^2}{9\eta\phi} (\rho_p - \rho_f)}
\end{equation}
%
on $\phi$ s'anomena coeficient de fregament i té un valor $\phi = 1$ en el cas de partícules esfèriques. Si el fluid que envolta la partícula té una densitat molt més baixa que la pròpia partícula, com per exemple un gra de pol·len que cau a través de l'aire, podem aproximar $\rho_f \approx 0$. En aquest cas, la velocitat de sedimentació s'anomena {\bf velocitat límit} $v_l$ o velocitat terminal, i ve donada per
%
\begin{equation}
    v_l = \frac{m g}{6 \pi \eta R}\, .
\end{equation}
%
La velocitat límit d'una persona caient a través de l'aire és d'uns 200 km/h, una velocitat excessivament gran per sobreviure a una caiguda des de gran altura. Tot i això, alguns animals més petits com els insectes tenen velocitats límits suficientment petites com per poder caure des d'altures arbitràriament grans sense fer-se mal.

\subsection{Alta velocitat: Força d’arrossegament}

Quan les velocitats són altes ($N_R > 1000$), l'arrossegament sol estar domina per la pròpia energia cinètica del fluid, i per tant és proporcional a $v^2$. És el cas, per exemple, de l'arrossegament que pateixen els ocells i els avions en volar. Aquesta és la raó per la qual els ocells necessiten aletejar i els avions necessiten mantenir potència en els motors quan volen mantenir alçada. Amb un raonament semblant al que hem fet servir en l'equació (\ref{eqn:Sustentacio3}) per a la força de sustentació, tenim
%
\begin{equation}
    \boxed{F_\text{arr} = \frac12 \rho \, C_x \,  S_f \, v^2 }
    \label{eqn:Sustentacio3}
\end{equation}
%
on $S_f$ és la superfície frontal de l'ala i $C_x$ és un coeficient aerodinàmic que dependrà de la seva forma. En aquest cas la potència que es consumeix en moure la partícula a velocitat constant serà
%
\begin{equation}
    \boxed{\mathcal{P} = F \, v = \frac12 \rho \, C_x \,  S_f \, v^3 }
    \label{eqn:PotenciaTurbulent}
\end{equation}
%
Podem veure que la potència consumida és proporcional a $v^2$ en el cas del flux laminar (\ref{eqn:PotenciaLaminar}), i proporcional a $v^3$ en el cas de flux turbulent (\ref{eqn:PotenciaTurbulent}).

\subsection{El teorema de la vieira}

Una curiositat sobre el moviment d'éssers vius dins els fluids viscosos és el que es coneix com a teorema de la vieira, a causa de la peculiar forma de desplaçament d'aquests mol·luscs (pectínids). Es pot demostrar que el moviment dels fluids amb nombre de Reynolds molt baix, on la viscositat és molt important, és totalment reversible. És a dir, si movem un objecte dins d'un fluid amb $N_R \ll 1$, i després fem el moviment invers, el fluid quedarà aproximadament amb la mateixa configuració que abans. Això no passa quan el nombre de Reynolds és més alt, i la turbulència comença a jugar un paper important.

El nom del teorema ve del fet que els bivalves pectínids, coneguts com a vieires, es desplacen obrint i tancant les conquilles repetidament. Això produeix un moviment net perquè es mouen en un ambient a Reynolds prou alt, i per tant els dos moviments (obrir i tancar la conquilla) no produeixen efectes contraris tot i ser inversos l'un de l'altre. En un ambient a Reynolds baix, aquest mètode no funcionaria perquè, en obrir la conquilla, la vieira retrocediria exactament la mateixa distància que havia avançat en tancar-se. Per això en el món microscòpic els organismes han de fer servir mecanismes de propulsió molt més sofisticats, com ara flagels o cilis.



\part{Calor i Termodinàmica}

\chapter{Termodinàmica}
\label{chap:Termodinamica}

\section{Principis de la termodinàmica}

La termodinàmica és la branca de la física que tracta sobre l'energia, entesa com la magnitud que ens pot permetre generar treball. També de les maneres en les quals aquesta energia es pot transmetre entre diferents sistemes, en forma de calor i treball, i de la relació que aquesta té amb la temperatura i amb magnituds encara més abstractes com ara l'entropia. Tot i ser un camp especialment matematitzat i abstracte de la física, la termodinàmica es va crear durant el segle XIX amb una finalitat molt pràctica: millorar les màquines tèrmiques utilitzades per la Revolució Industrial. 

En un moment on la invenció de la màquina de vapor va transformar enormement la vida de les persones, personatges com Sadi Carnot, William Thomson (conegut com a Lord Kelvin) i Rudolf Clausius van invertir grans esforços en entendre les lleis que podien permetre obtenir màquines més eficients per a la indústria, el transport, i la guerra. Això va permetre enunciar una sèrie de lleis empíriques, avui conegudes com a ``principis de la termodinàmica'' i una estructura matemàtica per tractar-les i fer prediccions. 

Més endavant, amb el naixement de la física estadística, les lleis de la termodinàmica es van poder relacionar amb l'estructura microscòpica de la matèria a través del treball de persones com Ludwig Boltzmann, James Clerk Maxwell, J. Willard Gibbs i Max Planck. Curiosament, una disciplina que havia començat per motius totalment pràctics ha acabat convertint-se en una de les àrees més esteses de la ciència, amb gran importància també en la química i la biologia. Actualment els principis de la termodinàmica es consideren entre les lleis més fonamentals de la natura, i apareixen en camps tant diversos com el funcionament de la cèl·lula, els dispositius electrònics o fins i tot la física dels forats negres.

En general, en termodinàmica dividirem l'Univers en dues parts fonamentals, el sistema que volem estudiar i el seu entorn:
%
\begin{itemize}
    \item Sistema termodinàmic: Es tracta de la regió de l'Univers que estem estudiant. Pot tractar-se d'un vas de precipitats on hi ha una reacció química, una màquina tèrmica o tot un ésser viu. 
    \item Entorn: És tota la resta de l'Univers, és a dir, tot el que envolta el nostre sistema termodinàmic. Depenent del tipus de sistema que considerem, podrà intercanviar energia o fins i tot matèria amb el seu entorn.
\end{itemize}

%
\begin{figure}[H]
    \centering
    \includegraphics[width=0.7\linewidth]{Figures/SistemaTermodinamic.pdf}
    \caption{Representació esquemàtica d'un sistema termodinàmic amb energia interna $U$, que pot intercanviar calor ($Q$) i treball ($W$) amb el seu entorn.}
    \label{fig:SistemaTermodinamic}
\end{figure}
%
El nostre sistema, representat esquemàticament a la Figura \ref{fig:SistemaTermodinamic}, tindrà una energia interna, que simbolitzarem amb la lletra $U$. Aquesta energia interna no inclou l'energia cinètica o potencial a causa del moviment del sistema com un tot, sinó que té en compte el moviment intern dels àtoms i molècules que el formen. Veurem que l'energia cinètica a causa del moviment intern dels àtoms es descriu macroscòpicament com una temperatura $T$. La unitat del Sistema Internacional (SI) per l'energia és el Joule (J), tot i que també se sol utilitzar l'erg (1 erg = 10$^{-7}$ J) del sistema cgs. Finalment, per a energies molt petites, de l'escala atòmica, es fa servir l'electró-Volt (1 eV = 1,602 $\cdot$ 10$^{-19}$ J).

El sistema pot intercanviar energia amb el seu entorn en dues formes diferents:
%
\begin{itemize}
    \item Treball ($W$): El treball és l'energia transferida al sistema de forma mecànica. És a dir, a causa de forces externes exercides sobre el sistema, com per exemple el treball de compressió. Si el treball contribueix a augmentar l'energia interna del sistema el considerarem positiu ($W>0$), mentre que si la fa disminuir el considerarem negatiu ($W<0$).
    \item Calor ($Q$): La calor és l'energia transferida al sistema de forma no mecànica, a causa d'una diferència de temperatura. Per exemple, un got d'aigua calenta transferirà part de la seva energia interna en forma de calor cap a l'ambient més fred. Novament, considerem $Q>0$ quan augmenta l'energia interna, i $Q<0$ quan la fa disminuir.
\end{itemize}
%
Tant el treball com la calor es mesuren en les mateixes unitats que l'energia interna.

\subsection{Principi Zero}

El principi zero de la termodinàmica tracta sobre l'equilibri termodinàmic entre dos sistemes, i ens permetrà definir una escala de temperatura. Però què és l'equilibri termodinàmic? 

Imaginem que tenim dos sistemes separats, per exemple un got amb aigua calenta i un altre got amb aigua freda. Si ara posem aquests dos sistemes en contacte, es produirà un intercanvi de calor entre ells. En el nostre cas, intuïtivament podem imaginar que hi haurà un flux de calor des de l'aigua calenta cap a la freda. Durant aquest procés, l'aigua calenta s'anirà refredant i l'aigua freda s'anirà escalfant. Això continuarà durant un temps, fins que les propietats dels dos sistemes deixaran de canviar i es mantindran constants en el temps. En aquest moment diem que els dos sistemes estan en equilibri termodinàmic. Podeu veure una representació esquemàtica a la Figura \ref{fig:EquilibriTermic}.
%
\begin{figure}[H]
    \centering
    \includegraphics[width=0.7\linewidth]{Figures/EquilibriTermic.pdf}
    \caption{La calor es transferirà del sistema més calent al més fred, fins que els dos tinguin la mateixa temperatura. En aquest moment, diem que els dos sistemes estan en equilibri termodinàmic l'un amb l'altre.}
    \label{fig:EquilibriTermic}
\end{figure}
%
Ens agradaria poder ``etiquetar'' d'alguna manera tots els sistemes que estan en equilibri entre ells. Per això, imaginem que tenim tres sistemes termodinàmics, que anomenarem A, B i C. El principi zero de la termodinàmica ens diu: 

\begin{mdframed}
\begin{it}
    Si dos sistemes estan en equilibri termodinàmic amb un tercer, llavors estan també en equilibri entre ells
\end{it}
\end{mdframed}

És a dir, si A està en equilibri amb B, i B està en equilibri amb C, llavors A està en equilibri amb C. Això a primera vista pot semblar evident, però no totes les interaccions entre diferents elements compleixen això. Per posar un exemple senzill, imaginem tres persones: l'Anna, la Berta i en Carles. Si l'Anna és amiga de la Berta, i la Berta és amiga d'en Carles, això no significa necessàriament que l'Anna sigui amiga d'en Carles.

El principi zero ens diu que l'equilibri termodinàmic entre sistemes defineix una {\it relació d'equivalència} entre tots els sistemes que estan en equilibri entre ells. Això ens permet assignar-los una ``etiqueta'', un número que sigui igual entre tots els sistemes en equilibri, i diferent si no ho estan. Aquest número es coneix com a temperatura, i se simbolitza amb la lletra $T$.

\subsection{Primer Principi}

El primer principi de la termodinàmica, també conegut com a principi de conservació de l'energia, estableix que
%
\begin{equation}
    \boxed{\Delta U = W + Q}
    \label{eqn:Stokes}
\end{equation}
%
És a dir, el canvi en l'energia interna d'un sistema correspon a la suma del treball i la calor que s'han transferit cap a ell. L'energia no es pot crear ni destruir, sinó que només es pot transformar i transferir d'un sistema a un altre en forma de treball i/o calor. En el cas particular d'un sistema aïllat, en el que no es permet l'intercanvi de treball ni calor, $U$ es mantindrà necessàriament constant. Per això, és comú enunciar el primer principi com

\begin{mdframed}
\begin{it}
    L’energia interna d’un sistema aïllat és constant.
\end{it}
\end{mdframed}

Aquest primer principi és extremadament important, ja que ens diu que és impossible construir una màquina que ens subministri energia infinita sense consumir res a canvi. El cas més típic és la generació d'energia elèctrica, una de les principals causes de la crisi energètica actual. Per tal de produir electricitat necessitem obtenir l'energia cremant combustibles fòssils, a partir de la fissió d'elements radioactius, o a partir del treball generat pel vent o un salt d'aigua. És a dir, en una central elèctrica només podrem extraure l'energia que hi haguem introduït prèviament en una altra forma, no és possible construir una màquina que produeixi energia ``gratuïtament''.

\subsection{Segon Principi}

El segon principi restringeix la direcció en la qual es poden produir els processos termodinàmics de forma espontània. En altres paraules, mostra que hi ha processos irreversibles, que poden ocórrer espontàniament en una direcció, però no en la contrària. Tornant a l'exemple del got amb aigua calenta i el got amb aigua freda, l'energia es transferirà espontàniament de l'aigua calenta a l'aigua freda, i acabarem tenint dos gots amb aigua tèbia. En canvi si comencem amb dos gots amb aigua tèbia, és impossible que l'energia flueixi en sentit contrari espontàniament fins a tornar a tenir un got amb aigua calenta i un altre amb aigua freda. Per això, una de les maneres d'enunciar el segon principi és

\begin{mdframed}
\begin{it}
    La calor no pot fluir de manera espontània d’un sistema més fred a un sistema més calent.
\end{it}
\end{mdframed}

És important dir que el segon principi no ens prohibeix construir una nevera, o refrigerador, que trasllada calor des d'una regió freda (l'interior de la nevera) cap a una regió més calenta (l'exterior de la nevera). La clau està en la paraula ``espontània''. En un refrigerador, la calor no flueix de manera espontània, sinó que es necessita consumir energia en forma de treball per tal que ho faci. En el cas d'una nevera, aquest treball el fa el motor d'un compressor que està endollat a la xarxa de subministrament elèctric.
%
\begin{figure}[H]
    \centering
    \includegraphics[width=0.49\linewidth]{Figures/Entropia1.pdf}
    \includegraphics[width=0.49\linewidth]{Figures/Entropia2.pdf}
    \caption{El sistema de l'esquerra està més ordenat, amb les partícules blaves i les vermelles perfectament separades, i per tant té una entropia baixa. En canvi, en el sistema de la dreta les boles estan mesclades de manera aleatòria, i la seva entropia és alta. Si deixem moure les boles a l'atzar, el sistema passarà de manera espontània d'una entropia baixa a una d'alta, però mai a l'inrevés.}
    \label{fig:Entropia}
\end{figure}
%
Per poder quantificar aquesta noció d'irreversibilitat dels processos, de vegades també coneguda com a ``fletxa del temps'', es va introduir una variable termodinàmica extraordinàriament important: l'{\bf entropia} ($S$). Es pot demostrar que el segon principi implica necessàriament que l'entropia total de l'Univers es pot mantenir constant o augmentar, però mai disminuir. Els processos reversibles no modifiquen l'entropia de l'Univers, mentre que els irreversibles la fan augmentar. D'aquesta manera, és totalment impossible tornar al lloc de partida, ja que això implicaria necessàriament disminuir l'entropia de l'Univers. Si volem forçar un procés no espontani, com hem descrit en el cas de la nevera, cal que generem una gran quantitat d'entropia addicional (en aquest cas per fer funcionar el compressor), per tal de compensar amb escreix la que la nevera ha perdut.

Tot i que l'entropia va ser introduïda per Clausius, pocs anys després Ludwig Boltzmann la va interpretar com el grau de ``desordre'' en un sistema a nivell microscòpic. És a dir, un sistema que estigui microscòpicament molt ordenat, com un cristall, tindrà menor entropia que un sistema desordenat, com un gas. D'aquesta manera, tenint en compte que és molt més probable arribar a un estat desordenat que a un estat ordenat, l'entropia sempre tendirà a augmentar. Això està representat a la Figura \ref{fig:Entropia}.

El concepte d'entropia com a mesura del desordre es va anar estenent al llarg de la història. Per exemple, Claude Shannon va definir l'entropia d'informació com una mesura de la quantitat d'informació continguda en un sistema. L'ecòleg català Ramon Margalef va introduir aquest concepte d'entropia en els camps de l'ecologia i les ciències del medi ambient com una mesura matemàtica del grau de biodiversitat en un ecosistema (com per exemple un bosc). L'edifici principal de la Facultat de Biologia de la Universitat de Barcelona, on s'imparteixen aquestes classes, porta el seu nom en honor a les seves contribucions.

\subsection{Tercer Principi}

Molts experiments de la física ens han fet arribar a la conclusió de que existeix un mínim absolut per la temperatura que un cos pot assolir. Per exemple, la llei dels gasos ideals $PV = nRT$ defineix una temperatura per la qual el volum d'un gas es reduiria a zero. A aquesta temperatura, -273,15 $^\circ$C se la coneix com a zero absolut, i defineix el zero de l'escala Kelvin de temperatura ($T = 0K$). Si fos possible arribar a aquesta temperatura, el moviment dels àtoms i molècules s'aturaria completament. De fet, en un sòlid cristal·lí a 0K tindríem un sistema perfectament ordenat, i per tant amb entropia zero. El tercer principi ens diu que això és impossible:

\begin{mdframed}
\begin{it}
    El zero absolut de temperatura no és accessible mitjançant un número finit de processos.
\end{it}
\end{mdframed}

És a dir, ens podem acostar més i més al zero absolut, però mai hi arribarem del tot. Actualment s'ha aconseguit arribar dins els laboratoris a temperatures extraordinàriament properes a 0K. Concretament, l'any 2021 es van assolir temperatures de l'ordre dels picokelvin, o 10$^{-12}$ K.

\section{Calorimetria}

La calorimetria és la part de la física que estudia la mesura de les transferències de calor, habitualment a través dels canvis de temperatura. Té una gran importància en la majoria de camps de la ciència experimental, particularment en tot allò que involucri química o ciència de materials. En el camp de les ciències biomèdiques, un exemple molt directe és el càlcul de l'energia continguda en un aliment. Així doncs, quan llegim que una barreta energètica conté 150 kcal,  això significa que aquesta seria la quantitat d'energia que s'alliberaria en forma de calor si cremem la barreta. Un altre cas típic és el de l'enginyeria forestal, on sovint ens interessa saber quina energia podem obtenir a partir d'una certa quantitat de biomassa. 

\subsection{Escales de temperatura}

Per tal de mesurar temperatures, és important dissenyar mètodes que ens permetin ``llegir'' valors en termes d'una altra propietat més senzilla d'observar. És a dir, necessitem trobar una {\bf propietat termomètrica} d'algun material. Anomenem propietat termomètrica a qualsevol tipus de propietat macroscòpica (com ara el volum, densitat, resistència elèctrica, color...) que varia al variar la temperatura. En els termòmetres tradicionals de mercuri, aquesta propietat és el volum d'una certa quantitat d'aquest metall, que fa que pujar i baixar el nivell del menisc dins un tub de vidre quan la temperatura canvia. En termòmetres digitals més moderns se sol utilitzar algun element electrònic, com per exemple un termoparell, que genera una petita diferència de potencial en funció dels canvis de temperatura. Això permet processar la informació digitalment i presentar el resultat en una pantalla.

Per tal de representar la temperatura a través d'un número, cal decidir-se per una escala de temperatura. D'aquesta manera, quan diem ``avui estem a 30 graus'', és important que tothom tingui clar quina escala estem utilitzant. A la majoria de països del món, se sobreentén que estem utilitzant l'escala centígrada o Celsius, i que aquests 30 graus signifiquen que fa força calor. Tot i això, als Estats Units s'utilitza l'escala Fahrenheit, i 30 graus significa un ambient molt fred. Les principals escales utilitzades són l'escala absoluta o Kelvin, l'escala centígrada o Celsius, i l'escala Fahrenheit. Aquestes són les seves definicions:

\begin{itemize}
    \item Escala absoluta (Kelvin): Zero absolut a 0K, i punt triple de l'aigua a 273,16 K. La separació entre graus és la mateixa que la Celsius. És la que s'utilitza més en física. 
    \item Escala centígrada (graus Celsius): Es fixa 0$^\circ$C en el punt de fusió del gel i 100$^\circ$C en el punt d'ebullició de l'aigua. 
    \item Escala Fahrenheit:  Utilitzada en els països anglosaxons, principalment els Estats Units. Es fixa 0$^\circ$F en el punt de congelació de la salmorra i 96$^\circ$F a la temperatura aproximada de la sang d'un cavall saludable. En aquesta escala, el punt de fusió del gel queda a 32$^\circ$F i el punt d'ebullició de l'aigua a 212$^\circ$F.
\end{itemize}

Per tal de traduir entre escales de temperatura, hi ha equacions de conversió com
%
\begin{equation}
\begin{split}
    [K] &= [^\circ C] + 273,15\, , \\
    [^\circ C] &= [K] - 273,15\, , \\
    [^\circ F] &= \frac{9}{5} [^\circ C] + 32\, , \\
    [^\circ C] &= \frac{5}{9} \left([^\circ F] - 32\right)\, .
\end{split}
\end{equation}

\subsection{Dilatació tèrmica}

La immensa majoria de materials, sòlids líquids i gasos, augmenta de volum quan augmenta la temperatura. Això s'utilitza, tal com hem vist en l'última secció, en els termòmetres de mercuri. Aquesta propietat sol portar problemes en construir grans estructures com carreteres, vies i ponts. Per això és important, per exemple, deixar petits espais entre els diversos trams de les vies d'un tren, perquè en cas contrari la dilatació tèrmica a l'estiu podria destruir les vies. En el món biològic, és important no confondre l'efecte purament físic de la dilatació tèrmica amb altres tipus de dilatacions. Per exemple la vasodilatació de venes i artèries, que és una transformació que el cos fa activament a través de mecanismes molt més complicats.

Experimentalment veiem que l'increment de longitud d'un objecte serà més gran com més llarg fos l'objecte inicialment. És a dir, si tenim una barra de ferro d'un metre i una altra de dos metres, la segona s'allargarà el doble que la primera en augmentar la temperatura. Per això, definim el coeficient de dilatació lineal $\alpha$ com
%
\begin{equation}
    \Delta L = \alpha L \Delta T\, ,
\end{equation}
%
on $\Delta T$ és l'increment de temperatura i $\Delta L$ és l'augment de longitud. Les unitats d'$\alpha$ són de l'invers de temperatura, és a dir K$^{-1}$. A la Taula \ref{tab:DilatacioLineal} es poden veure els coeficients de dilatació lineal de diversos materials comuns.
%
\begin{table}[H]
    \centering
    \begin{tabular}{|l|c|}
    \hline
                  & $\alpha$ \\ \hline
        Celuloide & 1,0 $\cdot$ 10$^{-6} \; K^{-1}$  \\ \hline
        Pirex     & 3,2 $\cdot$ 10$^{-6} \; K^{-1}$  \\ \hline
        Vidre     & 8,3 $\cdot$ 10$^{-6} \; K^{-1}$ \\ \hline
        Acer      & 12,7 $\cdot$ 10$^{-6} \; K^{-1}$  \\ \hline
        Gel (-5$^\circ$C) & 50,7 $\cdot$ 10$^{-6} \; K^{-1}$ \\ \hline
    \end{tabular}
    \caption{Coeficients de dilatació lineal d'algunes substàncies comunes.}
    \label{tab:DilatacioLineal}
\end{table}
%
Un cop coneixem la dilatació lineal, o com augmenten les longituds, ens podem preguntar com augmentaran les àrees i els volums dels cossos. Imaginem que tenim una superfície quadrada de dimensions $x$ i $y$, que tindrà àrea $A = xy$. Llavors, l'escalfem de manera que les dimensions esdevenen $x + \Delta x$ i $y + \Delta y$. L'increment d'àrea serà, doncs
%
\begin{equation}
    \Delta A = (x + \Delta x)(y + \Delta y) - xy = x \Delta y + y \Delta x + \Delta x \Delta y\, .
\end{equation}
%
Si suposem que els increments $\Delta x$ i $\Delta y$ són bastant petits respecte a les longituds inicials $x$ i $y$, veurem que l'últim terme és molt més petit que els anteriors. Per tant, en una bona aproximació el podem ignorar. Llavors, ens queda 
%
\begin{equation}
    \Delta A \approx x \Delta y + y \Delta x = 2 \alpha x y \Delta T = \beta A \Delta T\, ,
\end{equation}
%
on hem definit el coeficient de dilatació superficial $\beta \equiv 2 \alpha $. Per un raonament similar, els volums augmentaran segons el coeficient de dilatació cúbica (o volumètrica) $\gamma \equiv 3 \alpha$. Resumint, tenim
%
\begin{equation}
    \boxed{\Delta L = \alpha L \Delta T\, , \qquad \Delta A = \beta A \Delta T\, , \qquad  \Delta V = \gamma V \Delta T }
\end{equation}
\begin{equation}
    \boxed{\beta \equiv 2 \alpha\, , \quad  \gamma \equiv 3 \alpha }
\end{equation}

\subsection{Calor específic i calor latent}

Si subministrem calor a una substància, les coses més habituals que poden passar, exceptuant les reaccions químiques que puguin tenir lloc, són 

\begin{itemize}
    \item Pot augmentar la temperatura
    \item Pot donar-se una transició de fase
\end{itemize}

Però, quina és la calor necessària per augmentar la temperatura d'una substància? I per canviar-la de fase, per exemple per fondre el gel? Aquestes quantitats es coneixen amb el nom de capacitat calorífica, calor específic i calor latent. 

\paragraph{Capacitat calorífica:} Es defineix com el quocient entre el calor transferit i la variació de temperatura. És a dir
%
\begin{equation}
    \boxed{C = \frac{\Delta Q}{\Delta T}}
\end{equation}
%
Ho podem pensar com l'energia que hem de transferir a un objecte per tal d'augmentar la temperatura en un grau centígrad (o Kelvin, ja que és un increment). És una magnitud extensiva, que es mesura en J/K en el Sistema Internacional d'unitats. Podem definir la {\bf capacitat calorífica molar} com la capacitat calorífica per mol,
%
\begin{equation}
    \boxed{c = \frac{\Delta Q}{n \Delta T}}
\end{equation}
%
on $n$ és el nombre de mols de la substància. En aquest cas és una magnitud intensiva que es mesura en J/K·mol.

\paragraph{Calor específic:} És una magnitud intensiva que correspon a la capacitat calorífica per unitat de massa, 
%
\begin{equation}
    \boxed{c_e = \frac{\Delta Q}{m \Delta T}}
\end{equation}
%
El calor específic es mesura en J/K·kg en el SI, tot i que també es fan servir cal/K·g. Per exemple, el calor específic de l'aigua és de 4184 J/K·kg o bé 1 cal/K·g (la definició de caloria). És a dir, cal una caloria per augmentar la temperatura d'un gram d'aigua en un grau centígrad (o Kelvin). A la Taula \ref{tab:CalorsEspecifics} es mostren alguns exemples comuns de calors específics.
%
\begin{table}[H]
    \centering
    \begin{tabular}{|l|c|}
    \hline
                 & $c_e$ \\ \hline
        Alumini  & 0,898 $\cdot$ 10$^{3}$ J/K·kg \\ \hline
        Ferro    & 0,443 $\cdot$ 10$^{3}$ J/K·kg \\ \hline
        Nitrogen & 1,040 $\cdot$ 10$^{3}$ J/K·kg \\ \hline
        Aigua    & 4,184 $\cdot$ 10$^{3}$ J/K·kg \\ \hline
        Hidrogen & 14,25 $\cdot$ 10$^{3}$ J/K·kg \\ \hline
    \end{tabular}
    \caption{Calors específics d'algunes substàncies comunes.}
    \label{tab:CalorsEspecifics}
\end{table}
%
\paragraph{Calor latent:} El calor latent és la calor necessària per produir un canvi de fase, per unitat de massa, és a dir
%
\begin{equation}
    \boxed{L = \frac{\Delta Q_\text{canvi de fase}}{m}}
\end{equation}
%
És una magnitud intensiva i es mesura en J/kg en el SI d'unitats. A la Taula \ref{tab:CalorsLatents} es mostren alguns exemples comuns de calors latents.
%
\begin{table}[H]
    \centering
    \begin{tabular}{|l|c|c|}
    \hline
                 & $L_f$ (fusió) & $L_e$ (ebullició) \\ \hline 
        Mercuri  & 11,8 $\cdot$ 10$^{3}$ J/kg & 272 $\cdot$ 10$^{3}$ J/kg \\ \hline
        Nitrogen & 25,5 $\cdot$ 10$^{3}$ J/kg & 201 $\cdot$ 10$^{3}$ J/kg \\ \hline
        Or       & 64,4 $\cdot$ 10$^{3}$ J/kg & 1580 $\cdot$ 10$^{3}$ J/kg \\ \hline
        Etanol   & 104  $\cdot$ 10$^{3}$ J/kg & 854 $\cdot$ 10$^{3}$ J/kg \\ \hline
        Aigua    & 333  $\cdot$ 10$^{3}$ J/kg & 2255 $\cdot$ 10$^{3}$ J/kg \\ \hline
    \end{tabular}
    \caption{Calors latents d'algunes substàncies comunes.}
    \label{tab:CalorsLatents}
\end{table}
%
Naturalment, els dos tipus de transformacions (canvis de temperatura i de fase) poden passar l'un darrera l'altre si subministrem una quantitat de calor suficient. Durant les transicions de fase, tota la calor transferida al sistema s'invertirà en modificar l'estructura microscòpica del material, de forma que la temperatura romandrà constant. A la Figura \ref{fig:TvsQ} es mostra com seria l'evolució de la temperatura en funció de la calor transferida en escalfar un bloc de gel que inicialment està a -30$^\circ$C. 
%
\begin{figure}[H]
    \centering
    \includegraphics[width=0.9\linewidth]{Figures/TvsQ.pdf}
    \caption{Esquema dels canvis de temperatura en escalfar un bloc de gel fins a evaporar-lo. Noteu que la temperatura roman constant durant les transicions de fase, en les quals hi ha una mescla sòlid-líquid o líquid-gas.}
    \label{fig:TvsQ}
\end{figure}
%
El procés tindrà lloc en diverses etapes:

\begin{enumerate}
    \item En primer lloc, la calor s'inverteix en augmentar la temperatura del gel, fins arribar al seu punt de fusió a 0$^\circ$C.
    \item A continuació, tindrem una mescla de gel i aigua líquida que romandrà a 0$^\circ$C mentre tota l'energia es fa servir per trencar l'estructura cristal·lina del gel i transformar-lo en aigua líquida.
    \item Un cop el gel hagi desaparegut per complet la temperatura tornarà a augmentar, mentre la calor s'inverteix en augmentar la temperatura de l'aigua des dels 0$^\circ$C fins al seu punt d'ebullició a 100$^\circ$C.
    \item Arribats al punt d'ebullició, tota la calor subministrada serveix per vèncer la força d'interacció entre molècules d'aigua i convertir-la en vapor. Durant l'ebullició la temperatura es manté fixa a 100$^\circ$C fins que tot el líquid hagi desaparegut.
    \item Quan tinguem només vapor, podem tornar a augmentar la temperatura per sobre dels 100$^\circ$C. Podem escalfar el vapor fins que la temperatura sigui suficientment alta com per dissociar-lo o produir alguna reacció química amb les parets del dipòsit. Això no sol passar fins als 1200$^\circ$C aproximadament.
\end{enumerate}
%
\begin{figure}[H]
    \centering
    \includegraphics[width=0.8\linewidth]{Figures/DiagramaDeFases.pdf}
    \caption{Representació esquemàtica del diagrama de fases de l'aigua, on es poden veure les tres regions corresponents a sòlid (gel), líquid i gas (vapor).}
    \label{fig:DiagramaDeFases}
\end{figure}
%
Cal especificar que les temperatures de fusió i ebullició de 0$^\circ$C i 100$^\circ$C, respectivament, corresponen a una pressió d'1 atm. Si la pressió atmosfèrica és diferent, aquestes temperatures seran unes altres. Per exemple, al punt més alt de la superfície del nostre planeta, el cim del Mont Everest, la temperatura d'ebullició és només d'uns 68$^\circ$C. Això significa que és gairebé possible beure te bullent sense fer-se massa mal.

Podem representar gràficament, en funció de $P$ i $T$, en quin estat es trobarà una substància. Aquests gràfics es coneixen com a diagrames de fase. A la Figura \ref{fig:DiagramaDeFases} es mostra el diagrama de fases de l'aigua, que inclou les seves fases principals: gel, líquid i vapor. En realitat, existeixen moltes fases diferents de gel, però per simplicitat en aquest curs les ignorarem. Les línies que separen les diverses fases són les zones de coexistència, i corresponen al punt de fusió del gel i al punt d'ebullició de l'aigua. 

Hi ha també dos punts especialment importants en el diagrama de fases. En primer lloc, el punt triple, on es creuen la línia sòlid-líquid i la líquid-gas. En el cas de l'aigua, aquest punt es troba a 0.01$^\circ$C i 0,006033 atm, i és l'unic punt on poden coexistir en equilibri les tres fases de l'aigua simultàniament: sòlid, líquid i gas. Com hem vist, aquest punt triple es va utilitzar per definir l'escala absoluta de temperatura. L'altre punt s'anomena punt crític, i es troba al final de la línia de coexistència líquid-vapor, concretament a $T_c = 374 ^\circ $C i $P_c = $218 atm. Més enllà d'aquest punt, desapareix la distinció entre líquid i vapor, i el pas de l'un a l'altre es fa de forma contínua. És a dir, si tenim un tub tancat amb aigua i vapor separats per un menisc, en superar la temperatura crítica el menisc desapareixerà i obtindrem un fluid supercrític a mig camí entre un líquid i un gas. Els fluids supercrítics tenen una gran varietat d'aplicacions en la biotecnologia i la indústria, per exemple com a dissolvent per extraure la cafeïna dels grans de cafè.

\section{Mecanismes de transmissió de la calor}

Fins aquest moment hem estat parlant de transmissió de la calor entre diferents cossos i substàncies, però no hem entrat en detall sobre els mecanismes que permeten aquesta transmissió. És a dir, els fenòmens que permeten moure energia d'un lloc a l'altre a causa de diferències de temperatura. Això passa principalment a través de tres mecanismes: conducció, convecció i radiació. Veurem que els tres canals tenen una gran importància per als éssers vius.

\subsection{Conducció}

La conducció és possiblement el mecanisme més intuïtiu de transmissió de calor, molt important en els sòlids. En aquest cas, els dos cossos a diferent temperatura es troben en contacte directe, i l'agitació tèrmica del més calent es difon cap al més fred, tal com es mostra a la Figura \ref{fig:Conduccio}. 
%
\begin{figure}[H]
    \centering
    \includegraphics[width=0.7\linewidth]{Figures/Conduccio.pdf}
    \caption{Conducció de calor a través d'una barra, des d'una font a temperatura $T_1$ fins a una altra a temperatura $T_2 < T_1$. El flux de calor $H$ dependrà de la conductivitat tèrmica $\kappa$ de la barra, així com de la seva longitud $L$ i secció $S$.}
    \label{fig:Conduccio}
\end{figure}
%
El ritme de transmissió de calor $H$, mesurat en Watts (W) ve donat per la llei de Fourier 
%
\begin{equation}
    \boxed{H_\text{conducció} = \frac{\Delta Q}{\Delta t} = \kappa S \frac{\Delta T}{L}}
\end{equation}
%
on $S$ és l'àrea de la secció transversal del conductor central, $L$ és  la seva longitud i $\Delta T$ és la diferència de temperatures. El coeficient de proporcionalitat $\kappa$ s'anomena {\bf conductivitat tèrmica} i es mesura en W/m·K en el Sistema internacional d'unitats. A la Taula \ref{tab:Conductivitats} es mostren alguns exemples de conductivitat tèrmica de materials comuns.
%
\begin{table}[H]
    \centering
    \begin{tabular}{|l|c|}
    \hline
                & $\kappa$ \\ \hline
        Coure   & 400 W/m·K \\ \hline
        Alumini & 240 W/m·K \\ \hline
        Aigua   & 0,59 W/m·K \\ \hline
        Fusta   & 0,08 W/m·K \\ \hline
        Aire    & 0,024 W/m·K \\ \hline
    \end{tabular}
    \caption{Conductivitats tèrmiques d'algunes substàncies comunes.}
    \label{tab:Conductivitats}
\end{table}
%
Se solen anomenar bons conductors tèrmics els materials amb una alta conductivitat tèrmica, mentre que els que tenen baixa conductivitat es coneixen com a aïllants tèrmics. De la mateixa manera que en el cas de la llei de Poiseuille (vegeu el Capítol \ref{chap:Viscositat}), podem reescriure la llei de Fourier en forma anàloga a la llei d'Ohm definint la ``resistència tèrmica'' com
%
\begin{equation}
    R_T = \frac{L}{\kappa S}
\end{equation}
%
Quan tenim tot un sistema de conductors tèrmics amb bifurcacions, podem agrupar-los depenent de si estan connectats en sèrie o en paral·lel, tal com faríem en un circuit elèctric. Quan els conductors estan connectats en sèrie, les resistències tèrmiques se sumen 
%
\begin{equation}
    R_\text{sèrie} = \sum_i r_i \, ,
\end{equation}
%
mentre que si estan en paral·lel, la resistència equivalent serà
%
\begin{equation}
    \frac{1}{R_\text{paral·lel}} = \sum_i \frac{1}{r_i} \, .
\end{equation}
%

\subsection{Convecció}

La convecció és el mecanisme dominant en la transmissió de calor a través de líquids i gasos (fluids). En aquest cas, el que es mou no és només l'agitació tèrmica de la matèria, sinó que és la pròpia matèria la que es desplaça. Aquest desplaçament està produït pel principi d'Arquimedes (vegeu el Capítol \ref{chap:Hidrostatica}) o per un impuls extern (convecció forçada). 

Per entendre-ho millor, podem pensar en un sistema de calefacció per terra radiant en una habitació, i com s'acaba escalfant tot l'aire de la sala. Inicialment, l'aire que està en contacte directe amb el terra calent augmentarà la seva temperatura per conducció. En escalfar-se, a causa de la dilatació tèrmica, disminuirà la seva densitat i l'empenyiment d'Arquimedes el farà pujar cap al sostre. En pujar, serà reemplaçat per aire fred, més dens. D'aquesta manera, es produirà una circulació contínua d'aire calent que puja i aire fred que baixa, sovint de forma turbulenta. Això contribuirà a distribuir la calor provinent del terra per tota l'habitació. La Figura \ref{fig:Conveccio} mostra de manera molt esquemàtica aquest procés.
%
\begin{figure}[H]
    \centering
    \includegraphics[width=0.7\linewidth]{Figures/Conveccio.pdf}
    \caption{Transmissió de calor en un fluid per convecció, dins un recipient escalfat per la part de sota. El fluid de la part inferior s'ha escalfat per contacte amb el terra calent, i la seva densitat és menor que la del fluid fred de la part superior. Per això, tendeix a pujar mentre que el fluid fred tendeix a baixar. D'aquesta manera s'estableix una circulació que redistribueix la calor per tot el fluid.}
    \label{fig:Conveccio}
\end{figure}
%
En el cas de la convecció, el ritme de pèrdua de calor d'un objecte submergit en un fluid $H$, mesurat en Watts (W), ve donat per la llei de Newton 
%
\begin{equation}
    \boxed{H_\text{convecció} = \frac{\Delta Q}{\Delta t} = q S \Delta T}
\end{equation}
%
on $S$ és la superfície de l'objecte que està en contacte amb el fluid i $\Delta T$ és la diferència de temperatures. El coeficient de proporcionalitat $q$ s'anomena {\bf constant de transmissió per convecció} i es mesura en W/m$^2$·K en el Sistema internacional d'unitats.

La convecció és molt important en les ciències de la Terra a causa del seu paper en el clima, a través de corrents oceàniques, marines i atmosfèriques. Per exemple, en una tempesta, es forma el que es coneix com a cèl·lula convectiva, en la qual el calor latent que s'ha desprès de la condensació de l'aigua fa que l'aire s'escalfi i ascendeixi a gran velocitat a través de núvol.

Finalment, cal mencionar el concepte de {\bf convecció forçada}. Si l'empenyiment d'Arquimedes a causa de les diferències de densitat {\bf convecció natural} no és suficient per transmetre la calor de forma prou ràpida, podem utilitzar un ventilador o bomba per incrementar la circulació d'aire. Això s'utilitza en la majoria de sistemes de calefacció, i en forns per cuinar.

\subsection{Radiació}

Fins ara, hem vist que la conducció funciona principalment en els sòlids, i la convecció en els líquids i gasos. És a dir, a través de la matèria. Ens podem preguntar, doncs: És possible transmetre calor a través del buit? Clarament sí, ja que el calor de la nostra estrella, el Sol, ens arriba a la Terra a través d'espai buit. De fet, viatja a gran velocitat en forma de radiació electromagnètica, incloent-hi la llum visible. Això ens indica el tercer mecanisme de transmissió de calor: la radiació. 

Experimentalment, veiem que qualsevol objecte emet radiació electromagnètica a causa de la seva temperatura. Això ho podem veure en la nostra vida quotidiana en objectes molt calents, com per exemple una peça de ferro que ha estat dins el foc, en que una part significativa de l'emissió és de llum visible (normalment de color vermell). Quan passa això diem que el ferro està ``roent'' o ``al roig viu''. En objectes més freds, com el nostre propi cos, la majoria de l'emissió és en l'espectre infraroig, que només es pot veure amb càmeres especials.

L'estudi de la radiació tèrmica dels objectes va tenir una immensa importància, ja que va motivar el descobriment de la {\bf física quàntica} per Max Planck l'any 1900. La física quàntica va ser necessària per explicar la distribució de les diverses longituds d'ona que emet un cos a certa temperatura, una distribució anomenada {\bf espectre de Planck} 
%
\begin{equation}
    \boxed{E(\lambda, T) = \frac{2hc^2}{\lambda^5}\frac{1}{e^{h c / \lambda k T} - 1}}
    \label{eqn:LleiDePlanck}
\end{equation}
%
En aquesta equació, $E(\lambda, T)$ és la intensitat de llum, $\lambda$ és la longitud d'ona, $c = 2,998\cdot 10^{8}$ m/s és la velocitat de la llum, $h = 6,626 \cdot 10^{-34}$ m$^2$·kg/s s'anomena la constant de Planck i $k = 1,381 \cdot 10^{-23}$ m$^2$·kg·s$^{-2}$·K$^{-1}$ és la constant de Boltzmann. A la Figura \ref{fig:EspectreDePlanck} es pot veure una representació gràfica dels espectres de Planck a diverses temperatures.
%
\begin{figure}[H]
    \centering
    \includegraphics[width=0.8\linewidth]{Figures/EspectreDePlanck.pdf}
    \caption{Corbes d'intensitat lluminosa emesa per un cos a diferents temperatures, segons la llei de Planck (\ref{eqn:LleiDePlanck}). La zona entre 400 nm i 750 nm correspon a la llum visible, passant per tots els colors de l'arc de Sant Martí. }
    \label{fig:EspectreDePlanck}
\end{figure}
%
Podeu observar a la mateixa Figura que, com més alta és la temperatura, més petita és la longitud d'ona. És a dir, si anem escalfant un objecte, primer brillarà de color vermell, després groc-blanc, i finalment blau. A major temperatura, el pic es continuarà desplaçant cap a l'ultravioleta. La llei que ens diu la posició d'aquest màxim en funció de la temperatura és la {\bf llei del desplaçament de Wien}
%
\begin{equation}
    \boxed{\lambda_\text{max} = \frac{B}{T}}
\end{equation}
%
on $B = 2898 \, \mu$m·K s'anomena constant de desplaçament de Wien. Una observació interessant, en relació a la biologia, és que els nostres ulls s'han adaptat evolutivament per tenir una màxima sensibilitat a la llum de $\lambda = 550$ nm, que segons la llei de Wien correspon a una temperatura de 5270K, la temperatura de la superfície del Sol. És a dir, estudiant el funcionament dels nostres ulls podem conèixer la temperatura del Sol!

Finalment, el ritme total d'emissió d'energia $H$ per radiació d'un cos a temperatura $T$ ve donat per la {\bf llei de Stefan-Boltzmann}
%
\begin{equation}
    \boxed{H_\text{radiació} = e \sigma S T^4}
\end{equation}
%
en aquest cas $S$ és la superfície del cos i $\sigma = 5,670 \cdot 10^{-8}$ W·m$^{-2}$·K$^{-4}$ és la constant de Stefan-Boltzmann. El coeficient $e$ és un nombre entre 0 i 1 anomenat emissivitat, i indica si el cos és bon o mal emissor. Si el cos a temperatura $T_\text{cos}$ està en un ambient a temperatura $T_\text{amb}$, estarà emetent i absorbint radiació simultàniament fins que les dues temperatures s'igualin. per tant, en aquest cas tindrem 
%
\begin{equation}
    H = e \sigma S \left(T_\text{cos}^4 - T_\text{amb}^4\right)\, .
\end{equation}
%
Noteu que l'emissivitat $e$ és també el coeficient que es fa servir en l'absorció de radiació ambiental. Això s'anomena la {\bf llei de Kirchhoff}, i és necessari per tal que la transferència neta d'energia es faci zero en arribar a l'equilibri termodinàmic $T_\text{cos} = T_\text{amb}$. És a dir, si un cos és un bon absorbent és un bon emissor ($e$ proper a 1), i si és mal absorbent és mal emissor ($e$ proper a 0). Un absorbent/emissor perfecte, amb $e = 1$ és un cas idealitzat que es coneix com a cos negre. És per això que la radiació tèrmica es coneix també com a {\bf radiació de cos negre}.

\subsection{Lleis de Bergmann i Allen}

En el món dels éssers vius, les lleis de transmissió de la calor produeixen alguns efectes curiosos, que s'originen en la necessitat dels animals de mantenir una temperatura corporal en diversos climes. 

Això és degut a que la pèrdua de calor dels animals, principalment per convecció i radiació, és proporcional a la superfície exterior de l'animal $S$. Per altra banda, la producció de calor està repartida per tot el cos de l'animal, i per tant és proporcional al volum $V$. D'aquesta manera, si la longitud característica de l'animal és $L$, el quocient entre pèrdua i generació és aproximadament $S/V = L^2/L^3 = 1/L$. És a dir, com més gran sigui l'animal, més bé conservarà la temperatura.

És per això que, com a norma general, els animals solen ser més grossos en climes més freds i més petits en climes càlids. Aquesta norma, amb nombroses excepcions, es coneix com a {\bf llei de Bergmann}. També, per tal de minimitzar la superfície exterior, els animals de climes freds solen ser més rabassuts i arrodonits, amb potes gruixudes i curtes. En canvi, els animals de climes més càlids tendeixen a ser més estilitzats, amb potes llargues i primes. Aquest fenomen es coneix com a {\bf llei d'Allen}. 

També, per tal de protegir-se del fred, els animals han desenvolupat tota una sèrie de mecanismes d'aïllament tèrmic, fonamentalment la capa subcutània de greix, el pèl i les plomes. El greix actua simplement a causa de la seva baixa conductivitat tèrmica, que dificulta el flux de calor entre els òrgans interns i l'exterior. En el cas del pèl i les plomes, el seu funcionament es basa en crear una gran quantitat de cambres d'aire entre filaments, que dificulta enormement la circulació de l'aire i per tant redueix les pèrdues per convecció.

\subsection{Efecte hivernacle i canvi climàtic}

Finalment, la transmissió de calor provinent del Sol té una gran importància en el clima i els seus canvis. La Terra captura una bona part de l'energia lluminosa que ens arriba del Sol a través d'un mecanisme conegut com a {\bf efecte hivernacle}. Aquest efecte es produeix perquè la nostra atmosfera, a causa de diversos gasos que conté, com ara el diòxid de carboni i el vapor d'aigua, és transparent per algunes longituds d'ones electromagnètiques i opaca per algunes altres. 

Així, la radiació que ens arriba del Sol és majoritàriament en el rang visible, on l'atmosfera és transparent, i per tant no té cap problema per arribar a la superfície terrestre. En canvi, la superfície està molt més freda que el Sol, i emet en el rang infraroig, on l'atmosfera és opaca. Per tant, una part important de l'energia solar queda retinguda a la Terra, cosa que permet mantenir una temperatura ambiental compatible amb la vida. Desgraciadament l'excés d'emissions de diòxid de carboni provinents de la crema de combustibles fòssils des de la Revolució Industrial ha intensificat artificialment aquest efecte hivernacle. Això ens ha conduït a una crisi climàtica coneguda com a escalfament global, o senzillament canvi climàtic.

Un altre efecte que pot afectar l'absorció d'energia de la Terra és la capacitat de reflectir la llum solar, és a dir el color de la seva superfície. Aquesta reflectivitat es coneix com a {\bf albedo}. Per exemple, les acumulacions de gel als casquets polars són de color blanc i reflecteixen una part significativa de la llum solar cap a l'espai. Si la temperatura de la Terra augmenta i aquest gel es fon, la proporció de llum solar que quedarà absorbida serà major, cosa que contribuirà a empitjorar el problema climàtic encara més. 

%%%%%%%%%%%%%%%%%%%%%%%%%%%%%%%%%%%%%%%%%%%%%%%%%%%%%%%%%%%%%%%%%%%%%%%%%%
%%%%%%%%%%%%%%%%%%%%%%%%%%%%%%%%%%%%%%%%%%%%%%%%%%%%%%%%%%%%%%%%%%%%%%%%%%


\part{Ones i fenòmens ondulatoris}

\chapter{Ones mecàniques i sonores}
\label{chap:OnesMecaniques}

\section{Propietats i classificació}

Els fenòmens periòdics i ondulatoris estan presents en totes les àrees de la ciència, tot i que naturalment s'estudien sobretot a les assignatures de física. Podem definir una ona com una pertorbació en un conjunt d'objectes, partícules o camps, que es propaga en l'espai al llarg del temps sense propagar el propi medi. Definida així, és bastant complicat entendre a què ens referim exactament, de manera que ho exemplificarem amb una situació més o menys quotidiana: una ona en les grades d'un estadi de futbol, coneguda popularment com ``la ola''.

En aquesta ona humana, la pertorbació que propaguem és la posició del cos de cada persona. És a dir, una aficionada que es trobi en un seient concret de la graderia estarà inicialment asseguda. Quan observi que la persona del seu costat s'aixeca, s'aixecarà també, i es mantindrà dreta fins que la persona del seu costat es torni a asseure. En aquest moment, ella es tornarà a asseure novament. Si tota la resta d'aficionats i aficionades fan exactament el mateix, què és el que veurà un dron que estigui volant sobre l'estadi?

A la filmació d'aquest dron, que probablement després es retransmeti per televisió o per Internet, costarà molt distingir el moviment individual de cada persona. De fet, és possible que ni tan sols es puguin veure les persones amb claredat. El que sí que veurem és una gran onada, de molts metres d'amplada, que viatja a gran velocitat per les graderies, i pot arribar a donar diverses voltes a l'estadi. Això és una ona! Fixem-nos que, mentre l'ona ha viatjat lateralment al llarg de les grades, la nostra aficionada no s'ha mogut del seu seient. La pertorbació (la posició corporal de les persones) s'ha propagat, però el medi (les pròpies persones) no s'ha mogut del lloc.

Ara bé, el tipus de pertorbació que es propaga pot ser qualsevol cosa, no necessàriament la postura espacial d'una persona. Alguns exemples molt comuns són:
%
\begin{itemize}
    \item Ones mecàniques: El cas més senzill, quan el que pertorbem és el moviment físic d'un medi elàstic. Per exemple, el moviment d'una corda, les ones sonores, les ones sísmiques, o les deformacions en un objecte sòlid.
    %
    \item Ones sonores: Un cas particular d'ones mecàniques. Ones en la pressió de l'aire, que entre altres coses formen el so que sentim amb les nostres orelles. Aquestes ones de pressió formen la música, la parla, els ultrasons que s'utilitzen en una ecografia, i el sistema d'ecolocalització nocturna dels rat-penats. 
    %
    \item Ones electromagnètiques: En aquest cas pertorbem els camps electromagnètics, que poden existir fins i tot en el buit. L'exemple més natural és la llum visible, amb tots els colors de l'Arc de Sant Martí. Però aquestes ones inclouen totes les regions de l'espectre electromagnètic: ones de ràdio, microones, infraroig, llum visible, llum ultravioleta, raigs X i raigs gamma ($\gamma$).
    %
    \item Funció d'ona quàntica: En aquest cas parlem d'ones de probabilitat. En aquest cas la funció d'ona $\psi(x_i)$ ens descriu l'amplitud complexa de la probabilitat de trobar el nostre sistema físic en una configuració determinada. És la base del funcionament de tota la química, en particular en el comportament de biomolècules. En el cas concret del Grau de Biotecnologia, s'estudia en l'assignatura optativa de Física Quàntica.
    %
    \item Ones gravitacionals: Un exemple molt més exòtic, i sense aplicació en ciències biològiques, però que ha guanyat molta fama en l'última dècada. Es tracta d'ones en la deformació de l'espaitemps. Es prediuen a partir de les equacions d'Einstein, publicades l'any 1915, i no es van detectar fins un segle més tard, el 14 de setembre del 2015, als detectors LIGO (Laser Interferometer Gravitational-wave Observatory).
\end{itemize}
%
Es distingeixen en general dos tipus d'ones, transversals i longitudinals, tal com es mostra a la Figura \ref{fig:LongitudinalTransversal}. Les seves diferències són:
%
\begin{itemize}
    \item Ones transversals: El cas més comú. En aquest cas la pertorbació del medi és perpendicular a la direcció de propagació de l'ona. Les ones en una corda, les electromagnètiques o les ones gravitacionals són exemples d'ones transversals.
    %
    \item Ones longitudinals: La pertorbació es produeix en la mateixa direcció que la propagació de l'ona. Un exemple típic són les ones sonores en l'aire.
\end{itemize}
%
\begin{figure}[H]
    \centering
    \includegraphics[width=\linewidth]{Figures/LongitudinalTransversal.pdf}
    \caption{Representació d'una ona longitudinal (a dalt), on el desplaçament de les partícules es produeix en la direcció de propagació, i una ona transversal (a baix), on les partícules es mouen perpendicularment a la direcció de l'ona.}
    \label{fig:LongitudinalTransversal}
\end{figure}
%

\section{Propagació de les ones}
%
Per poder descriure una ona, en general resumim la propietat que pertorbem a través d'una funció d'ona. Aquesta funció dependrà de l'espai i del temps, i pot representar el desplaçament lateral d'una corda, la pressió de l'aire, o el camp electromagnètic. Per agafar un exemple el més visual possible, per a la majoria d'aquest tema prendrem un pols que viatja al llarg d'una corda. La funció d'ona en aquest cas serà $y(x, t)$, és a dir l'altura de la corda en una posició $x$, a l'instant de temps $t$, com es veu a la Figura \ref{fig:FuncioOna}. 
%
\begin{figure}[H]
    \centering
    \includegraphics[width=0.8\linewidth]{Figures/FuncioOna.pdf}
    \caption{La funció d'ona pot representar qualsevol pertorbació que es propagui, en aquest cas l'altura de cada un dels punts d'una corda. S'hi poden veure marcades l'amplitud $A$ i la longitud d'ona $\lambda$. La longitud d'ona es pot definir com la distància entre dues crestes, o la distància entre dos nodes ascendents (o descendents).}
    \label{fig:FuncioOna}
\end{figure}
%
Assumirem també que totes les nostres ones es mouen en un medi no dispersiu, és a dir, on totes les ones es mouen a la mateixa velocitat $v$. En una corda, la velocitat de propagació ve donada per
%
\begin{equation}
    \boxed{v = \sqrt{\frac{\tau}{\mu}}}
    \label{eqn:VelocitatCorda}
\end{equation}
%
on $\tau$ és la tensió de la corda (mesurada en Newtons), i $\mu$ la seva densitat lineal (en quilograms per metre). La interpretació d'aquesta equació és força intuïtiva, ja que variar la tensió de les cordes és el mètode habitual per afinar qualsevol instrument de corda. Més endavant veurem la relació exacta entre la velocitat $v$ i l'afinació d'una corda. Però, què passa si un pols passa des d'un tros de corda amb velocitat $v_1$ a un altre tros amb velocitat $v_2$? En aquest cas, entenem que les dues cordes tindran la mateixa tensió (estan lligades una darrere l'altra), i que la diferència de velocitat està causada per diferents densitats lineals (una corda és més gruixuda que l'altra).

El que veurem que passa és que una part de l'ona continua el seu camí a través de la segona corda (l'ona transmesa) i una altra part tornarà enrere (l'ona reflectida). Això és exactament el que passa amb la llum a través d'un vidre. És a dir, part de la llum que conté la nostra imatge es transmetrà a través del vidre (i per tant la gent que hi hagi a l'altre costat ens podrà veure), i una altra part es reflectirà (de manera que nosaltres mateixos ens veurem reflectits al vidre). Si el pols incident arriba des de l'esquerra amb una alçada (o amplitud) $A_I$, les amplituds dels polsos transmès $A_T$ i reflectit $A_R$ venen donades per les següents expressions, que no demostrarem en aquest curs, 
%
\begin{equation}
    \boxed{A_R = \frac{v_2 - v_1}{v_1 + v_2} A_I\, , \quad A_T = \frac{2v_2}{v_1 + v_2} A_I}
    \label{eqn:ReflexioTransmissio}
\end{equation}
%
Podem analitzar per separat dos casos. En el primer, la segona corda és més pesada ($\mu_2 > \mu_1$) i per tant la seva velocitat de propagació serà menor ($v_2 < v_1$), tal com es veu a la Figura \ref{fig:CordesPrimaGruixuda}. En aquest cas, veiem que el pols reflectit s’inverteix (amb menor amplitud), i el pols transmès també té menor amplitud i és més estret que l'incident.
%
\begin{figure}[H]
    \centering
    \includegraphics[width=0.8\linewidth]{Figures/CordesPrimaGruixuda.pdf}
    \caption{Transmissió d'un pols en la interfase entre dues cordes. La segona és més massiva que la primera ($\mu_2 > \mu_1$) i per tant la seva velocitat de propagació és més baixa. El pols reflectit s’inverteix amb menor amplitud, i el pols transmès també té menor amplitud i és més estret que l'incident. }
    \label{fig:CordesPrimaGruixuda}
\end{figure}
%
En el segon cas, la segona corda és més lleugera ($\mu_2 < \mu_1$) i per tant la seva velocitat de propagació serà major ($v_2 > v_1$), tal com es veu a la Figura. Ara, el pols reflectit ja no s’inverteix (però segueix tenint menor amplitud), però el pols transmès té major amplitud i es fa més llarg.
%
\begin{figure}[H]
    \centering
    \includegraphics[width=0.8\linewidth]{Figures/CordesGruixudaPrima.pdf}
    \caption{Transmissió d'un pols en la interfase entre dues cordes. La segona és més lleugera que la primera ($\mu_2 < \mu_1$) i per tant la seva velocitat de propagació és més alta. El pols reflectit es manté en la mateixa direcció (no s'inverteix), i el pols transmès té major amplitud i es fa més llarg.}
    \label{fig:CordesGruixudaPrima}
\end{figure}
%
\begin{figure}[H]
    \centering
    \includegraphics[width=0.7\linewidth]{Figures/ReflexioExtremFix.pdf}
    \caption{Si l'extrem de la corda està fix, el pols reflectit és de la mateixa amplitud que l'incident, però invertit. Podem imaginar aquest cas com el límit $v_2 \to 0$.}
    \label{fig:ReflexioExtremFix}
\end{figure}
%
Podem preguntar-nos què passa en els casos extrems, és a dir, quan el segon medi és impossible de moure (extrem fix) o no costa gens de moure (extrem lliure). Com veurem més endavant a la secció d'ones sonores, aquests casos seran anàlegs a un tub tancat o obert per un extrem, en un instrument de vent. A la Figura \ref{fig:ReflexioExtremFix} es pot veure el cas amb un extrem fix, i a la Figura \ref{fig:ReflexioExtremLliure} el cas amb un extrem lliure.
%
\begin{figure}[H]
    \centering
    \includegraphics[width=0.7\linewidth]{Figures/ReflexioExtremLliure.pdf}
    \caption{Si l'extrem de la corda està lliure, el pols reflectit és de la mateixa amplitud que l'incident, i no s'inverteix. Podem imaginar aquest cas com el límit $v_2 \to \infty$.}
    \label{fig:ReflexioExtremLliure}
\end{figure}
%

\subsection{Ones viatgeres}

En general, en un sistema unidimensional, la funció d'ona sempre es podrà escriure de forma general com
%
\begin{equation}
    y(x, t) = f(x - vt) + g(x + vt)\, ,
    \label{eqn:FuncioOna}
\end{equation}
%
on $f(x)$ és la forma (a $t=0$) d'una ona que viatja cap a la dreta, i $g(x)$ és la forma (també a $t=0$) de l'ona que viatja cap a l'esquerra. Anem a veure per què això és així. A la Figura \ref{fig:OnaViatgera} podem veure un pols, amb una forma complicada, que es propaga cap a la dreta, és a dir, tenim $y(x,t) = f(x - v t)$. 
%
\begin{figure}[H]
    \centering
    \includegraphics[width=0.7\linewidth]{Figures/OnaViatgera.pdf}
    \caption{Representació d'un pols que es desplaça cap a la dreta amb velocitat $v$, representada com $y(x, t) = f(x - vt)$. Per a $t=0$, el pic es troba a la posició $x=x_0$, i en qualsevol temps posterior a $x=x_0+vt$.}
    \label{fig:OnaViatgera}
\end{figure}
%
A l'instant inicial, quan $t = 0$, tindrem  $y(x,0) = f(x)$, com esperàvem. En aquest moment, ens podem fixar en un punt concret del perfil, per exemple el pic que es troba a $x = x_0$. És a dir $y(x_0, 0) = f(x_0)$ és màxima. En un temps posterior $t$, si l'ona es propaga cap a la dreta a una velocitat $v$, esperem que aquest màxim estigui a una posició $x = x_0 + vt$. Podem comprovar que efectivament és així: $y(x_0 + vt, t) = f(x_0 + vt - vt) = f(x_0)$, altre cop el valor màxim. Amb un raonament similar, podem veure que $g(x)$ és la forma de l'ona que es propaga cap a l'esquerra.
%
\begin{figure}[H]
    \centering
    \includegraphics[width=0.8\linewidth]{Figures/FuncioOna.pdf}
    \caption{Forma característica d'una ona sinusoidal, del tipus $f(x) = A \sin (k x - \phi)$. S'hi poden veure marcades l'amplitud $A$ i la longitud d'ona $\lambda = 2\pi/k$.}
    \label{fig:OnaSinusoidal}
\end{figure}
%
La forma més senzilla que pot tenir una ona és la d'una ona sinusoidal, que és la que obtindrem si la font de la pertorbació segueix un moviment harmònic simple. En aquest cas, la forma de l'ona serà $f(x) = A \sin (k x - \phi)$, on $k = 2\pi/\lambda$ s'anomena nombre d'ona, i es mesura en radians per metre (rad/m). La quantitat $\lambda$ és la longitud d'ona, la longitud entre dos pics o dues valls de l'ona, tal i com es mostra a la Figura \ref{fig:OnaSinusoidal}. Si suposem que l'ona es propaga cap a la dreta, tindrem que
%
\begin{equation}
    y_{\rightarrow}(x, t) = A \sin \left[ k (x - vt) + \phi \right] = A \sin ( k x - \omega t + \phi )\, ,
\end{equation}
%
mentre que si l'ona es propaga cap a l'esquerra, tindrem que $g(x) = A \sin (k x - \phi)$, i per tant 
%
\begin{equation}
    y_{\leftarrow}(x, t) = A \sin \left[ k (x + vt) + \phi \right] = A \sin ( k x + \omega t + \phi )\, .
\end{equation}
%
En aquestes relacions, $\omega = k v = 2\pi / T$ és la freqüència angular, mesurada en radians per segon (rad/s). $T$ és el període, és a dir el temps que un punt de la corda triga a fer un cicle complet. Cada cicle correspon a $2\pi$ radians, i per tant podem definir la freqüència com $f = 1 / T = \omega / 2 \pi$. La freqüència $f$ (de vegades denotada també per la lletra grega $\nu$) es mesura en Hertz (Hz) o cicles per segon. A continuació es mostra una petita taula amb algunes relacions entre aquestes quantitats, juntament amb les seves unitats en el Sistema Internacional d'unitats.
%
\begin{equation}
    \boxed{
    \begin{split}
        \omega &= \frac{2\pi}{T} = 2\pi f\, , \quad  \qquad \text{(Mesurada en rad/s)}\\
        k &= \frac{2\pi}{\lambda} = \frac{\omega}{v}\, , \qquad  \qquad \text{(Mesurada en rad/m)}\\
        f &= \frac{\omega}{2\pi} = \frac{v}{\lambda} = \frac{1}{T}\, , \qquad \text{(Mesurada en Hz, o s$^{-1}$)} \\
        v &= \frac{\lambda}{T} = \frac{\omega}{k} = \lambda f \, .  \qquad \text{(Mesurada en m/s)}
    \end{split}
    }
\end{equation}

\section{Ones estacionàries}

Pensem ara en la corda d'un violí, que està subjecta pels seus dos extrems. En aquest cas, les ones viatgeres que s'hi propaguin es reflectiran contínuament en arribar a cada un dels extrems. D'aquesta manera, de forma espontània tindrem dues ones iguals, amb la mateixa amplitud i longitud d'ona, propagant-se en direccions oposades. D'aquesta manera, podem escriure la funció d'ona total com a  
%
\begin{equation}
     y_T(x, t) = y_{\rightarrow}(x, t) + y_{\leftarrow}(x, t) = A \sin ( k x - \omega t + \phi ) + A \sin ( k x + \omega t + \phi )\, .
\end{equation}
%
Aquí podem utilitzar una identitat trigonomètrica bastant coneguda,
%
\begin{equation}
    \sin A + \sin B = 2 \sin \left(\frac{A + B}{2}\right)  \cos \left(\frac{A - B}{2}\right)\, ,
\end{equation}
%
que ens condueix a la forma 
%
\begin{equation}
     y_T(x, t) = 2 A \sin ( k x + \phi ) \cos ( \omega t )\, .
\end{equation}
%
El resultat és el que es coneix com a ona estacionària, en la qual la dependència espacial i temporal estan multiplicades. Anem a analitzar aquesta forma en detall, ja que és extremadament important en la majoria de sistemes físics que involucren ones i ressonància, entre ells els instruments musicals. Fixem-nos primer en la part purament espacial, $2 A \sin ( k x + \phi )$. Com veiem a la Figura, és una forma sinusoidal semblant a la que teníem en el cas de les ones viatgeres. Però aquesta ona no viatja! En lloc d'això, l'únic que fa la part temporal és modular (inflar i desinflar) aquesta forma sinusoidal periòdicament. En certa manera, podem pensar que cada punt de la corda segueix un moviment harmònic simple, amb una amplitud que depèn de la seva posició espacial. Aquest efecte s'intenta representar a la Figura \ref{fig:OnaEstacionaria} afegint línies semitransparents, corresponents a la posició de la corda en diversos instants de temps.
%
\begin{figure}[H]
    \centering
    \includegraphics[width=0.9\linewidth]{Figures/OnaEstacionaria.pdf}
    \caption{Esquema d'una ona estacionària en una corda, on les línies semitransparents intenten representar diferents posicions al llarg del temps. Els nodes (els punts que no oscil·len) estan marcats amb esferes negres, i els ventres (els punts que oscil·len amb màxima amplitud) amb esferes vermelles.}
    \label{fig:OnaEstacionaria}
\end{figure}
%
Hi ha dos tipus de punts especialment importants, al llarg de la corda:
%
\begin{itemize}
    \item Ventres: Són punts en els que la corda oscil·la amb la màxima amplitud. A la Figura estan marcats de color vermell. En aquests punts, les ones viatgeres $y_{\rightarrow}(x, t)$ i $y_{\leftarrow}(x, t)$ se sumen l'una a l'altra en tot moment, de manera que l'amplitud resultant és el doble: $2A$. Aquest efecte es coneix com a {\bf interferència constructiva}.
    \item Nodes: Són punts en els que la corda no oscil·la en absolut. A la Figura estan marcats de color negre. En aquests punts, les ones viatgeres $y_{\rightarrow}(x, t)$ i $y_{\leftarrow}(x, t)$ es resten l'una a l'altra en tot moment. Aquest efecte es coneix com a {\bf interferència destructiva}.
\end{itemize}

\subsection{Interferència}
-
\begin{figure}[H]
    \centering
    \includegraphics[width=0.9\linewidth]{Figures/Interferencia.pdf}
    \caption{Interferència entre les ones provinents de dues fonts, marcades amb punts negres, que emeten en sincronia. Esquerra: Funció d'ona en un instant determinat de temps.  Dreta: Representació de la fase i amplitud en cada punt. La fase està representada pel color, mentre que l'amplitud correspon a la claror/foscor. Es poden apreciar les hipèrboles d'interferència constructiva i destructiva.}
    \label{fig:Interferencia}
\end{figure}
%
El fenomen d'interferència és especialment aparent quan les ones es propaguen en més d'una dimensió espacial. Si dues fonts en posicions $\vec r_1$ i $\vec r_2$ emeten ones amb la mateixa fase, la interferència en un punt $\vec r$ serà totalment constructiva quan la diferència de distàncies a una i altra font sigui un nombre enter de longituds d'ona, és a dir
%
\begin{equation}
    |\vec r - \vec r_1| - |\vec r - \vec r_2| = n \lambda\, ,\quad  n \in \mathbb{Z}\, .
    \label{eqn:InterferenciaConstructiva}
\end{equation}
%
Per altra banda, la interferència serà totalment destructiva quan la diferència de distàncies a una i altra font sigui un nombre senar (imparell) de semilongituds d'ona, de manera que les dues ones es cancel·lin perfectament, com
%
\begin{equation}
    |\vec r - \vec r_1| - |\vec r - \vec r_2| = (2n -1) \frac{\lambda}{2}\, ,\quad  n \in \mathbb{Z}\, .
    \label{eqn:InterferenciaDestructiva}
\end{equation}
%
Els lloc geomètrics que correspon a les equacions (\ref{eqn:InterferenciaConstructiva}) i (\ref{eqn:InterferenciaDestructiva}) són hipèrboles (o hiperboloides en 3D). Aquestes hipèrboles es poden veure en forma de franges a la Figura \ref{fig:Interferencia}.

\subsection{Modes d'oscil·lació: harmònics}

La formació d'ones estacionàries en una corda subjecta pels dos extrems produeix un efecte molt curiós: la corda no pot sonar amb qualsevol freqüència, sinó que només pot produir unes notes molt concretes! Aquest efecte, més enllà del seu ús en els instruments musicals, és d'una importància enorme en la ciència moderna. És el responsable, per exemple, de que la llum làser s'emeti amb un ``color'' molt concret. També, és l'origen dels nivells d'energia dels electrons en àtoms i molècules, i per tant de tota la química. Els famosos orbitals atòmics són les ``notes'' que ressonen en la funció d'ona quàntica!

En el cas d'una corda subjecta pels dos extrems, els extrems no poden oscil·lar, estan atrapats. Per tant, per força han de coincidir exactament amb dos dels nodes. Això fa que només hi puguin ``ressonar'' una llista de longituds d'ona, que es coneixen com a modes d'oscil·lació, o ``harmònics'' de la corda. A la Figura \ref{fig:ModesExtremsFixes} es veuen representats aquests primers sis harmònics. 
%
\begin{figure}[H]
    \centering
    \includegraphics[width=0.9\linewidth]{Figures/ModesExtremsFixes.pdf}
    \caption{Esquema dels primers sis harmònics d'una corda fixa pels dos extrems.}
    \label{fig:ModesExtremsFixes}
\end{figure}
%
La distància entre dos nodes consecutius és $\lambda/2$. Així doncs, si la longitud total de la corda és $L$ les longituds d'ona permeses $\lambda_n$ seran
%
\begin{equation}
    \lambda_n = \frac{2}{n} L\, , \qquad \qquad n = 1, 2, 3, \dots \qquad \text{(Corda fixa pels dos extrems)}
\end{equation}
%
Com que totes les ones es propaguen a la mateixa velocitat $v$ per la corda, les freqüències corresponents es relacionen amb la del primer harmònic (anomenat harmònic fonamental) com
%
\begin{equation}
    f_n = \frac{v}{\lambda_n} = n f_1 \, , \quad n = 1, 2, 3, \dots  \qquad \text{(Corda fixa pels dos extrems)}
\end{equation}
%
Com podem veure, les freqüències de tots els harmònics són múltiples enters de la freqüència del mode fonamental ($n = 1$). Podem pensar que $n$ és el nombre de ventres que hi ha a la corda.
%
\begin{figure}[H]
    \centering
    \includegraphics[width=\linewidth]{Figures/PentagramaHarmonics.pdf}
    \caption{Representació en un pentagrama dels primers 16 harmònics d'una corda fixa pels dos extrems. Hem assumit que l'harmònic fonamental és un do$_1$, amb una freqüència aproximada de 65,41 Hz.}
    \label{fig:PentagramaHarmonics}
\end{figure}
%
En un instrument musical, l'harmònic fonamental és la nota que volem que soni quan vibrar la corda. Tot i això, no només sonarà la freqüència fonamental sinó que també s'excitaran els harmònics superiors, amb diferents amplituds. La combinació de totes aquestes freqüències sonant alhora és el que dóna el timbre característic a cada instrument, i fan que un violí soni diferent a una flauta, tot i que toquin la mateixa nota. A la Figura \ref{fig:PentagramaHarmonics} es pot veure la correspondència dels harmònics successius d'una corda amb les notes musicals que produeixen.

En el cas que la corda tingui un extrem lliure, aquest haurà de correspondre a un ventre enlloc d'un node. Això ens força a tenir $L = n \lambda_n / 2 + \lambda_n / 4 $, on l'últim terme correspon a la distància entre un node i un ventre. Ens queda  
%
\begin{equation}
    \lambda_n = \frac{4}{2n - 1} L\, , \qquad \qquad n = 1, 2, 3, \dots \qquad \text{(Corda lliure per un extrem)}
\end{equation}
%
Novament,com que totes les ones es propaguen a la mateixa velocitat $v$ per la corda, les freqüències corresponents es relacionen amb la del primer harmònic (anomenat harmònic fonamental) com
%
\begin{equation}
    f_n = \frac{v}{\lambda_n} = (2n - 1) f_1 \, , \quad n = 1, 2, 3, \dots  \qquad \text{(Corda lliure per un extrem)}
\end{equation}
%
En aquest cas, les freqüències de tots els harmònics són múltiples senars (imparells) de la freqüència del mode fonamental ($n = 1$). A la Figura \ref{fig:ModesExtremLliure} es veuen representats els primers sis harmònics amb un extrem lliure. 
%
\begin{figure}[H]
    \centering
    \includegraphics[width=0.9\linewidth]{Figures/ModesExtremsLliures.pdf}
    \caption{Esquema dels primers sis harmònics d'una corda amb un extrem lliure.}
    \label{fig:ModesExtremLliure}
\end{figure}
%
El cas amb els dos extrems lliures és molt semblant al cas amb dos extrems fixos, i normalment no s'estudia separadament. Això passa perquè podem passar de l'un a l'altre senzillament desplaçant l'ona una longitud $\lambda/4$ cap a l'esquerra o la dreta. Per tant, les freqüències dels harmònics amb dos extrems lliures són les mateixes que amb els dos extrems fixos. Quan estudiem les ones sonores a la propera secció veurem que el desplaçament de l'aire té precisament una fase de $\pi/2$ respecte la seva pressió, de manera que nodes i ventres són intercanviables depenent de la magnitud que considerem. Podem resumir tots els casos com:
%
\begin{itemize}
    \item Dos extrems fixos: $\qquad \qquad \; f_n = n f_1, \qquad \qquad f_1 = \frac{v}{2L}$.
    \item Dos extrems lliures: $\qquad \quad \;\; f_n = n f_1, \qquad \qquad f_1 = \frac{v}{2L}$.
    \item Un extrem fix i un de lliure: $f_n = (2n - 1) f_1, \quad f_1 = \frac{v}{4L}$.
\end{itemize}
%
Els harmònics superiors són múltiples del fonamental quan els dos extrems són iguals (fix-fix o lliure-lliure), mentre que són múltiples senars quan els extrems són diferents (fix-lliure o lliure-fix). L'harmònic fonamental $f_1$ té la meitat de freqüència en el cas d'extrems diferents.

\section{Ones sonores}

Les ones sonores són un cas particular d'ones mecàniques, on la matèria que es desplaça ja no és una corda, sinó el medi que ens envolta. Aquest medi és normalment l'aire en el cas dels éssers humans, però també pot ser l'aigua del mar, o fins i tot la roca en el cas de les ones sísmiques. És important estudiar les ones sonores separadament, perquè juguen un paper molt important en la nostra vida quotidiana i en la tecnologia. Formen el so que sentim amb les nostres orelles, i són el canal de transmissió de la parla, la música, i els mecanismes d'ecolocalització d'una gran varietat d'animals. També són utilitzades en aplicacions biomèdiques de diagnòstic per la imatge com l'ecografia, i la teràpia per ones de xoc. 
%
\begin{figure}[H]
    \centering
    \includegraphics[width=\linewidth]{Figures/CompressioRarefaccio.pdf}
    \caption{Comparació entre l'ona de desplaçament $s(x,t)$ de l'aire i l'ona de pressió $P(x,t)$ en el so. A les regions on els desplaçaments convergeixen, es produeix una compressió, és a dir un màxim de pressió. En canvi, les regions on els desplaçaments divergeixen s'anomenen rarefaccions, i corresponen als mínims de pressió. Això es tradueix en un desfasament d'un quart de cicle entre les dues ones.}
    \label{fig:CompressioRarefaccio}
\end{figure}
%
Quan estudiem una ona sonora podem considerar, com a mínim, dues magnituds diferents com a funció d'ona. Ens podem fixar en el desplaçament de l'aire $s(x, t)$ o en la pressió de l'aire $P(x, t)$. En el primer cas, és útil imaginar que hi ha partícules de pols en suspensió. Quan l'ona sonora travessi l'aire, una partícula que inicialment estava en una posició $x = x_0$, oscil·larà molt lleugerament de manera que la seva posició sigui $x = x_0 + s(x, t)$. Per intensitats sonores quotidianes, aquesta oscil·lació tindrà una amplitud minúscula, imperceptible a ull nu. Per altra banda, a causa d'aquests petits desplaçaments de l'aire, la partícula de pols experimentarà petites variacions de pressió de la forma $P(x, t)$. És a dir, en els llocs on l'aire s'hagi acumulat a causa dels petits desplaçaments es produirà una {\bf compressió} i la pressió serà una mica més gran. Per altra banda, els llocs on l'aire s'hagi enrarit s'anomenen punts de {\bf rarefacció} i la pressió serà una mica menor. Si observem la Figura \ref{fig:CompressioRarefaccio}, podem veure que els màxims de $P(x, t)$ estan desplaçats un quart de cicle respecte els màxims de $s(x, t)$.

Observant la representació gràfica a la Figura \ref{fig:CompressioRarefaccio}, es veu que els màxims de $s(x, t)$ corresponen al màxim desplaçament de les partícules cap a la dreta, mentre que els mínims de $s(x, t)$ corresponen al màxim desplaçament de les partícules cap a l'esquerra. Això provoca que els indrets que tenen un màxim de $s(x, t)$ a l'esquerra i un mínim de $s(x, t)$ a la dreta pateixin una acumulació de partícules, que ocasiona una compressió, i per tant un màxim de $P(x, t)$. D'igual manera, els punts amb un mínim de $s(x, t)$ a l'esquerra i un màxim de $s(x, t)$ a la dreta són les rarefaccions, o mínims de $P(x,t)$. Això es pot escriure matemàticament afegint un desfasament d'un quart de cicle entre $s(x, t)$ i $P(x, t)$, com
%
\begin{equation}
\begin{split}
    s(x, t) &= s_0 \sin\left(k x - \omega t + \phi \right)\, , \\
    P(x, t) &= P_0 \sin\left(k x - \omega t + \phi - \frac{\pi}{2}\right)\, .
\end{split}
\end{equation}

Això implica que podem fer aplicar el mateix formalisme matemàtic dels harmònics d'una corda a les ones sonores en un tub, com per exemple en un instrument de vent. En aquest cas tindrem ones sonores en un tub, que pot tenir extrems oberts o tancats. Un extrem obert haurà de coincidir amb un node de l'ona de pressió $P(x, t)$, que equival a un ventre de l'ona de desplaçament $s(x, t)$. En canvi, un extrem tancat serà un ventre de pressió i un node de desplaçament. Per tant, els harmònics seran:
%
\begin{itemize}
    \item Dos extrems oberts: $\qquad \qquad \quad \;\; f_n = n f_1, \qquad \qquad f_1 = \frac{v}{2L}$.
    \item Dos extrems tancats: $\qquad \qquad \quad f_n = n f_1, \qquad \qquad f_1 = \frac{v}{2L}$.
    \item Un extrem obert i un de tancat: $f_n = (2n - 1) f_1, \quad f_1 = \frac{v}{4L}$.
\end{itemize}
%
Fixem-nos que, de manera exactament igual al cas de la corda vibrant, els harmònics superiors són múltiples del fonamental quan els dos extrems són iguals (obert-obert o tancat-tancat), mentre que són múltiples senars quan els extrems són diferents (obert-tancat o tancat-obert). L'harmònic fonamental $f_1$ té la meitat de freqüència en el cas d'extrems diferents, tal com es veu a la Figura \ref{fig:PicsHarmonics}.
%
\begin{figure}[H]
    \centering
    \includegraphics[width=0.9\linewidth]{Figures/PicsHarmonics.pdf}
    \caption{Pics d'amplitud corresponents als harmònics d'un tub, que podríem detectar amb un analitzador d'espectre. Els blaus corresponen als que veuríem en un tub amb els extrems iguals (obert-obert, o tancat-tancat), mentre que els vermells apareixen quan els extrems del tub són diferents (obert-tancat o tancat-obert). És important adonar-se que entre dos pics consecutius sempre hi ha una distància $\Delta f = v_\text{so}/ 2L$, on $L$ és la longitud del tub.}
    \label{fig:PicsHarmonics}
\end{figure}

\subsection{Velocitat del so i intensitat sonora}

L'expressió per la velocitat del so en un material (normalment un fluid) té una forma semblant a la de les ones en una corda: és l'arrel quadrada de la dificultat de deformar el material dividida entre la densitat de massa. La podem expressar com
%
\begin{equation}
    \boxed{v_\text{so} = \sqrt{\frac{K}{\rho}}}
\end{equation}
%
on $K$ és el mòdul de compressibilitat (l'invers de la compressibilitat $\kappa$) i $\rho$ és la densitat. En el cas de l'aire, es pot fer servir l'aproximació d'un gas ideal i reescriure l'equació com
%
\begin{equation}
    v_\text{aire} = \sqrt{\gamma \frac{P}{\rho}}\, ,
\end{equation}
%
on $P$ és la pressió atmosfèrica i $\gamma = \frac{7}{5} = $1,4 s'anomena l'índex adiabàtic per a gasos diatòmics (com el nitrogen i l'oxigen de l'aire). En condicions normals, la velocitat del so en l'aire és de 344 m/s, o uns 1238 km/h. És a dir, el so viatja un quilòmetre cada tres segons aproximadament, cosa que es pot fer servir per estimar la distància a una tempesta comptant el temps des que observem un llamp fins que en sentim el tro. A la Taula \ref{tab:VelocitatDelSo} es poden veure exemples de la velocitat del so en altres medis.
%
\begin{table}[H]
    \centering
    \begin{tabular}{|l|c|c|c|}
    \hline
              & $\rho$ & $K$& $v_\text{so}$\\ \hline
        Aire  & 1,2 kg/m$^3$ & 1,43$\cdot 10^5$ Pa & 344 m/s\\ \hline
        Aigua & 1000 kg/m$^3$ & 2,0$\cdot 10^9$ Pa & 1500 m/s\\ \hline
        Sang  & 1056 kg/m$^3$ & 2,6$\cdot 10^9$ Pa & 1570 m/s\\ \hline
        Ferro & 7900 kg/m$^3$ & 2,0$\cdot 10^{11}$ Pa & 5120 m/s\\ \hline
    \end{tabular}
    \caption{Velocitats del so en diferents medis. Observem que el so es propaga molt més ràpidament en líquids i sòlids que en gasos.}
    \label{tab:VelocitatDelSo}
\end{table}
%
Les ones sonores, com qualsevol altra ona, transporten una certa energia per unitat de temps. És a dir, una font sonora (per exemple una persona parlant) està generant una potència $\mathcal{P}$, mesurada en Watts. Una altra persona que l'estigui escoltant està absorbint una part d'aquesta potència, concretament la que arribi a l'àrea del timpà de la seva orella. Per això, sovint és útil parlar de l'{\bf intensitat sonora} $I$, és a dir la potència per unitat de superfície que transporta una ona (mesurada en Watts per metre quadrat). Com que vivim en un espai tridimensional, quan el so s'hagi allunyat fins a una distància $r$ de la font, la seva potència s'haurà distribuït uniformement sobre una esfera de superfície $S = 4 \pi r^2$, de manera que
%
\begin{equation}
    \boxed{I = \frac{\mathcal{P}}{S} = \frac{\mathcal{P}}{4 \pi r^2}}
\end{equation}
%
Fixem-nos que la intensitat sonora es redueix de manera inversament proporcional al quadrat de la distància a la font. Això és un fenomen que trobem en qualsevol tipus de transmissió d'energia, també en la llum i la radiació provinent de materials radioactius. També passa, per raons similars, amb la intensitat dels camps elèctrics i gravitatoris. En el cas de les ones sonores, la intensitat es relaciona amb l'amplitud de les oscil·lacions de pressió com
%
\begin{equation}
    \boxed{I = \frac{P_0^2}{2\rho v_\text{so}}}
\end{equation}
%
que, en unitats del Sistema Internacional, equival a
%
\begin{equation}
    I = 0,0012 \, P_0^2 \qquad \text{(Unitats del SI)}\, .
\end{equation}
%

\subsection{Sensació sonora: el decibel}

Els sons que rebem en la nostra vida quotidiana tenen un enorme rang d'intensitats. Hi ha sons extremadament febles, com el de la nostra pròpia respiració (uns $10^{-11}$ W/m$^2$), i sons extraordinàriament potents, com el motor a reacció d'un avió comercial (aproximadament 1 W/m$^2$). Això és una potència 10$^{11}$ vegades més gran! Si la nostra percepció de la intensitat d'un soroll fos proporcional a la potència, seria molt poc pràctic per a la nostra vida. És molt difícil tractar amb una sensació molts ordres de magnitud més gran que una altra. 

Per això, la nostra sensació sonora augmenta cada cop més lentament a mesura que la potència que rebem es fa més gran. D'aquesta manera, cada cop que la potència que ens arriba es multiplica per 10, la nostra sensació augmenta una quantitat similar. És a dir, per a nosaltres, augmentar des de $10^{-11}$ W/m$^2$ fins a $10^{-10}$ W/m$^2$ es percep com un augment de volum similar a passar augmentar des de $0.1$ W/m$^2$ fins a $1$ W/m$^2$. Encara que en el segon cas l'augment absolut sigui molt més gran, l'augment relatiu és el mateix (un factor 10): és el que es coneix com una escala logarítmica. D'aquesta manera, definim una magnitud anomenada sensació sonora $\beta$ com
%
\begin{equation}
    \boxed{\beta = 10 \log_{10} \left(\frac{I}{I_0}\right)}
\end{equation}
%
on $I$ és la intensitat sonora en W/m$^2$, i $I_0 = 10^{-12}$ W/m$^2$ és una intensitat de referència, que es considera el llindar de percepció humana. La unitat de sensació sonora $\beta$ és el decibel (dB), molt utilitzat tant en ciència com en situacions quotidianes. D'aquesta definició, es pot veure que cada cop que la intensitat sonora es {\bf multiplica} per 10, a la sensació sonora se li {\bf sumen} 10 dB. El logaritme té aquesta propietat de convertir multiplicacions en sumes. A la Taula \ref{tab:TaulaSensacionsSonores} es poden veure alguns exemples de la sensació sonora associada a diferents sons.
%
\begin{table}[H]
    \centering
    \begin{tabular}{|l|c|c|c|}
    \hline
                          & Distància & $\beta$ & Apreciació  \\ \hline
        Motor a reacció     & 50 m  & 140 dB & llindar del dolor \\ \hline
        Mascletà valenciana & 50 m  & 130 dB & extremadament fort \\ \hline
        Martell pneumàtic   & 10 m  & 110 dB &  \\ \hline
        Crit humà           & 1,5 m & 100 dB & molt fort \\ \hline
        Clàxon              & 10 m  & 90 dB  &  \\ \hline
        Secador de cabell   & 20 cm & 80 dB  & fort \\ \hline
        Cotxe en marxa      &       & 70 dB  &  \\ \hline
        Conversa            & 1 m   & 60 dB  & moderat \\ \hline
        Oficina             &       & 50 dB  &  \\ \hline
        Biblioteca          &       & 40 dB  &  \\ \hline
        Dormitori tranquil  &       & 30 dB  & fluix \\ \hline
        Xiuxiueig           & 1 m   & 20 dB  &  \\ \hline
        Respiració humana   & 5 m   & 10 dB  & gairebé inaudible \\ \hline
    \end{tabular}
    \caption{Sensacions sonores, en decibels, corresponents a diverses situacions quotidianes. Cada augment de 10 dB correspon a multiplicar per 10 la intensitat sonora.}
    \label{tab:TaulaSensacionsSonores}
\end{table}
%
Aquesta particularitat de la percepció en escala logarítmica dels nostres sentits no és única en la sensació sonora. De fet, moltes de les nostres percepcions de l'ambient són també en escala logarítmica, per tal de poder assimilar variacions enormes en la magnitud física que percebem. Un altre exemple típic és el de la freqüència sonora de les notes musicals. Cada cop que dupliquem la freqüència, és a dir la multipliquem per 2, la nostra percepció és que hem augmentat en una octava. És a dir, percebem el mateix canvi en l'escala musical en passar de 440 Hz (la$_3$) fins a 880 Hz (la$_4$), que en passar de 880 Hz (la$_4$) a 3520 Hz (la$_5$), tot i que el canvi de freqüència és més gran en el segon cas.

\subsection{Efecte Doppler}

Finalment, un efecte molt important en les ones sonores és l'aparent canvi de freqüència (o de nota musical) que percebem quan l'emissor o el receptor es mouen. Això es nota molt sovint quan un avió, o un cotxe, passa prop nostre a gran velocitat. Mentre que la persona pilotant l'avió sent una nota aproximadament constant durant tota l'estona, nosaltres sentim una nota més aguda quan s'acosta i una nota més greu quan s'allunya. Això produeix el característic so ``{\it nyiiiiiiiaoooooooo}'', amb la baixada ràpida de to en el moment que l'avió ens avança.

Per entendre aquest fenomen, és útil introduir el concepte de {\bf front d'ona}. Podem pensar en els fronts d'ona com els llocs on la pressió es troba en un màxim. D'aquesta manera, quan l'avió emet un soroll, els fronts d'ona són esferes concèntriques que augmenten de radi a la velocitat del so. Aquests fronts esfèrics tenen el seu centre en el lloc des del qual s'han emès. Si l'avió està aturat, això tindria un aspecte similar al que es veu a la Figura \ref{fig:Doppler0}.
%
\begin{figure}[H]
    \centering
    \includegraphics[width=0.4\linewidth]{Figures/Doppler0.pdf}
    \caption{Fronts d'ona emesos per una font sonora, en aquest cas un avió en marxa, quan està aturat. Els fronts són esferes concèntriques centrades en l'emissor.}
    \label{fig:Doppler0}
\end{figure}
%
En canvi, si l'avió està volant a una certa velocitat, està perseguint els seus propis fronts d'ona. Això farà que els centres dels diversos fronts estiguin desplaçats entre ells, i per tant que es trobin més junts davant de l'avió i més separats darrere seu. Una persona que es trobi davant seu els rebrà molt junts, i per tant sentirà una freqüència més alta, mentre que una persona que es trobi darrere de l'avió rebrà els fronts més separats, és a dir una freqüència més baixa. Això és el que s'ha intentat representar a la Figura \ref{fig:Doppler08}.
%
\begin{figure}[H]
    \centering
    \includegraphics[width=0.7\linewidth]{Figures/Doppler08.pdf}
    \caption{Fronts d'ona que exemplifiquen l'efecte Doppler, quan l'emissor i el receptor es mouen l'un cap a l'altre. A causa del moviment de l'emissor, els fronts estan més junts davant seu que darrere seu. Això farà que el receptor, l'altre avió, percebi un so més agut. Un cop emissor i receptor s'hagin creuat, el receptor rebrà els fronts d'ona més separats, i per tant un so més greu. }
    \label{fig:Doppler08}
\end{figure}
%
Per tal de fer els càlculs, anem a suposar que l'emissor (l'avió) es mou a una velocitat $v_E$ en direcció al receptor, mentre que el receptor es mou amb velocitat $v_R$ en direcció a l'emissor. Si algun d'aquests elements es mogués en sentit contrari, simplement canviaríem el signe de la seva velocitat, com explicarem més endavant. En primer lloc, calcularem la distància $d_f$ que hi ha entre els fronts d'ona entre ells. Si l'emissor emet una oscil·lació cada període $T_E = 1/f_E$, la distància entre fronts serà
%
\begin{equation}
    d_f = (v_\text{so} - v_E) T_E\, , 
\end{equation}
%
és a dir, la distància que el front ha recorregut durant el temps $T_E$ menys la distància que l'emissor ha avançat durant aquest temps. Per altra banda, el receptor es desplaça a velocitat $v_\text{so} + v_R$ respecte als fronts, i per tant els rebrà amb un període $T_R = 1/f_R$ donat per
%
\begin{equation}
    T_R = \frac{d_f}{v_\text{so} + v_R} = \frac{v_\text{so} - v_E}{v_\text{so} + v_R} T_E\, .
\end{equation}
%
Reescrivint aquesta expressió en termes de la freqüència, per a emissor i receptor movent-se l'un cap a l'altre, ens queda
%
\begin{equation}
    f_R = \frac{v_\text{so} + v_R}{v_\text{so} - v_E} f_E\, .
\end{equation}
%
Podem generalitzar ara per signes arbitraris de les velocitats, i escriure l'equació com 
%
\begin{equation}
    \boxed{f_R = \frac{v_\text{so} \pm v_R}{v_\text{so} \mp v_E} f_E}
\end{equation}
%
Els signes del numerador i denominador, respectivament, dependran de la direcció en la que es moguin emissor i receptor entre ells. A continuació es mostra un resum de casos:
%
\begin{itemize}
    \centering
    \item $\; E \longrightarrow \quad \longleftarrow R \quad \bullet$ : $(+/-)$
    \item $\; E \longrightarrow \quad \longrightarrow R \quad \bullet$ : $(-/-)$
    \item $\; E \longleftarrow \quad \longleftarrow R \quad \bullet$ : $(+/+)$
    \item $\; E \longleftarrow \quad \longrightarrow R \quad \bullet$ : $(-/+)$
\end{itemize}
%
Un cas molt curiós es dóna quan la velocitat de l'emissor és la velocitat del so ($v_E = v_\text{so}$). En aquest cas la fórmula ens diu que $f_R \to \infty$. Això el que vol dir és que els fronts d'ona estan avançant exactament a la mateixa velocitat que l'avió, i es van acumulant l'un sobre l'altre, com es veu a la Figura \ref{fig:Doppler1}. 
%
\begin{figure}[H]
    \centering
    \includegraphics[width=0.6\linewidth]{Figures/Doppler1.pdf}
    \caption{Fronts d'ona al voltant d'un avió viatjant exactament a la velocitat del so. Podem observar que tots els fronts s'han acumulat a la part del davant de l'avió, cosa que generarà el bang sònic. Una persona que es trobi a terra no sentirà cap soroll fins que l'avió no arribi, moment en el qual sentirà el bang sònic. }
    \label{fig:Doppler1}
\end{figure}
%
D'aquesta manera, algú que estigui davant de l'avió no sentirà res fins que l'avió no arribi, moment en el qual rebrà tots els fronts d'ona de cop en forma d'{\bf ona de xoc}. Això es percep com un soroll molt fort, semblant a una explosió, conegut com a {\bf bang sònic}. Aquest so és extremadament potent i pot trencar fàcilment els vidres dels edificis, per això els avions eviten volar a velocitats supersòniques prop de zones habitades.

Finalment, si l'avió vola més ràpid que la velocitat del so ($v_E > v_\text{so}$), avançarà contínuament els seus propis fronts d'ona. D'aquesta manera, només hi haurà ones sonores en una regió cònica al seu darrere, anomenada {\bf con de Mach}. En fotografies d'avions supersònics, aquest con pot produir la condensació de l'aire i veure's com un núvol cònic que persegueix l'avió. Tal com es veu a la Figura \ref{fig:Doppler2}, fent servir trigonometria bàsica podem deduir que l'angle del con es relaciona amb la velocitat de l'emissor com
%
\begin{equation}
    \boxed{\sin \theta = \frac{v_\text{so}}{v_E} = \frac{1}{M}}
\end{equation}
%
on el nombre $M = v_E / v_\text{so}$ s'anomena nombre de Mach, i ens indica la velocitat de l'avió en unitats de la velocitat del so. 
%
\begin{figure}[H]
    \centering
    \includegraphics[width=0.6\linewidth]{Figures/Doppler2.pdf}
    \caption{Con de Mach darrere un avió viatjant a $v = 2v_\text{so}$ (Mach 2). En aquest cas l'avió està avançant els seus propis fronts d'ona, de manera que es genera una ona de xoc de forma cònica darrere seu, on està contingut tot el soroll. La regió fora del con de Mach està encara en silenci.}
    \label{fig:Doppler2}
\end{figure}
%

Per exemple, un avió viatjant a Mach 2 es mou a dues vegades la velocitat del so. En vols comercials, sobretot els transoceànics, és comú que la velocitat s'indiqui a les pantalles de la cabina com a ``Mach 0,8'', per exemple. Els antics avions Concorde transportaven passatgers a Mach 2, mentre que la llançadora espacial de la NASA arriba a Mach 27. Com a curiositat, el primer objecte de fabricació humana que va superar la velocitat del so és el fuet. La punta d'aquest objecte, quan es fa servir correctament, acumula una quantitat immensa d'energia que fa que assoleixi velocitats supersòniques. L'espetec característic d'aquesta arma correspon al bang sònic que produeix la punta. 

Curiosament, l'efecte del bang sònic no és exclusiva de les ones sonores, sinó que té un anàleg en les ones electromagnètiques, com la llum. A priori, això pot semblar impossible, ja que la teoria de la relativitat prohibeix que cap objecte superi la velocitat de la llum en el buit. La clau està en aquesta última part: en el buit. Com veurem en el pròxim capítol, la llum es propaga més lentament a través de materials transparents com l'aigua o el vidre. En aquests medis, sí que és possible que una partícula molt energètica es desplaci més ràpid que la llum en el medi, tot i que segueixi essent menor que la del buit. En aquest cas, l'equivalent al bang sònic s'anomena {\bf radiació de Cherenkov}, i es pot veure en filmacions de reactors nuclears submergits en aigua. La gran quantitat de partícules que es desprenen de la reacció nuclear formen una llum blava característica. Aquesta brillantor blava també s'ha observat en accidents de criticalitat, i és un indicador de nivells de radiació catastròfics per als éssers vius.

\subsection{Ecolocalització i percepció del so}

Molts animals, com ara els rat-penats o els dofins, fan servir ones sonores per orientar-se en l'espai i identificar les preses. Aquest mètode s'anomena {\bf ecolocalització}, i consisteix en enviar una senyal sonora en forma de crit, i mesurar el temps que triga l'ona reflectida (l'eco) en tornar. Mentre més temps trigui, més lluny està l'objecte que volen detectar. També, en alguns casos, poden utilitzar l'efecte Doppler per tal de deduir la velocitat a la qual es mou una presa. Aquest principi s'ha explotat tecnològicament en els sonars de vaixells i submarins.

Els éssers humans podem sentir típicament sons entre 20 Hz i 20 kHz. Els sons de freqüències més baixes de 20 Hz s'anomenen infrasons, i els de freqüències més altes de 20 kHz s'anomenen ultrasons. Tot i això, la màxima freqüència audible disminueix ràpidament amb l'edat. Altres animals poden emetre i sentir sons fora d'aquest rang. Algunes balenes poden comunicar-se en infrasons, mentre que la majoria d'animals que utilitzen ecolocalització treballen en ultrasons. Alguns rat-penats poden arribar als 200 kHz, una freqüència totalment inaudible pels éssers humans.


\chapter{Ones electromagnètiques i òptica}
\label{chap:OnesElectromagnetiques}

\section{Electromagnetisme i llum}

Les ones electromagnètiques són lleugerament més complexes del que hem estudiat fins ara, però són probablement les més importants a la vida de les persones, a la ciència i a la tecnologia. L'exemple més quotidià és la llum, que ens permet percebre imatges amb els nostres ulls. També són ones electromagnètiques les ones de ràdio, les microones, la radiació infraroja, la radiació ultravioleta, els raigs X i la radiació gamma. Veurem un resum de totes aquestes parts de l'espectre properament.
%
\begin{figure}[H]
    \centering
    \includegraphics[width=0.9\linewidth]{Figures/OnaEM.pdf}
    \caption{Esquema dels camps elèctric $\vec E$ (en vermell) i el camp magnètic $\vec H$ (en blau) en una ona electromagnètica plana i linealment polaritzada.}
    \label{fig:OnaEM}
\end{figure}
%
En general, i salvant certs matisos, la radiació electromagnètica es genera a partir del moviment de partícules carregades elèctricament. En aquest cas, ens trobem davant d'una ona transversal, en la que oscil·len simultàniament dos camps vectorials: el camp elèctric ($\vec E$) i el camp magnètic ($\vec H$). Aquests camps oscil·len en fase, però en direccions perpendiculars entre ells, i perpendiculars també a la direcció de propagació. Això és pot veure més clarament a la Figura \ref{fig:OnaEM}, que mostra el que s'anomena una ona plana linealment polaritzada (el cas més senzill). La direcció de propagació de l'ona, i el flux d'energia, venen donats pel {\bf vector de Poynting}
%
\begin{equation}
    \vec S = \vec E \times \vec H\, ,
\end{equation}
%
on $\mu_0$ és una constant que s'anomena permeabilitat del buit. Tot i que ara és un coneixement bastant universal, la naturalesa electromagnètica de la llum no es va conèixer fins a mitjan segle XIX. Aquest descobriment va ser conseqüència del treball per part de James Clerk Maxwell, qui va publicar el 1865 una sèrie d'equacions per descriure l'electricitat i el magnetisme. Aquestes {\bf equacions de Maxwell}, més endavant simplificades per Oliver Heaviside, es troben entre les més famoses de la física. 

Un cop formulades les seves equacions, Maxwell es va adonar que permetien solucions en forma d'ona. No només això, sinó que va ser capaç de calcular la velocitat d'aquesta ona en termes de constants físiques ja conegudes, anomenades permitivitat elèctrica del buit ($\varepsilon_0$) i permeabilitat magnètica del buit ($\mu_0$). Aquesta velocitat va resultar ser
%
\begin{equation}
    c = \frac{1}{\sqrt{\varepsilon_0 \mu_0}} = 2,99792458 \cdot 10^8 \text{ m/s} \, ,
\end{equation}
%
que coincidia perfectament amb el valor ja conegut de la velocitat de la llum, que ja havia estat mesurada el 1676 per Ole Rømer. Davant d'això, Maxwell va proposar per primer cop que la llum té una naturalesa electromagnètica. No va ser fins el 1887 en què Heinrich Hertz va poder generar aquestes ones (en aquest cas ones de ràdio) a partir d'un dispositiu purament electromagnètic.

\subsection{Espectre electromagnètic}

Depenent de la seva freqüència, o longitud d'ona, les diverses ones electromagnètiques tenen diferents propietats i diferents noms. A grans trets, aquests trams de l'espectre són:
%
\begin{itemize}
    \item Ones de ràdio ($\lambda >$ 1 m): Són les ones de longitud d'ona més alta, i freqüència més baixa. Es fan servir en tot tipus de telecomunicacions, com ara xarxes de telefonia mòbil, WiFi, Bluetooth, GPS, ràdio i televisió... També en detecció per radar.
    %
    \item Microones (1 m $ \geq \lambda >$ 100 $\mu$m): També s'utilitzen per telecomunicacions i deteccions per radar. Les ones de $\lambda$ = 12 cm tenen la propietat d'escalfar diverses substàncies, entre elles l'aigua, i per tant es fan servir en forns microones.
    %
    \item Infraroig (IR) (100 $\mu$m $ \geq \lambda >$ 700 nm): Radiació propera a l'espectre visible, però més energètica. Té múltiples aplicacions tecnològiques, com ara els comandaments a distància. És el tipus de radiació tèrmica que emetem els éssers humans i qualsevol objecte a temperatura ambient.
    %
    \item Llum visible (700 nm $ \geq \lambda >$ 400 nm): És la zona de l'espectre electromagnètic que detectem amb els nostres ulls, i que ens permet formar imatges i veure el nostre entorn. Va des del color vermell ($\sim$700 nm) fins al violeta ($\sim$400 nm), passant per tots els colors de l'arc de Sant Martí.
    %
    \item Ultravioleta (UV) (400 nm $ \geq \lambda >$ 10 nm): De longitud d'ona més curta que la llum visible, tot i que és detectable per alguns animals com les abelles. És la responsable de les cremades solars, i es fa servir sovint per eliminar patògens com virus i bacteris.
    %
    \item Raigs X (10 nm $ \geq \lambda >$ 10 pm): És una radiació altament energètica, amb longituds d'ona a l'escala atòmica. Per això, és molt útil en l'exploració de l'estructura molecular. També es fa servir extensivament en el diagnòstic per la imatge en medicina.
    %
    \item Raigs gamma (10 pm $ \geq \lambda$): El cas més energètic. Típicament es produeix en reaccions nuclears i fenòmens astrofísics extrems. Es fa servir en algunes aplicacions mèdiques, com ara la tomografia per emissió de positrons (PET).
\end{itemize}

Les ones a partir de l'ultravioleta es coneixen com a radiacions ionitzants, perquè els seus fotons són suficientment energètics per extraure electrons d'àtoms i molècules. Aquestes radiacions són molt perilloses per als éssers vius, perquè poden alterar l'estructura de biomolècules, en particular les cadenes d'ADN. Això pot inutilitzar les cèl·lules, o produir mutacions cancerígenes.   

\subsection{Òptica geomètrica}

Durant la majoria d'aquest capítol, ens centrarem en l'anàlisi de la llum, el que es coneix com a òptica. Per estudiar la majoria de fenòmens òptics, és molt útil pensar en la propagació de les ones en termes de fronts, que ja hem vist al capítol anterior, i raigs de llum:
%
\begin{itemize}
    \item Fronts d'ona: Són superfícies imaginàries corresponents a una mateixa fase de l'ona. Per exemple, els punts on el camp elèctric està en un màxim. Són perpendiculars a la direcció de propagació de l'ona.
    %
    \item Raigs de llum: Les ``trajectòries'' que segueix la llum, que són perpendiculars als fronts d'ona. És la idea intuïtiva que tenim d'un raig de llum del Sol, o d'un raig làser.
\end{itemize}
%
\begin{figure}[H]
    \centering
    \includegraphics[width=0.6\linewidth]{Figures/FrontsRaigs.pdf}
    \caption{Representació esquemàtica d'una ona esfèrica propagant-se des d'una font lluminosa puntual al centre. Els fronts d'ona es representen de color blau, i els raigs de llum de color vermell. Els raigs sempre són perpendiculars als fronts d'ona.}
    \label{fig:FrontsRaigs}
\end{figure}
%
A la Figura \ref{fig:FrontsRaigs} es veu una representació esquemàtica dels fronts d'ona i els raigs que surten d'una font lluminosa. Podem comprovar que els raigs de llum són perpendiculars als fronts d'ona en tot moment. La distància entre dos fronts d'ona consecutius és d'una longitud d'ona $\lambda$. Una eina  bastant útil i intuïtiva per imaginar la propagació d'una ona a través de l'espai, basant-se en el concepte de front d'ona, és el {\bf principi de Huygens}. Tot i que el seu ús en aquest capítol serà limitat, el principi de Huygens està darrere de fenòmens com la difracció i la formulació de la mecànica quàntica en termes de la integral de camins. Es pot formular així:

\begin{mdframed}
\begin{it}
    Cada punt d’un front d’ona es comporta com emissor d’ones esfèriques secundàries. A cada instant posterior, el front d’ona és l'envolupant de les ones secundàries emeses anteriorment.
\end{it}
\end{mdframed}
%
En altres paraules, podem imaginar que tots els punts d'un front d'ona estan generant, al seu torn, una ona esfèrica que s'allunya d'ells. Totes aquestes petites ones esfèriques són tangents al nou front d'ona al cap d'un temps curt. Aquest efecte està exemplificat a la Figura \ref{fig:PrincipiDeHuygens}.
%
\begin{figure}[H]
    \centering
    \includegraphics[width=0.75\linewidth]{Figures/PrincipiDeHuygens.pdf}
    \caption{Esquema del principi de Huygens, per un front d'ona pla (a l'esquerra) i un front d'ona esfèric (a la dreta). De cada punt del front principal (blau fosc), generem un nou front esfèric secundari (blau clar). L'envolupant d'aquests fronts secundaris forma el nou front principal.}
    \label{fig:PrincipiDeHuygens}
\end{figure}
%
\subsection{Reflexió i refracció}

El formalisme dels fronts d'ona i els raigs de llum ens permet estudiar una gran quantitat de fenòmens òptics a partir d'arguments geomètrics, com la trigonometria. Això s'anomena ``òptica geomètrica'', i ens permet estudiar la llum sense haver de recórrer a les equacions diferencials de les ones. La immensa majoria d'aplicacions òptiques es poden estudiar d'aquesta manera.

El primer que ens podem preguntar és què passa quan la llum troba un canvi de medi durant la seva propagació. És a dir, quin serà el comportament d'un raig de llum (raig incident) en interaccionar amb la superfície de separació entre dos materials, com per exemple aire i aigua, o aire i vidre. Intuïtivament, esperem que part de la llum continuï el seu viatge cap al segon medi (raig refractat), i una altra part es reflecteixi de tornada cap a l'aire (raig reflectit). Però quins seran exactament aquests raigs?
%
\begin{figure}[H]
    \centering
    \includegraphics[width=0.5\linewidth]{Figures/SnellPerpendicular.pdf}
    \caption{Quan una ona passa d'un medi a un altre amb un índex de refracció major, la seva velocitat disminueix, i la distància entre fronts es fa més curta.}
    \label{fig:SnellPerpendicular}
\end{figure}
%
Per poder respondre a aquesta pregunta cal saber que la llum, que viatja a velocitat $c$ en el buit, es propaga més lentament a través de materials transparents (dielèctrics). Això passa a causa de la interacció de les ones electromagnètiques amb l'estructura microscòpica del material, que fan que la velocitat de la llum disminueixi de manera efectiva, depenent de la seva freqüència. Per quantificar aquest relentiment, es defineix {\bf l'índex de refracció} $n$ d'un material com
%
\begin{equation}
    n = \frac{c}{v}\, 
\end{equation}
%
on $c$ és la velocitat de la llum en el buit i $v$ la velocitat de propagació en el material. L'índex de refracció $n$ sempre serà major o igual a 1. Per exemple, l'índex de refracció de l'aigua és proper a 1,33, i el del vidre sol estar prop d'1,5. Quan la llum passa d'un medi a un altre, la seva freqüència es mantindrà, ja que els fronts d'ona arriben a la interfase al mateix ritme amb el que marxen. Altrament, els fronts d'ona s'hi acumularien. Tot i això, si en el segon medi la velocitat de propagació és menor, la separació entre fronts (la longitud d'ona $\lambda$) serà menor: els fronts no han pogut allunyar-se gaire quan arriba el següent front. Això s'exemplifica a la Figura \ref{fig:SnellPerpendicular}.
%
\begin{figure}[H]
    \centering
    \includegraphics[width=0.65\linewidth]{Figures/FrontsDeSnell.pdf}
    \caption{Càlcul de la distància entre les interseccions dels fronts d'ona amb la interfase entre dos medis amb diferents índexs de refracció.}
    \label{fig:FrontsDeSnell}
\end{figure}
%
Anem ara a estudiar el que passa quan un raig de llum arriba a la superfície (o interfase) entre dos medis, amb índexs de refracció $n_1$ i $n_2$, formant un cert angle. Aquest angle, que anomenarem angle d'incidència $\theta_1$, es mesura entre la direcció del raig de llum i la direcció perpendicular a la superfície. Podem observar que l'angle entre la interfase i els fronts d'ona també és $\theta_1$. Llavors, la distància $d_1$ entre els punts d'intersecció de dos fronts amb la interfase serà
%
\begin{equation}
    d_1 = \frac{c T}{n_1 \sin \theta_1}\, ,
\end{equation}
%
on $T$ és el període de l'ona (el mateix per als dos medis) i $n_1$ és l'índex de refracció del primer medi. Això es pot veure dibuixat a la Figura \ref{fig:FrontsDeSnell}. Ara, els fronts d'ona s'han d'unir de forma contínua en els dos medis, de manera que la distància entre les seves interseccions amb la interfase ha de ser la mateixa pels dos costats. En altres paraules, $d_1 = d_2$, és a dir
%
\begin{equation}
    \frac{c T}{n_1 \sin \theta_1} = \frac{c T}{n_2 \sin \theta_2}\, .
\end{equation}
%
Reescrivint aquesta equació obtenim la {\bf llei de Snell}, possiblement la llei més fonamental en l'òptica geomètrica:
%
\begin{equation}
    \boxed{n_1 \sin \theta_1 = n_2 \sin \theta_2}
\end{equation}
%
\begin{figure}[H]
    \centering
    \includegraphics[width=0.6\linewidth]{Figures/RaigsSnell.pdf}
    \caption{Esquema dels angles d'incidència $\theta_1 = \theta_i$, reflexió $\theta_r$ i refracció $\theta_2$.}
    \label{fig:RaigsSnell}
\end{figure}
%
És a dir, els raigs de llum canvien de direcció en creuar la interfase entre dos medis transparents amb diferents índexs de refracció. Aquest fenomen s'anomena {\bf refracció de la llum}. A més, part de la llum es reflectirà de tornada cap al medi original. En aquest cas, l'index de refracció original i el final són el mateix, de manera que l'angle d'incidència $\theta_i = \theta_1$ i el de reflexió $\theta_r$ són el mateix:
%
\begin{equation}
    \boxed{\theta_i = \theta_r}
\end{equation}
%
Això coincideix amb la idea intuïtiva de la reflexió de la llum en un mirall. A la Figura \ref{fig:RaigsSnell} es pot veure un esquema dels tres raigs de llum (incident, reflectit i refractat) en passar de l'aire ($n_1 \approx 1$) a un medi amb $n = $1,7. És important adonar-se que si $n_2 < n_1$ existirà un {\bf angle límit} $\theta_1 = \theta_L$ d'incidència que compleixi 
%
\begin{equation}
    \sin \theta_L = \frac{n_2}{n_1}\, ,
\end{equation}
%
i per tant $\theta_2 = 90^\circ$. És a dir, l'angle refractat surt paral·lelament a la interfase entre medis. Per qualsevol $\theta_1 > \theta_L$, senzillament no existeix cap valor possible per a $\theta_2$, ja que l'equació $\sin \theta_2 > 1$ no té solució. En aquest cas, no hi haurà raig refractat i tota la llum es reflectirà de tornada cap al primer medi. Això passa, per exemple, quan mirem la superfície des de dintre l'aigua. A partir de cert angle, la superfície es veu com un mirall, i no permet veure cap imatge de l'exterior de l'aigua.
%
\begin{figure}[H]
    \centering
    \includegraphics[width=0.65\linewidth]{Figures/DispersioPrisma.pdf}
    \caption{Representació esquemàtica del fenomen de dispersió d'un raig de llum blanca en un prisma. En aquest cas, $n_\text{vermell} < n_\text{blau}$.}
    \label{fig:DispersioPrisma}
\end{figure}
%
Un altre fenomen curiós és el de la {\bf dispersió de la llum}. Atès que l'índex de refracció de la llum en certs materials depèn de la seva freqüència, els diferents colors que formen un raig de llum blanca es desvien de manera diferent en trobar-se amb el material. Això fa possible, per exemple, obtenir tots els colors de l'arc de Sant Martí quan la llum blanca travessa un prisma de vidre, com es veu a la Figura \ref{fig:DispersioPrisma}.

\section{Lents convergents i divergents}

El fet de poder deformar els raigs de llum fent servir materials amb diversos índexs de refracció fa possible una enorme quantitat d'aplicacions tecnològiques que augmentin la nostra capacitat de visió: els instruments òptics. Això inclou ulleres, lentilles, lupes, microscopis, telescopis, i tota mena d'aparells derivats. Aquests instruments fan servir diversos dispositius com prismes i miralls, però el seu element fonamental és la lent. Una lent és, bàsicament, una peça aplanada de vidre o algun material similar limitada per dues superfícies esfèriques. És a dir, les seves dues cares són casquets esfèrics. Per estudiar l'efecte de les lents sobre la llum, farem servir els diagrames de raigs, com els de les Figures \ref{fig:LentConvergent} i \ref{fig:LentDivergent}. En aquesta secció és molt important mantenir un conveni consistent de convencions i criteris de signes, per evitar confusions. La primera convenció és que sempre assumirem que {\bf la llum es desplaça d'esquerra a dreta}, i que l'observador (per exemple el nostre ull) es troba a la dreta del diagrama. La llum emesa, la lent i l'observador estan alineats al llarg d'una línia imaginària anomenada {\bf eix òptic}. Tenim generalment dos tipus de lents: {\bf convergents i divergents}. 
%
\begin{figure}[H]
    \centering
    \includegraphics[width=\linewidth]{Figures/LentConvergent.pdf}
    \caption{En una lent convergent ($f > 0$), els raigs de llum que arriben paral·lels a l'eix òptic es creuen al focus principal $F'$. }
    \label{fig:LentConvergent}
\end{figure}
%
Una lent convergent tendeix a ``focalitzar'' els raigs de llum. Aquestes lents són més gruixudes al centre que a les vores. Els raigs que incideixen paral·lels a l'eix òptic de de l'esquerra passaran tots per un punt a la dreta, anomenat {\bf focus principal $F'$}. La distància entre $F'$ i la lent s'anomena {\bf distància focal} i se simbolitza amb la lletra $f$. Es pot veure un esquema a la Figura \ref{fig:LentConvergent}. A la mateixa distància $f$ però a l'esquerra de la lent definim un altre punt, anomenat {\bf focus secundari $F$}. Els raigs incidents que passin per aquest focus secundari sortiran paral·lels a l'eix òptic per la dreta.
%
\begin{figure}[H]
    \centering
    \includegraphics[width=\linewidth]{Figures/LentDivergent.pdf}
    \caption{En una lent divergent ($f < 0$), els raigs que arriben paral·lels a l'eix òptic no es creuen, però les seves perllongacions cap enrere es troben al focus $F'$.}
    \label{fig:LentDivergent}
\end{figure}
%
Una lent divergent, en canvi, tendeix a ``separar'' els raigs que la travessen. Aquestes lents són més primes al centre que a les vores. En aquest cas, el focus principal $F'$ es troba a l'esquerra de la lent, el costat des d'on incideix la llum. Els raigs que incideixen paral·lels a l'eix òptic es dispersaran de tal manera que semblin provenir de $F'$, és a dir que la seva perllongació cap a l'esquerra passarà pel focus principal. En aquest cas, per convenció, la distància entre la lent i $F'$ (la distància focal) es considera negativa: $f < 0$. En canvi, el focus secundari $F$ es troba a la dreta. Els raigs que arribin a la lent en la direcció de $F$ sortiran paral·lels a l'eix òptic. Es pot veure un esquema a la Figura \ref{fig:LentDivergent}. 

Existeix una expressió, anomenada {\bf equació del constructor de lents}, que ens proporciona la distància focal $f$ d'una lent a partir del seu índex de refracció i dels radis de les seves superfícies. Aquesta equació es pot escriure com
%
\begin{equation}
    \boxed{\frac{1}{f} = (n - 1) \left(\frac{1}{R_1} - \frac{1}{R_2}\right)}
\end{equation}
%
En aquesta equació, $n$ és l'índex de refracció del vidre de la lent i $R_1$ i $R_2$ són els radis dels casquets esfèrics que en formen les superfícies. Novament, cal seguir un criteri de signes. En aquest cas, considerarem que un radi és positiu quan es corba cap a l'esquerra en forma de C, i negatiu quan es corba cap a la dreta en forma de D. Més endavant veurem un resum detallat de tots aquests convenis. Si la lent es troba submergida en un medi amb índex de refracció $n_\text{medi} \neq 1$, hem d'utilitzar l'equació del constructor de lents amb un índex de refracció relatiu $n = n_\text{lent}/n_\text{medi}$. L'invers de la distància focal s'anomena potència de la lent, com
%
\begin{equation}
    p = \frac{1}{f} \qquad \text{(Mesurada en diòptries)}\, .
\end{equation}
%
Si la focal està mesurada en metres, les unitats de la seva potència són les {\bf diòptries}, de símbol dp, molt utilitzades en aplicacions pràctiques. A més, a nivell pràctic és molt útil quan situem dues lents primes de potències $p_1$ i $p_2$ molt a prop l'una de l'altra, ja que el conjunt actua com una sola lent de potència $p_T = p_1 + p_2$. És a dir, en la composició de lents, les seves potències se sumen. És molt habitual utilitzar les diòptries en les lents de les ulleres de correcció de la miopia i la hipermetropia. Per corregir la miopia s'utilitzen lents divergents, de potència negativa (per exemple, -3 dp), mentre que per corregir la hipermetropia s'utilitzen lents convergents.

\subsection{Diagrames de raigs}

Quan observem un objecte a través d'una lent, els raigs de llum procedents de cada un dels seus punts es desviaran, i semblaran procedir d'un punt diferent. Això farà que veiem una {\bf imatge} de l'objecte, que en general serà diferent de la posició real de l'objecte. L'objectiu dels diagrames de raigs és poder deduir quina serà aquesta imatge: on es formarà, de quina mida, i en quina orientació. En general, la manera més visual de generar aquesta imatge és seguint un conjunt de tres regles senzilles sobre el comportament de tres tipus determinats de raigs. Per a més claredat, assignarem un color a cada un d'aquests raigs, per exemplificar-los a la Figura \ref{fig:DiagramaRaigs}. \\

\noindent
Aquestes regles són:
%
\begin{itemize}
    \item Regla 1 (vermell): Un raig que passa pel centre de la lent no es desvia.
    \item Regla 2 (verd): Un raig que entra paral·lel a l'eix surt en la direcció de $F'$.
    \item Regla 3 (blau): Un raig que entra en la direcció de $F$ surt paral·lel a l'eix.
\end{itemize}
%
\begin{figure}[H]
    \centering
    \includegraphics[width=\linewidth]{Figures/DiagramaRaigsConvergent.pdf}
    
    \vspace{1cm}
    
    \includegraphics[width=\linewidth]{Figures/DiagramaRaigsDivergent.pdf}
    \caption{Dos exemples de diagrames de raigs per una lent convergent (a dalt) i una lent divergent (a baix), seguint les tres regles amb els seus colors identificatius. En el cas de la lent convergent, la imatge és real, invertida i de la mateixa mida que l'objecte ($s' > 0, \, y' = - y$). En el cas de la lent divergent, la imatge és virtual, dreta, i més petita que l'objecte ($s' < 0, \, y > y' > 0$).}
    \label{fig:DiagramaRaigs}
\end{figure}
%
És important tenir en compte que un raig en la direcció de $F$ o $F'$ no necessàriament passarà per cap d'aquests punts, però sí que ho farà la seva perllongació. Si dibuixem els raigs corresponents a aquestes tres regles, provinents d'un punt de l'objecte, els raigs es creuaran en el punt on es genera la imatge de l'objecte. Normalment, se sol representar esquemàticament aquest objecte com una fletxa vertical.

La distància de l'objecte a la lent, en un diagrama de raigs, se sol simbolitzar amb la lletra $s$. Es considera positiu ($s > 0$) quan aquest es troba al costat esquerre de la lent, i l'anomenem objecte real. Si l'objecte està al costat dret de la lent, considerem $s < 0$ i l'anomenem objecte virtual. Això és un cas bastant exòtic que només es pot aconseguir quan el propi objecte és una imatge generada per una altra lent. En canvi, la distància de la imatge a la lent se simbolitza amb $s'$. Si la imatge es forma a la dreta de la lent, en diem imatge real i $s' > 0$. Si la imatge es forma a l'esquerra de la lent en diem imatge virtual i $s' < 0$. Només una imatge real es pot projectar en una pantalla. Una imatge virtual no es pot projectar en una pantalla perquè la llum realment no hi passa. 

L'equació que relaciona la focal $f$ de la lent amb les posicions de l'objecte $s$ i la imatge $s'$ en l'eix òptic s'anomena equació de la lent prima, es dedueix de trigonometria bàsica, i té la forma
%
\begin{equation}
    \boxed{\frac{1}{f} = \frac{1}{s} + \frac{1}{s'}}
    \label{eqn:lent_prima}
\end{equation}
%
Les altures de l'objecte ($y$) i la imatge ($y'$) es poden obtenir a partir de
%
\begin{equation}
    \boxed{\frac{s'}{s} = - \frac{y'}{y}}
\end{equation}
%
En aquest cas, $y, y' > 0$ significa cap a dalt i $y, y' < 0$ cap a baix. Com que el conveni de signes en òptica geomètrica és notablement complex, a la següent secció presentem un resum de convencions per a totes les variables utilitzades.

\subsection{Conveni de signes en òptica}

A la Figura \ref{fig:ConveniSignes} veiem un esquema bàsic d'un diagrama de raigs, on s'especifiquen els signes corresponents a les diverses variables depenent de la seva posició. Addicionalment, la Taula \ref{tab:ConvencionsOptica} resumeix el significat dels signes positius i negatius en cada una de les variables, per tal de simplificar els càlculs i evitar errors. En qualsevol cas, és molt convenient comparar els signes obtinguts de les expressions matemàtiques amb el que hem obtingut gràficament mitjançant el diagrama de raigs.
%
\begin{figure}[H]
    \centering
    \includegraphics[width=\linewidth]{Figures/ConveniSignes.pdf}
    \caption{Conveni de signes en els diagrames de raigs.}
    \label{fig:ConveniSignes}
\end{figure}
%
%
\begin{table}[H]
    \centering
    \begin{tabular}{|l|l|l|l|l|}
    \hline
                & {\bf Radis de la lent}           &             & {\bf Posició de l'objecte} \\ \hline
        $R > 0$ & Corbat cap a l'esquerra: $\big($ & $s > 0$     & A l'esquerra: objecte real \\ \hline
        $R < 0$ & Corbat cap a la dreta: $\big)$   & $s < 0$     & A la dreta: objecte virtual \\ \hline \hline
                & {\bf Distància focal}            &             & {\bf Posició de la imatge} \\ \hline       
        $f > 0$ & Lent convergent                  & $s' > 0$    & A la dreta: imatge real \\ \hline
        $f < 0$ & Lent divergent                   & $s' < 0$    & A l'esquerra: imatge virtual \\ \hline \hline
                & {\bf Potència de la lent}        &             & {\bf Altura de l'objecte/imatge} \\ \hline
        $p > 0$ & Lent convergent                  & $y, y' > 0$ & Objecte o imatge drets \\ \hline
        $p < 0$ & Lent divergent                   & $y, y' < 0$ & Objecte o imatge invertits \\ \hline
    \end{tabular}
    \caption{Conveni de signes per a les variables que utilitzem en òptica geomètrica.}
    \label{tab:ConvencionsOptica}
\end{table}

\subsection{Aberracions òptiques}

Abans d'acabar aquesta secció cal mencionar que, a la pràctica, els aparells òptics pateixen un conjunt d'efectes que perjudiquen la qualitat de les imatges. Aquests efectes s'anomenen aberracions òptiques, i fan que les imatges que n'obtenim apareguin distorsionades, desenfocades, o amb colors que no corresponen a la realitat. Per a aplicacions de precisió, això és un problema molt important que cal tenir en compte. Per tal de solucionar aquests problemes els instruments d'alta qualitat incorporen més elements òptics, com ara lents addicionals o vidres de diferents tipus. Això fa que els sistemes puguin ser molt més complexos que els casos idealitzats que hem estudiat en aquest capítol. A més, aquests elements correctius fan que augmenti enormement la dificultat de producció i el preu de mercat d'aquests instruments. Les aberracions òptiques més comunes són:

\begin{itemize}
    \item Aberració cromàtica: L'índex de refracció normalment depèn del color de la llum (dispersió). Això fa que la potència d'una lent depengui del color, i que apareguin arcs de Sant Martí a les imatges. 
    \item Aberració esfèrica: L'aproximació de lent prima no és exacta. Els raigs que incideixen paral·lels a l'eix òptic d'una lent amb superfícies esfèriques no convergeixen exactament al mateix focus. 
    \item Coma: Deformació dels objectes allunyats de l'eix òptic, que fa que la imatge de les fonts puntuals tingui forma de ``coma'' com un cometa. 
    \item Astigmatisme: Apareix quan la curvatura de la lent no és la mateixa en les direccions horitzontal i vertical. Les imatges apareixen deformades en una de les direccions.
    \item Curvatura de camp: Apareix quan els raigs a angles grans convergeixen més a prop de la lent. Provoca que la imatge d'un objecte pla aparegui corbada. 
    \item Distorsió: Apareix quan la potència de la lent varia amb la distància a l'eix òptic. El resultat és la formació d'una imatge deformada radialment.
\end{itemize}


\section{L'ull i els aparells òptics}

L'òptica i les lents no només són utilitzades en tecnologia, sinó que han aparegut de forma natural en els ulls dels animals a través de la selecció natural. La capacitat de percebre la llum ambiental procedent del Sol, i poder formar imatges del nostre entorn, suposa un enorme avantatge evolutiu. Per això, els ulls s'han perfeccionat enormement al llarg de l'evolució, i han arribat a ser aparells òptics extremadament sofisticats, especialment tenint en compte els materials i teixits disponibles per als éssers vius. Hi ha molts tipus d'ulls, cadascun amb la seva estructura i detalls de funcionament. En aquest capítol ens centrarem en els ulls dels vertebrats, sobretot mamífers, ocells i rèptils, tot i que alguns mol·luscs com el pop tenen ulls similars. Aquests ulls es basen en generar una imatge sobre una superfície, anomenada retina, coberta de cèl·lules fotosensibles.

Des de l'antiguitat, els éssers humans hem fet servir mètodes per millorar la nostra capacitat de visió. Hi ha certes dades que ens porten a pensar en l'ús de lents rudimentàries i lupes des de l'antiguitat, per tal de veure-hi millor o concentrar els raigs de sol. El 1609 Galileo Galilei va veure els satèl·lits de Júpiter amb un telescopi. Més tard, el 1665 Robert Hooke va ser capaç de veure les cèl·lules del suro fent servir un microscopi. Antonie van Leeuwenhoek va impulsar el naixement de la microbiologia, i va ser capaç de veure microorganismes per primera vegada. 

\subsection{L'ull i les ametropies}

La Figura \ref{fig:UllAmbEtiquetes} mostra un dibuix molt esquemàtic d'un ull humà. Els elements més rellevants són:
%
\begin{itemize}
    \item Cristal·lí: Lent convergent de focal variable, que fa convergir els raigs sobre la retina. Per tal de modificar la seva potència, canvia de curvatura per l'acció de músculs ciliars.
    %
    \item Retina: Superfície al fons de l'ull coberta de cèl·lules sensibles a la llum. Aquestes cèl·lules es coneixen com a cons i bastons.
    %
    \item Nervi òptic: Transmet el senyal des de les cèl·lules fotosensibles de la retina cap al cervell, amb la finalitat de processar la imatge.
    %
    \item Iris: Membrana que limita l’entrada de la llum, actua com el diafragma d'una càmera. És el responsable del color dels ulls.
    %
    \item Pupil·la: Obertura al centre de l'iris per on entra la llum cap al cristal·lí. Es dilata o es contrau depenent de la quantitat de llum disponible.
    %
    \item Còrnia: Part frontal transparent que cobreix l'iris.
\end{itemize}
%
El funcionament bàsic de l'ull és molt semblant al d'una càmera fotogràfica. Fa servir una lent convergent d'alta potència, el cristal·lí, per projectar una imatge invertida de l'exterior sobre una superfície fotoreceptora, la retina. La retina està recoberta de cèl·lules sensibles a la llum, els cons i els bastons, que codifiquen la imatge en senyals nervioses que s'envien a través del nervi òptic cap al cervell, on són interpretades. D'entre les cèl·lules fotosensibles de la retina, els bastons són molt més sensibles i abundants, mentre que els cons són responsables de la visió en color.
%
\begin{figure}[H]
    \centering
    \includegraphics[width=0.9\linewidth]{Figures/UllAmbEtiquetes.pdf}
    \caption{Dibuix esquemàtic de l'ull humà, amb les parts més importants indicades. Es pot veure com es forma la imatge d'un objecte sobre la retina.}
    \label{fig:UllAmbEtiquetes}
\end{figure}
%
De manera igual a una càmera fotogràfica, l'ull ha de ser capaç de veure objectes a distàncies diferents. Això significa que s'ha d'adaptar per tal que la imatge es projecti exactament sobre la retina. En altres paraules, tots els raigs procedents d'un mateix punt de l'objecte s'han de focalitzar en un mateix punt de la retina per tal de tenir una imatge nítida. Si això no fos així, la llum procedent del mateix punt de l'objecte es distribuiria per una àrea extensa de la retina, i la imatge es veuria desenfocada. Aquesta adaptació per tal d'ajustar-se a objectes a diferents distàncies s'anomena enfocament en el cas de les càmeres artificials, i {\bf acomodació} en el cas de l'ull. L'acomodació s'aconsegueix de diverses maneres: en els mamífers, ocells i rèptils es modifica la forma del cristal·lí, mentre que en els peixos i amfibis el cristal·lí manté la forma però es desplaça respecte a la retina. En general, el poder d'acomodació d'un ull té un rang limitat. És a dir, tindrem una distància màxima i una mínima a les que podrem enfocar:
%
\begin{itemize}
    \item Punt llunyà ($x_\infty$): Distància màxima a la que podem enfocar. En el cas d'un ull humà saludable, podem veure a distàncies arbitràriament grans, i per tant el punt llunyà és a l'infinit $x_\infty \to \infty$ . Veurem que això no és així en un ull miop.
    %
    \item Punt proper ($x_p$): Distància mínima a la que podem enfocar. En el cas d'un ull humà saludable, és d'uns $x_p =$ 25 cm. Aquest punt pot ser més llunyà en un ull hipermetrop, o més proper en un ull miop. 
\end{itemize}
%
La mida típica d'un ull humà és de 2,5 cm de diàmetre. Per tant, fent servir l'equació de la lent prima (\ref{eqn:lent_prima}), podem veure que la potència del cristal·lí en un ull saludable varia entre 40 dp (enfocant a l'infinit) i 44 dp (enfocant al punt proper). El mecanisme per modificar la curvatura, i per tant la potència, del cristal·lí és deformar-lo mitjançant els músculs ciliars. Quan els músculs estan relaxats, enfoquem al punt llunyà, mentre que la màxima acomodació permet enfocar al punt proper. A la pràctica, els ulls humans sovint pateixen ametropies o trastorns de refracció, és a dir, problemes en la capacitat d'acomodació que limiten la nostra capacitat d'enfocar objectes a diferents distàncies. Les més comunes són la miopia, la hipermetropia, l'astigmatisme i la presbícia. Ens centrarem en les primeres dues. 

La miopia que fa que la imatge dels objectes llunyans es formi a una distància massa curta del cristal·lí, abans d'arribar a la retina, tal com es veu a la Figura \ref{fig:Miopia}. És a dir, el punt llunyà de l'ull no és a una distància infinita. Si una persona miop intenta observar un objecte a una distància major que el seu punt llunyà, en veurà una imatge desenfocada. Aquest problema està causat normalment per la forma del globus ocular, que és massa allargat, o per una potència excessiva del cristal·lí. Per corregir la miopia es fan servir lents divergents, en forma d'ulleres o lents de contacte, que de forma efectiva redueixen la potència òptica del cristal·lí i permeten enfocar a distàncies grans.
%
\begin{figure}[H]
    \centering
    \includegraphics[width=0.9\linewidth]{Figures/MiopiaSenseCorregir.pdf}
    
    \vspace{0.2cm}
    
    \includegraphics[width=0.9\linewidth]{Figures/MiopiaCorregida.pdf}
    \caption{Esquema d'un ull amb miop (a dalt) que focalitza els raigs de llum abans de la retina. Per corregir la miopia s'utilitzen lents divergents (a baix) per poder formar imatges enfocades d'objectes llunyans.}
    \label{fig:Miopia}
\end{figure}

%
\begin{figure}[H]
    \centering
    \includegraphics[width=0.9\linewidth]{Figures/HipermetropiaSenseCorregir.pdf}
    
    \vspace{0.2cm}
    
    \includegraphics[width=0.9\linewidth]{Figures/HipermetropiaCorregida.pdf}
    \caption{Esquema d'un ull amb hipermetrop (a dalt) que focalitza els raigs de llum després de la retina. Per corregir la hipermetropia s'utilitzen lents convergents (a baix) per poder formar imatges enfocades d'objectes propers.}
    \label{fig:Hipermetropia}
\end{figure}
%
Per altra banda, la hipermetropia és el cas oposat a la miopia. En aquest cas, les imatges d'objectes propers es formen a una distància massa llarga del cristal·lí, mes enllà de la retina, tal com es veu a la Figura \ref{fig:Hipermetropia}. Això fa que el punt proper de l'ull estigui més allunyat del que es considera normal (25 cm). Novament, pot estar causada per la forma del globus ocular o per una poca capacitat d'acomodació de l'ull. Es corregeix amb lents convergents, que augmenten la capacitat de focalitzar els raigs de llum.

\subsection{La lupa i el microscopi}

Una limitació de l'ull és la seva capacitat de resolució, anomenada {\bf agudesa visual}. Un ull típicament pot distingir entre dos punts que es trobin separats, com a mínim, en un angle de 0,03$^\circ$ ($5\cdot10^{-4}$ rad). Si acostem un objecte fins al punt proper de 25 cm, i fent servir que $\sin \theta \sim \tan \theta \sim \theta$ per a angles petits, podem veure que la nostra agudesa visual ens permet resoldre fins a 0,125 mm. És a dir, a ull nu mai podrem distingir dos punts que estiguin a una distància menor de 0,125 mm, sempre els veurem com un sol punt. Això és problemàtic en el camp de la biologia, ja que l'escala típica d'una cèl·lula és de micròmetres (1 $\mu$m = 10$^{-6}$ m). Si ens interessen les estructures a nivell molecular la situació és encara més extrema, ja que ens situem en escales de nanòmetres (1 nm = 10$^{-9}$ m) o \AA ngströms (1 \AA = 10$^{-10}$ m). Per tal de millorar la nostra capacitat de visió a escales molt petites hi ha una varietat d'instruments òptics, els més utilitzats són la lupa i el microscopi. 

\paragraph{Lupa:} Una lupa senzilla és, simplement, una lent divergent. Si situem la mostra que volem observar a una distància lleugerament menor a la focal de la lupa ($s \lesssim f$), això produeix una imatge virtual de l'objecte molt augmentada ($y' \gg y$), i a una distància major ($|s'| \gg |s|$), com es veu a la Figura \ref{fig:Lupa}. D'aquesta manera, si acostem l'ull suficientment a la lupa, podem enfocar l'objecte a una distància molt menor que el punt proper, mantenint-ne la resolució angular. Per això, es defineix {\bf l'augment angular} $M$ de la lupa com el quocient entre l'angle que ocuparia l'objecte si el situéssim al punt proper $x_p$ i l'angle que ocupa a una distància $f$ de la lupa. Per semblança de triangles, podem veure que
%
\begin{equation}
    \boxed{M = \frac{\theta \text{ a una distància } f}{\theta \text{ a una distància } x_p} = \frac{x_p}{f}}
\end{equation}

\begin{figure}[H]
    \centering
    \includegraphics[width=\linewidth]{Figures/Lupa.pdf}
    \caption{Diagrama de raigs que mostra el funcionament d'una lupa. Si situem l'objecte entre la lent i el seu focus secundari, es forma una imatge virtual molt augmentada.}
    \label{fig:Lupa}
\end{figure}

\paragraph{Microscopi:} El model més senzill d'un microscopi òptic fa servir dues lents convergents consecutives, muntades al llarg d'un tub de suport. La primera lent, {\bf l'objectiu}, té una focal curta i genera una imatge real i invertida de la mostra dins el tub. La segona lent, {\bf l'ocular}, actua com una lupa que augmenta la imatge procedent de l'objectiu (que ara juga el paper d'objecte), tal com es veu a la Figura \ref{fig:Microscopi}. D'aquesta manera, l'augment total del microscopi $M_\text{total}$ es defineix com el producte dels augments de l'objectiu $M_\text{obj}$ i de l'ocular $M_\text{ocu}$. En aquest cas considerem que la mostra està a una distància de l'objectiu lleugerament major que la seva focal ($s_1 \gtrsim f_\text{obj}$), i que aquest genera una imatge real a una distància de l'ocular lleugerament menor a la seva focal ($s_2 \lesssim f_\text{ocu}$). D'aquesta manera, tenim
%
\begin{equation}
    \boxed{M_\text{total} = M_\text{obj} \cdot M_\text{ocu} = \frac{-s_1'}{s_1} \cdot  \frac{-s_2'}{s_2} = -\frac{s_1'}{f_\text{obj}} \cdot \frac{x_p}{f_\text{ocu}}}
\end{equation}
%
\begin{figure}[H]
    \centering
    \includegraphics[width=\linewidth]{Figures/Microscopi.pdf}
    \caption{Diagrama de raigs que mostra el funcionament d'un microscopi òptic senzill. Els raigs corresponents a l'objectiu s'han representat amb línies més primes, i els de l'ocular amb línies gruixudes.}
    \label{fig:Microscopi}
\end{figure}
%
A la pràctica, els microscopis solen ser molt més complicats que aquest model senzill, i tant l'objectiu com l'ocular estan formats per diverses lents. Això millora la qualitat de la imatge i redueix les aberracions que es produirien si féssim servir una sola lent. Molts microscopis, a més, tenen diversos objectius amb diversos augments que podem anar intercanviant mentre fem una observació.

\section{Limitacions de l'òptica}

Tot i que els microscopis òptics ens permeten veure estructures molt més petites que a ull nu, també tenen una limitació en la mida dels objectes que ens permeten resoldre. Aquesta limitació prové del fet que la llum és una ona, i per tant pateix un fenomen anomenat difracció. Quan les ones de llum passen a través d'estructures de mida comparable a la longitud d'ona, a causa del principi de Huygens, es produeixen patrons d'interferència entre diverses parts de l'ona. Això fa que, per exemple, la imatge d'una obertura circular adopti una forma característica en forma de cercles concèntrics anomenada disc d'Airy, tal com es veu a la Figura \ref{fig:Difraccio}. A causa d'això, dos punts massa propers no es podran distingir clarament. 
%
\begin{figure}[H]
    \centering
    \includegraphics[width=\linewidth]{Figures/Difraccio.pdf}
    \caption{Simulació del patró de difracció de la llum en passar per una petita obertura circular (esquerra) on es pot veure la formació del disc d'Airy. Quan tenim dues obertures molt juntes (dreta), en aquest cas $d = 3r$, no podem resoldre-les bé. }
    \label{fig:Difraccio}
\end{figure}
%
El límit de resolució d'un microscopi ve donat pel límit de difracció d'Abbe, és a dir
%
\begin{equation}
    D_\text{min} = \frac{\lambda}{2 n \sin \theta} = \frac{\lambda}{2 \text{NA}}\, ,
\end{equation}
%
on $\lambda$ és la longitud d'ona, $n$ és l'índex de refracció del medi, i $\theta$ és l'angle entre el centre de la mostra i la vora de l'objectiu. Normalment es fa servir el nombre NA $= n \sin \theta$, anomenat apertura numèrica de l'objectiu, que sol ser proper a 1. Per tant, si considerem llum blava amb una longitud d'ona d'uns 500 nm, veiem que el límit de resolució estarà entorn als 200 nm. Això vol dir que un microscopi òptic no podrà tenir mai un augment més gran que $M_\text{max} \approx 500$.

\subsection{El microscopi electrònic}

Si volem distingir clarament objectes per sota del micròmetre, necessitem microscopis que treballin amb longituds d'ona molt més petites que les de la llum visible. Una possible solució és utilitzar la dualitat ona-partícula de la mecànica quàntica. Segons la teoria quàntica, les partícules es comporten com una ona de probabilitat. La seva longitud d'ona ve donada per la llei enunciada per Louis de Broglie,
%
\begin{equation}
\lambda = \frac{h}{p}\, ,
\end{equation}
%
on $h = 6,626 \cdot 10^{-34}$ m$^2$·kg/s és la constant de Planck i $p$ és el moment lineal de la partícula. Si accelerem aquest electró, de massa $m_e = 9,109 \cdot 10^{-31}$ kg i càrrega $e = 1,602 \cdot 10^{-19}$ C, a través d'una diferència de potencial $V$, tindrem
%
\begin{equation}
    p = m_e v\, , \quad K = eV = \frac12 m_e v^2 = \frac{p^2}{2m_e}\, ,
\end{equation}
%
i per tant
%
\begin{equation}
    \lambda = \frac{h}{\sqrt{2 e m_e V}}\, .
\end{equation}
%
Per a una diferència de potencial de 50 kV, obtenim una longitud d'ona de $\lambda = 5,5 \cdot 10^{-3}$ nm, que ens dóna un límit de resolució de $D_\text{min} \approx 0,2$ nm. Això és 1000 vegades més petit que amb el microscopi òptic! 
%
\begin{figure}[H]
    \centering
    \includegraphics[width=0.6\linewidth]{Figures/MicroscopiElectronic.jpeg}
    \caption{Microscopi electrònic de l'any 1974, en exposició al vestíbul de la Facultat de Física de la Universitat de Barcelona.}
    \label{fig:MicroscopiElectronic}
\end{figure}
%
Els microscopis electrònics, com el que es veu a la Figura \ref{fig:MicroscopiElectronic}, són utilitzats molt habitualment en totes les àrees de la ciència i la tecnologia, especialment en l'àmbit biomèdic. N'hi ha de molts tipus, i amb diferents prestacions. Com a norma general, aquests microscopis tenen un canó d'electrons que accelera aquestes partícules utilitzant un alt voltatge. Després fan servir una successió de lents electromagnètiques que utilitzen camps elèctrics i magnètics per desviar els feixos d'electrons cap a un detector, passant per la mostra que es vol observar.

\subsection{Difracció de raigs X}

\begin{figure}[H]
    \centering
    \includegraphics[width=0.6\linewidth]{Figures/Photo51.jpg}
    \caption{Fotografia 51, obtinguda l'any 1952 per Rosalind Franklin i Raymond Gosling. Aquesta fotografia va permetre el descobriment de l'estructura de doble hèlix de l'ADN.}
    \label{fig:Photo51}
\end{figure}
%
Per a estructures molt petites, en què ens interessa saber la disposició espacial dels àtoms en un cristall o una molècula, podem recórrer a la difracció. És a dir, el mateix fenomen de les ones que limita la resolució dels microscopis pot ser utilitzat al nostre favor. L'estudi matemàtic de la difracció, amb tècniques com les transformades de Fourier, ens permet relacionar els patrons de difracció que observem amb les estructures que els han causat. Per fer això, normalment s'utilitzen els raigs X, ja que tenen longituds d'ona properes a l'escala atòmica. Es poden obtenir de diverses maneres, entre elles els tubs de raigs X i els sincrotrons.

La difracció de raigs X va tenir una importància cabdal en el descobriment de l'estructura de l'ADN. Aquest descobriment va ser possible gràcies als experiments de difracció de raigs X realitzats per Rosalind Franklin i el seu estudiant Raymond Gosling al King's College de Londres. La imatge més famosa és la {\bf Fotografia 51}, que es va obtenir el dia 2 de maig de l'any 1952. Aquesta imatge, que es veu a la Figura \ref{fig:Photo51}, va ser decisiva per determinar que l'ADN estava organitzat en una estructura de doble hèlix. La Fotografia 51 ha esdevingut un símbol del naixement de la biologia moderna.



\end{document}